\input{preamble}

% OK, start here.
%
\begin{document}

\title{Théorie des déformations}


\maketitle

\phantomsection
\label{section-phantom}

\tableofcontents

\section{Introduction}
\label{section-introduction}

\noindent
Le but de ce chapitre est de donner une introduction (relativement) élémentaire à
la théorie des déformations des modules, des morphismes, etc. Nous y traitons
les résultats que l'on peut démontrer au moyen du complexe cotangent naïf. Dans
le chapitre consacré au complexe cotangent, nous étendrons un peu ces résultats.
Le lecteur averti pourra consulter le traité d'Illusie sur ce
sujet ; voir \cite{cotangent}.





\section{Déformations d'anneaux et complexe cotangent naïf}
\label{section-deformations}

\noindent
Dans cette section, nous utilisons le complexe cotangent naïf pour faire un peu
de théorie des déformations. Partons d'un homomorphisme surjectif d'anneaux $A' \to A$
dont le noyau est un idéal $I$ de carré nul. Supposons en outre
donnés un homomorphisme d'anneaux $A \to B$, un $B$-module $N$ et une application $A$-linéaire
$c : I \to N$. Nous cherchons dans cette section à déterminer
le point d'interrogation qui complète le diagramme suivant :
\begin{equation}
\label{equation-to-solve}
\vcenter{
\xymatrix{
0 \ar[r] & N \ar[r] & {?} \ar[r] & B \ar[r] & 0 \\
0 \ar[r] & I \ar[u]^c \ar[r] & A' \ar[u] \ar[r] & A \ar[u] \ar[r] & 0
}
}
\end{equation}
et à mesurer en outre l'unicité de la solution (si elle existe). Plus précisément,
nous cherchons une surjection de $A'$-algèbres $B' \to B$ dont le noyau soit
un idéal de carré nul identifié à
$N$, telle que $A' \to B'$ induise l'application donnée $c$.
Nous dirons que $B'$ est une {\it solution} de (\ref{equation-to-solve}).

\begin{lemma}
\label{lemma-huge-diagram}
Étant donné un diagramme commutatif
$$
\xymatrix{
& 0 \ar[r] & N_2 \ar[r] & B'_2 \ar[r] & B_2 \ar[r] & 0 \\
& 0 \ar[r]|\hole & I_2 \ar[u]_{c_2} \ar[r] &
A'_2 \ar[u] \ar[r]|\hole & A_2 \ar[u] \ar[r] & 0 \\
0 \ar[r] & N_1 \ar[ruu] \ar[r] & B'_1 \ar[r] & B_1 \ar[ruu] \ar[r] & 0 \\
0 \ar[r] & I_1 \ar[ruu]|\hole \ar[u]^{c_1} \ar[r] &
A'_1 \ar[ruu]|\hole \ar[u] \ar[r] & A_1 \ar[ruu]|\hole \ar[u] \ar[r] & 0
}
$$
dont les faces avant et arrière sont des solutions de (\ref{equation-to-solve}), on a :
\begin{enumerate}
\item Il existe un élément canonique de
$\Ext^1_{B_1}(\NL_{B_1/A_1}, N_2)$
dont l'annulation est une condition nécessaire et suffisante pour qu'il existe
un homomorphisme d'anneaux $B'_1 \to B'_2$ complétant le diagramme.
\item S'il existe un homomorphisme $B'_1 \to B'_2$ complétant le diagramme,
l'ensemble de tous ces homomorphismes est un espace principal homogène sous
$\Hom_{B_1}(\Omega_{B_1/A_1}, N_2)$.
\end{enumerate}
\end{lemma}

\begin{proof}
Soit $E = B_1$ considéré comme un ensemble.
Considérons la surjection $A_1[E] \to B_1$, de noyau $J$, utilisée
pour définir le complexe cotangent naïf par la formule
$$
\NL_{B_1/A_1} = (J/J^2 \to \Omega_{A_1[E]/A_1} \otimes_{A_1[E]} B_1)
$$
dans
Algèbre, section \ref{algebra-section-netherlander}.
Puisque $\Omega_{A_1[E]/A_1} \otimes B_1$ est un
$B_1$-module libre, on a
$$
\Ext^1_{B_1}(\NL_{B_1/A_1}, N_2) =
\frac{\Hom_{B_1}(J/J^2, N_2)}
{\Hom_{B_1}(\Omega_{A_1[E]/A_1} \otimes B_1, N_2)}
$$
Nous allons construire une obstruction dans le module de droite.
Soit $J' = \Ker(A'_1[E] \to B_1)$. Remarquons qu'il existe une surjection
$J' \to J$ dont le noyau est $I_1A'_1[E]$.
Pour tout $e \in E$, notons $x_e \in A_1[E]$ la variable correspondante.
Choisissons un relèvement $y_e \in B'_1$ de l'image de $x_e$ dans $B_1$ et
un relèvement $z_e \in B'_2$ de l'image de $x_e$ dans $B_2$.
Ces choix déterminent des homomorphismes de $A'_1$-algèbres
$$
A'_1[E] \to B'_1 \quad\text{et}\quad A'_1[E] \to B'_2
$$
Le premier donne une application $J' \to N_1$, $f' \mapsto f'(y_e)$,
et le second une application $J' \to N_2$, $f' \mapsto f'(z_e)$.
Un calcul montre que ces applications annulent $(J')^2$.
Comme le carré gauche du diagramme (où interviennent $c_1$ et $c_2$)
est commutatif, ces applications coïncident sur $I_1A'_1[E]$ comme applications à valeurs dans $N_2$.
Remarquons que $B'_1$ est la somme amalgamée de $J' \to A'_1[E]$ et $J' \to N_1$.
Ainsi, si les applications $J' \to N_1 \to N_2$ et $J' \to N_2$ coïncident, on
obtient un homomorphisme $B'_1 \to B'_2$ complétant le diagramme.
On définit donc l'obstruction comme la classe de l'application
$$
J/J^2 \to N_2,\quad f \mapsto f'(z_e) - \nu(f'(y_e))
$$
où $\nu : N_1 \to N_2$ est l'application donnée et où $f' \in J'$
est un relèvement de $f$. Les remarques précédentes montrent que cette application est bien définie.
Nous sommes libres
de remplacer nos choix de $z_e$ par $z_e + \delta_{2, e}$
et de $y_e$ par $y_e + \delta_{1, e}$, pour certains $\delta_{i, e} \in N_i$.
Cela remplace l'application précédente par
$$
f \mapsto f'(z_e + \delta_{2, e}) - \nu(f'(y_e + \delta_{1, e})) =
f'(z_e) - \nu(f'(z_e)) +
\sum (\delta_{2, e} - \nu(\delta_{1, e}))\frac{\partial f}{\partial x_e}
$$
Cela signifie exactement que l'on modifie l'application $J/J^2 \to N_2$
par la composée $J/J^2 \to \Omega_{A_1[E]/A_1} \otimes B_1 \to N_2$,
où la seconde application envoie $\text{d}x_e$ sur
$\delta_{2, e} - \nu(\delta_{1, e})$. Notre obstruction est donc bien définie
et elle est nulle si et seulement s'il existe un relèvement.

\medskip\noindent
L'assertion (2) résulte de l'observation suivante : étant donnés deux homomorphismes
$\varphi, \psi : B'_1 \to B'_2$ complétant le diagramme,
$\varphi - \psi$ se factorise par une application $D : B_1 \to N_2$ qui
est une $A_1$-dérivation :
\begin{align*}
D(fg) & = \varphi(f'g') - \psi(f'g') \\
& =
\varphi(f')\varphi(g') - \psi(f')\psi(g') \\
& =
(\varphi(f') - \psi(f'))\varphi(g') + \psi(f')(\varphi(g') - \psi(g')) \\
& =
gD(f) + fD(g)
\end{align*}
Ainsi, $D$ correspond à une unique application $B_1$-linéaire
$\Omega_{B_1/A_1} \to N_2$. Réciproquement, une telle application linéaire
donne une dérivation $D$ et, étant donné un homomorphisme d'anneaux $\psi : B'_1 \to B'_2$
complétant le diagramme,
l'application $\psi + D$ est un autre homomorphisme d'anneaux qui le complète.
\end{proof}

\begin{lemma}
\label{lemma-choices}
S'il existe une solution de (\ref{equation-to-solve}), l'ensemble des
classes d'isomorphisme de solutions est un espace principal homogène sous
$\Ext^1_B(\NL_{B/A}, N)$.
\end{lemma}

\begin{proof}
Remarquons d'emblée que, pour deux solutions $B'_1$ et $B'_2$
de (\ref{equation-to-solve}), le lemme \ref{lemma-huge-diagram} fournit un
élément d'obstruction $o(B'_1, B'_2) \in \Ext^1_B(\NL_{B/A}, N)$
à l'existence d'un homomorphisme $B'_1 \to B'_2$. Cet élément est clairement
l'obstruction à l'existence d'un isomorphisme et distingue donc
les classes d'isomorphisme. Pour achever la démonstration, il suffit alors de
montrer qu'étant donnés une solution $B'$ et un élément
$\xi \in \Ext^1_B(\NL_{B/A}, N)$,
on peut trouver une seconde solution $B'_\xi$ telle que
$o(B', B'_\xi) = \xi$.

\medskip\noindent
Soit $E = B$ considéré comme un ensemble. Considérons la surjection $A[E] \to B$ de noyau
$J$ utilisée pour définir le complexe cotangent naïf par la formule
$$
\NL_{B/A} = (J/J^2 \to \Omega_{A[E]/A} \otimes_{A[E]} B)
$$
dans Algèbre, section \ref{algebra-section-netherlander}.
Puisque $\Omega_{A[E]/A} \otimes B$ est un $B$-module libre, on a
$$
\Ext^1_B(\NL_{B/A}, N) =
\frac{\Hom_B(J/J^2, N)}
{\Hom_B(\Omega_{A[E]/A} \otimes B, N)}
$$
On peut donc représenter $\xi$ par la classe d'un morphisme $\delta : J/J^2 \to N$.

\medskip\noindent
Pour tout $e \in E$, notons $x_e \in A[E]$ la variable correspondante.
Choisissons un relèvement $y_e \in B'$ de l'image de $x_e$ dans $B$.
Ces choix déterminent un homomorphisme de $A'$-algèbres $\varphi : A'[E] \to B'$.
Soit $J' = \Ker(A'[E] \to B)$. Remarquons que $\varphi$ induit une application
$\varphi|_{J'} : J' \to N$ et que $B'$ est la somme amalgamée, comme dans le
diagramme suivant :
$$
\xymatrix{
0 \ar[r] & N \ar[r] & B' \ar[r] & B \ar[r] & 0 \\
0 \ar[r] & J' \ar[u]^{\varphi|_{J'}} \ar[r] & A'[E] \ar[u] \ar[r] &
B \ar[u]_{=} \ar[r] & 0
}
$$
Soit $\psi : J' \to N$ la somme de l'application $\varphi|_{J'}$ et de la
composée
$$
J' \to J'/(J')^2 \to J/J^2 \xrightarrow{\delta} N.
$$
La somme amalgamée suivant $\psi$ est alors une autre extension d'anneaux $B'_\xi$
qui s'insère dans un diagramme comme ci-dessus. Un calcul montre que
$o(B', B'_\xi) = \xi$, comme voulu.
\end{proof}

\begin{lemma}
\label{lemma-extensions-of-algebras}
Soit $A$ un anneau. Soit $B$ une $A$-algèbre. Soit $N$ un $B$-module.
L'ensemble des classes d'isomorphisme d'extensions de $A$-algèbres
$$
0 \to N \to B' \to B \to 0
$$
où $N$ est un idéal de carré nul est canoniquement en bijection avec
$\Ext^1_B(\NL_{B/A}, N)$.
\end{lemma}

\begin{proof}
Pour le démontrer, appliquons les résultats précédents au cas où
(\ref{equation-to-solve}) est donné par le diagramme
$$
\xymatrix{
0 \ar[r] & N \ar[r] & {?} \ar[r] & B \ar[r] & 0 \\
0 \ar[r] & 0 \ar[u] \ar[r] & A \ar[u] \ar[r]^{\text{id}} & A \ar[u] \ar[r] & 0
}
$$
Le lemme résulte donc du lemme \ref{lemma-choices}
et du fait qu'il existe une solution, à savoir $N \oplus B$.
(Voir la remarque ci-dessous pour une construction directe de la bijection.)
\end{proof}

\begin{remark}
\label{remark-extensions-of-algebras}
Soient $A \to B$ et $N$ comme dans le lemme \ref{lemma-extensions-of-algebras}.
Soit $\alpha : P \to B$ une présentation de $B$ sur $A$ ; voir
Algèbre, section \ref{algebra-section-netherlander}.
Pour $J = \Ker(\alpha)$, le complexe cotangent naïf $\NL(\alpha)$
associé à $\alpha$ est le complexe $J/J^2 \to \Omega_{P/A} \otimes_P B$.
On a
$$
\Ext^1_B(\NL(\alpha), N) =
\Coker\left(\Hom_B(\Omega_{P/A} \otimes_P B, N) \to \Hom_B(J/J^2, N)\right)
$$
car $\Omega_{P/A}$ est un module libre.
Considérons une extension $0 \to N \to B' \to B \to 0$ comme dans le lemme.
Puisque $P$ est une algèbre polynomiale sur $A$, on peut
relever $\alpha$ en un homomorphisme de $A$-algèbres $\alpha' : P' \to B'$.
Alors $\alpha'|_J : J \to N$ se factorise sous la forme $J \to J/J^2 \to N$
car $N$ est de carré nul dans $B'$. Le lemme envoie notre extension
sur la classe de cette application $J/J^2 \to N$ dans le conoyau affiché.
\end{remark}












\begin{lemma}
\label{lemma-extensions-of-algebras-functorial}
Soient des homomorphismes d'anneaux $A \to B \to C$, un $B$-module $M$, un $C$-module $N$,
une application $B$-linéaire $c : M \to N$, et des extensions de
$A$-algèbres à noyaux de carré nul
\begin{enumerate}
\item[(a)] $0 \to M \to B' \to B \to 0$ correspondant à
$\xi \in \Ext^1_B(\NL_{B/A}, M)$, et
\item[(b)] $0 \to N \to C' \to C \to 0$ correspondant à
$\zeta \in \Ext^1_C(\NL_{C/A}, N)$.
\end{enumerate}
Voir le lemme \ref{lemma-extensions-of-algebras}.
Il existe alors un homomorphisme de $A$-algèbres $B' \to C'$ compatible avec
$B \to C$ et $c$ si et seulement si $\xi$ et $\zeta$
s'envoient sur le même élément de
$\Ext^1_B(\NL_{B/A}, N)$.
\end{lemma}

\begin{proof}
L'énoncé a un sens car on dispose des applications
$$
\Ext^1_B(\NL_{B/A}, M) \to \Ext^1_B(\NL_{B/A}, N)
$$
induites par l'application $M \to N$, ainsi que de
$$
\Ext^1_C(\NL_{C/A}, N) \to \Ext^1_B(\NL_{C/A}, N) \to \Ext^1_B(\NL_{B/A}, N)
$$
où la première flèche utilise le foncteur de restriction $D(C) \to D(B)$
et la seconde l'application canonique de complexes
$\NL_{B/A} \to \NL_{C/A}$. L'énoncé du lemme se déduit du
lemme \ref{lemma-huge-diagram} appliqué au diagramme
$$
\xymatrix{
& 0 \ar[r] & N \ar[r] & C' \ar[r] & C \ar[r] & 0 \\
& 0 \ar[r]|\hole & 0 \ar[u] \ar[r] &
A \ar[u] \ar[r]|\hole & A \ar[u] \ar[r] & 0 \\
0 \ar[r] & M \ar[ruu] \ar[r] & B' \ar[r] & B \ar[ruu] \ar[r] & 0 \\
0 \ar[r] & 0 \ar[ruu]|\hole \ar[u] \ar[r] &
A \ar[ruu]|\hole \ar[u] \ar[r] & A \ar[ruu]|\hole \ar[u] \ar[r] & 0
}
$$
et d'une compatibilité entre les constructions des démonstrations
des lemmes \ref{lemma-extensions-of-algebras} et \ref{lemma-huge-diagram},
dont nous omettons l'énoncé et la démonstration. (Voir la remarque ci-dessous pour un argument direct.)
\end{proof}

\begin{remark}
\label{remark-extensions-of-algebras-functorial}
Soient $A \to B \to C$, $M$, $N$, $c : M \to N$,
$0 \to M \to B' \to B \to 0$, $\xi \in \Ext^1_B(\NL_{B/A}, M)$,
$0 \to N \to C' \to C \to 0$ et $\zeta \in \Ext^1_C(\NL_{C/A}, N)$ comme dans le
lemme \ref{lemma-extensions-of-algebras-functorial}.
En prenant la somme amalgamée le long de $c : M \to N$, on construit une extension
$$
\xymatrix{
0 \ar[r] & N \ar[r] & B'_1 \ar[r] & B \ar[r] & 0 \\
0 \ar[r] & M \ar[u]^c \ar[r] & B' \ar[u] \ar[r] & B \ar[u] \ar[r] & 0
}
$$
en posant $B'_1 = (N \times B')/M$, où $M$ est plongé
antidiagonalement. En prenant le produit fibré le long de $B \to C$, on construit une extension
$$
\xymatrix{
0 \ar[r] & N \ar[r] & C' \ar[r] & C \ar[r] & 0 \\
0 \ar[r] & N \ar[u] \ar[r] & B'_2 \ar[u] \ar[r] & B \ar[u] \ar[r] & 0
}
$$
en posant $B'_2 = C' \times_C B$ (produit fibré d'anneaux). Une chasse au diagramme élémentaire
montre qu'il existe un homomorphisme de $A$-algèbres $B' \to C'$
compatible avec $B \to C$ et $c$ si et seulement si $B'_1$ est isomorphe
à $B'_2$ comme extension de $A$-algèbres de $B$ par $N$. Pour vérifier
le lemme \ref{lemma-extensions-of-algebras-functorial},
il suffit donc de montrer que $B'_1$ correspond, par la bijection du
lemme \ref{lemma-extensions-of-algebras},
à l'image de $\xi$ par l'application
$\Ext^1_B(\NL_{B/A}, M) \to \Ext^1_B(\NL_{B/A}, N)$,
et que $B'_2$ correspond à l'image de $\zeta$ par l'application
$\Ext^1_C(\NL_{C/A}, N) \to \Ext^1_B(\NL_{B/A}, N)$.
La première de ces assertions résulte immédiatement de la construction
de la classe dans la remarque \ref{remark-extensions-of-algebras}.
Pour la seconde, choisissons un diagramme commutatif
$$
\xymatrix{
Q \ar[r]_\beta & C \\
P \ar[u]^\varphi \ar[r]^\alpha & B \ar[u]
}
$$
de $A$-algèbres, tel que $\alpha$ soit une présentation de $B$ sur $A$
et $\beta$ une présentation de $C$ sur $A$. Voir la
remarque \ref{remark-extensions-of-algebras} et les références qui y sont données.
Posons $J = \Ker(\alpha)$ et $K = \Ker(\beta)$. L'application $\varphi$
induit une application de complexes $\NL(\alpha) \to \NL(\beta)$
et, en particulier, $\bar\varphi : J/J^2 \to K/K^2$.
Choisissons un homomorphisme de $A$-algèbres $\beta' : Q \to C'$
qui relève $\beta$. Alors
$\alpha' = (\beta' \circ \varphi, \alpha) : P \to B'_2 = C' \times_C B$
est un relèvement de $\alpha$. Pour ces choix, la composée de l'application
$K/K^2 \to N$ induite par $\beta'$ et de l'application $\bar\varphi : J/J^2 \to K/K^2$
est la restriction de $\alpha'$ à $J/J^2$.
En explicitant les constructions de nos classes dans la
remarque \ref{remark-extensions-of-algebras},
on voit bien que
$B'_2$ correspond à l'image de $\zeta$ par l'application
$\Ext^1_C(\NL_{C/A}, N) \to \Ext^1_B(\NL_{B/A}, N)$.
\end{remark}

\begin{lemma}
\label{lemma-parametrize-solutions}
Soient $0 \to I \to A' \to A \to 0$, $A \to B$ et $c : I \to N$ comme dans
(\ref{equation-to-solve}). Notons $\xi \in \Ext^1_A(\NL_{A/A'}, I)$
l'élément correspondant à l'extension $A'$ de $A$ par $I$ par le
lemme \ref{lemma-extensions-of-algebras}. L'ensemble des classes d'isomorphisme
de solutions est canoniquement en bijection avec la fibre de
$$
\Ext^1_B(\NL_{B/A'}, N) \to \Ext^1_A(\NL_{A/A'}, N)
$$
au-dessus de l'image de $\xi$.
\end{lemma}

\begin{proof}
D'après le lemme \ref{lemma-extensions-of-algebras} appliqué à $A' \to B$ et au
$B$-module $N$, les éléments $\zeta$ de $\Ext^1_B(\NL_{B/A'}, N)$
paramètrent les extensions $0 \to N \to B' \to B \to 0$ de $A'$-algèbres.
D'après le lemme \ref{lemma-extensions-of-algebras-functorial} appliqué
à $A' \to A \to B$ et $c : I \to N$, il existe un homomorphisme de $A'$-algèbres
$A' \to B'$ compatible avec $c$ et $A \to B$ si et seulement si
$\zeta$ s'envoie sur $\xi$. Cela revient bien sûr à dire que $B'$ est une
solution de (\ref{equation-to-solve}).
\end{proof}

\begin{remark}
\label{remark-parametrize-solutions}
Remarquons que, dans la situation du lemme \ref{lemma-parametrize-solutions},
on a
$$
\Ext^1_A(\NL_{A/A'}, N) =
\Ext^1_B(\NL_{A/A'} \otimes_A^\mathbf{L} B, N) =
\Ext^1_B(\NL_{A/A'} \otimes_A B, N)
$$
La première égalité résulte de
Compléments d'algèbre, lemme \ref{more-algebra-lemma-tensor-hom-adjoint}, et
la seconde de
Compléments d'algèbre, lemme \ref{more-algebra-lemma-tensor-NL}.
On dispose d'applications de complexes
$$
\NL_{A/A'} \otimes_A B \to \NL_{B/A'} \to \NL_{B/A}
$$
qui forment presque un triangle distingué ; voir
Algèbre, lemme \ref{algebra-lemma-exact-sequence-NL}.
Si elles formaient un triangle distingué, on conclurait
que l'image de $\xi$ dans $\Ext^2_B(\NL_{B/A}, N)$
serait l'obstruction à l'existence d'une solution de
(\ref{equation-to-solve}).
\end{remark}

\noindent
Si notre homomorphisme d'anneaux $A \to B$ est d'intersection complète locale, il existe
une solution. Il s'agit d'un résultat de relèvement ; remarquons que,
pour les homomorphismes syntomiques, nous avons démontré un résultat de relèvement assez fort dans
Lissage des morphismes d'anneaux, Proposition \ref{smoothing-proposition-lift-smooth}.

\begin{lemma}
\label{lemma-existence-lci}
Si $A \to B$ est un homomorphisme d'anneaux d'intersection complète locale, alors
il existe une solution de (\ref{equation-to-solve}).
\end{lemma}

\begin{proof}[Première démonstration]
Écrivons $B = A[x_1, \ldots, x_n]/J$. D'après Compléments d'algèbre, définition
\ref{more-algebra-definition-local-complete-intersection},
l'idéal $J$ est Koszul-régulier. Il s'ensuit que $J$ est $H_1$-régulier et
quasi-régulier ; voir Compléments d'algèbre, section \ref{more-algebra-section-ideals}.
Soit $J' \subset A'[x_1, \ldots, x_n]$
l'image réciproque de $J$. Notons $I[x_1, \ldots, x_n]$ le
noyau de $A'[x_1, \ldots, x_n] \to A[x_1, \ldots, x_n]$.
D'après Compléments d'algèbre, lemme
\ref{more-algebra-lemma-conormal-sequence-H1-regular-ideal}, on a
$I[x_1, \ldots, x_n] \cap (J')^2 = J'I[x_1, \ldots, x_n] =
JI[x_1, \ldots, x_n]$. On obtient donc une suite exacte courte
$$
0 \to I \otimes_A B \to J'/(J')^2 \to J/J^2 \to 0
$$
Puisque $J/J^2$ est projectif (Compléments d'algèbre, lemme
\ref{more-algebra-lemma-quasi-regular-ideal-finite-projective}),
on peut choisir un scindage de cette suite
$$
J'/(J')^2 = I \otimes_A B \oplus J/J^2 
$$
Soit $(J')^2 \subset J'' \subset J'$ l'ensemble des éléments qui s'envoient dans le
second facteur de la décomposition précédente. Alors
$$
0 \to I \otimes_A B \to A'[x_1, \ldots, x_n]/J'' \to B \to 0
$$
est une solution de (\ref{equation-to-solve}) avec $N = I \otimes_A B$.
Le cas général s'obtient en prenant la somme amalgamée le long de l'application donnée
$I \otimes_A B \to N$.
\end{proof}

\begin{proof}[Seconde démonstration]
Veuillez lire la remarque \ref{remark-parametrize-solutions}
avant cette démonstration. D'après
Compléments d'algèbre, lemme \ref{more-algebra-lemma-transitive-lci-at-end},
les applications $\NL_{A/A'} \otimes_A B \to \NL_{B/A'} \to \NL_{B/A}$
forment bien un triangle distingué dans $D(B)$.
Il suffit donc de montrer que $\Ext^2_B(\NL_{B/A}, N)$ s'annule.
D'après Compléments d'algèbre, lemme \ref{more-algebra-lemma-lci-NL},
le complexe $\NL_{B/A}$ est parfait, d'amplitude de Tor contenue dans
$[-1, 0]$. Il s'ensuit que notre $\Ext^2$ s'annule, par exemple
d'après Compléments d'algèbre, lemme \ref{more-algebra-lemma-splitting-unique}, partie (1).
\end{proof}






\section{Épaississements d'espaces annelés}
\label{section-thickenings-spaces}

\noindent
Dans les quelques sections qui suivent, nous utiliserons les notions suivantes :
\begin{enumerate}
\item Un faisceau d'idéaux $\mathcal{I} \subset \mathcal{O}_{X'}$ sur
un espace annelé $(X', \mathcal{O}_{X'})$ est {\it localement nilpotent}
si toute section locale de $\mathcal{I}$ est localement nilpotente.
Comparer avec Algèbre, Item \ref{algebra-item-ideal-locally-nilpotent}.
\item Un {\it épaississement} d'espaces annelés est un morphisme
$i : (X, \mathcal{O}_X) \to (X', \mathcal{O}_{X'})$ d'espaces annelés
tel que
\begin{enumerate}
\item $i$ induise un homéomorphisme $X \to X'$ ;
\item l'application $i^\sharp : \mathcal{O}_{X'} \to i_*\mathcal{O}_X$
soit surjective ; et
\item le noyau de $i^\sharp$ soit un faisceau d'idéaux localement nilpotent.
\end{enumerate}
\item Un {\it épaississement du premier ordre} d'espaces annelés est un épaississement
$i : (X, \mathcal{O}_X) \to (X', \mathcal{O}_{X'})$ d'espaces annelés
tel que $\Ker(i^\sharp)$ soit de carré nul.
\item Il est clair comment définir les {\it morphismes d'épaississements},
les {\it morphismes d'épaississements au-dessus d'un espace annelé de base}, etc.
\end{enumerate}
Si $i : (X, \mathcal{O}_X) \to (X', \mathcal{O}_{X'})$ est un épaississement
d'espaces annelés, nous identifions les espaces topologiques sous-jacents
et considérons $\mathcal{O}_X$, $\mathcal{O}_{X'}$ et
$\mathcal{I} = \Ker(i^\sharp)$ comme des faisceaux sur $X = X'$. On obtient
une suite exacte courte
$$
0 \to \mathcal{I} \to \mathcal{O}_{X'} \to \mathcal{O}_X \to 0
$$
de $\mathcal{O}_{X'}$-modules. D'après
Modules, lemme \ref{modules-lemma-i-star-equivalence},
la catégorie des $\mathcal{O}_X$-modules est équivalente à la catégorie
des $\mathcal{O}_{X'}$-modules annulés par $\mathcal{I}$. En particulier,
si $i$ est un épaississement du premier ordre,
$\mathcal{I}$ est un $\mathcal{O}_X$-module.

\begin{situation}
\label{situation-morphism-thickenings}
Un morphisme d'épaississements $(f, f')$ est donné par un diagramme commutatif
\begin{equation}
\label{equation-morphism-thickenings}
\vcenter{
\xymatrix{
(X, \mathcal{O}_X) \ar[r]_i \ar[d]_f & (X', \mathcal{O}_{X'}) \ar[d]^{f'} \\
(S, \mathcal{O}_S) \ar[r]^t & (S', \mathcal{O}_{S'})
}
}
\end{equation}
d'espaces annelés, dont les flèches horizontales sont des épaississements. Dans cette
situation, on pose
$\mathcal{I} = \Ker(i^\sharp) \subset \mathcal{O}_{X'}$ et
$\mathcal{J} = \Ker(t^\sharp) \subset \mathcal{O}_{S'}$.
Puisque $f = f'$ sur les espaces topologiques sous-jacents, nous identifierons
les foncteurs images réciproques (topologiques) $f^{-1}$ et $(f')^{-1}$.
Remarquons que $(f')^\sharp : f^{-1}\mathcal{O}_{S'} \to \mathcal{O}_{X'}$
induit en particulier une application $f^{-1}\mathcal{J} \to \mathcal{I}$
et donc une application de $\mathcal{O}_{X'}$-modules
$$
(f')^*\mathcal{J} \longrightarrow \mathcal{I}
$$
Si $i$ et $t$ sont des épaississements du premier ordre, alors
$(f')^*\mathcal{J} = f^*\mathcal{J}$ et l'application précédente devient une
application $f^*\mathcal{J} \to \mathcal{I}$.
\end{situation}

\begin{definition}
\label{definition-strict-morphism-thickenings}
Dans la situation \ref{situation-morphism-thickenings}, on dit que $(f, f')$ est un
{\it morphisme strict d'épaississements}
si l'application $(f')^*\mathcal{J} \longrightarrow \mathcal{I}$ est surjective.
\end{definition}

\noindent
Le lemme suivant montre en particulier qu'un morphisme
$(f, f') : (X \subset X') \to (S \subset S')$
d'épaississements de schémas est strict si et seulement si $X = S \times_{S'} X'$.

\begin{lemma}
\label{lemma-strict-morphism-thickenings}
Dans la situation \ref{situation-morphism-thickenings}, le morphisme $(f, f')$
est un morphisme strict d'épaississements si et seulement si
(\ref{equation-morphism-thickenings}) est cartésien dans la catégorie
des espaces annelés.
\end{lemma}

\begin{proof}
Omis.
\end{proof}




\section{Modules sur les épaississements du premier ordre d'espaces annelés}
\label{section-modules-thickenings}

\noindent
Dans cette section, nous présentons quelques préliminaires à la théorie des déformations
des modules. Soit $i : (X, \mathcal{O}_X) \to (X', \mathcal{O}_{X'})$
un épaississement du premier ordre d'espaces annelés. Nous utiliserons librement la notation
introduite dans la section \ref{section-thickenings-spaces} ; en particulier, nous
identifierons les espaces topologiques sous-jacents.
Dans cette section, nous considérons des suites exactes courtes
\begin{equation}
\label{equation-extension}
0 \to \mathcal{K} \to \mathcal{F}' \to \mathcal{F} \to 0
\end{equation}
de $\mathcal{O}_{X'}$-modules, où $\mathcal{F}$, $\mathcal{K}$ sont des
$\mathcal{O}_X$-modules et $\mathcal{F}'$ un $\mathcal{O}_{X'}$-module.
Dans cette situation, on dispose d'une application canonique de $\mathcal{O}_X$-modules
$$
c_{\mathcal{F}'} :
\mathcal{I} \otimes_{\mathcal{O}_X} \mathcal{F}
\longrightarrow
\mathcal{K}
$$
où $\mathcal{I} = \Ker(i^\sharp)$.
En effet, pour des sections locales $f$ de $\mathcal{I}$ et $s$
de $\mathcal{F}$, on pose $c_{\mathcal{F}'}(f \otimes s) = fs'$,
où $s'$ est une section locale de $\mathcal{F}'$ relevant $s$.

\begin{lemma}
\label{lemma-inf-map}
Soit $i : (X, \mathcal{O}_X) \to (X', \mathcal{O}_{X'})$
un épaississement du premier ordre d'espaces annelés. Supposons données des
extensions
$$
0 \to \mathcal{K} \to \mathcal{F}' \to \mathcal{F} \to 0
\quad\text{et}\quad
0 \to \mathcal{L} \to \mathcal{G}' \to \mathcal{G} \to 0
$$
comme dans (\ref{equation-extension}),
ainsi que des applications $\varphi : \mathcal{F} \to \mathcal{G}$ et
$\psi : \mathcal{K} \to \mathcal{L}$.
\begin{enumerate}
\item S'il existe une application de $\mathcal{O}_{X'}$-modules
$\varphi' : \mathcal{F}' \to \mathcal{G}'$ compatible avec $\varphi$
et $\psi$, alors le diagramme
$$
\xymatrix{
\mathcal{I} \otimes_{\mathcal{O}_X} \mathcal{F}
\ar[r]_-{c_{\mathcal{F}'}} \ar[d]_{1 \otimes \varphi} &
\mathcal{K} \ar[d]^\psi \\
\mathcal{I} \otimes_{\mathcal{O}_X} \mathcal{G}
\ar[r]^-{c_{\mathcal{G}'}} &
\mathcal{L}
}
$$
est commutatif.
\item L'ensemble des applications de $\mathcal{O}_{X'}$-modules
$\varphi' : \mathcal{F}' \to \mathcal{G}'$ compatibles avec $\varphi$
et $\psi$ est, s'il est non vide, un espace principal homogène sous
$\Hom_{\mathcal{O}_X}(\mathcal{F}, \mathcal{L})$.
\end{enumerate}
\end{lemma}

\begin{proof}
L'assertion (1) résulte immédiatement de la description des applications.
Pour (2), si $\varphi'$ et $\varphi''$ sont deux applications
$\mathcal{F}' \to \mathcal{G}'$ compatibles avec $\varphi$
et $\psi$, alors $\varphi' - \varphi''$ se factorise sous la forme
$$
\mathcal{F}' \to \mathcal{F} \to \mathcal{L} \to \mathcal{G}'
$$
L'application du milieu provient d'un unique élément de
$\Hom_{\mathcal{O}_X}(\mathcal{F}, \mathcal{L})$, d'après
Modules, lemme \ref{modules-lemma-i-star-equivalence}.
Réciproquement, étant donné un élément $\alpha$ de ce groupe, on peut ajouter la
composée (affichée ci-dessus avec $\alpha$ au milieu)
à $\varphi'$. Nous omettons certains détails.
\end{proof}

\begin{lemma}
\label{lemma-inf-obs-map}
Soit $i : (X, \mathcal{O}_X) \to (X', \mathcal{O}_{X'})$
un épaississement du premier ordre d'espaces annelés. Supposons données des
extensions
$$
0 \to \mathcal{K} \to \mathcal{F}' \to \mathcal{F} \to 0
\quad\text{et}\quad
0 \to \mathcal{L} \to \mathcal{G}' \to \mathcal{G} \to 0
$$
comme dans (\ref{equation-extension}),
ainsi que des applications $\varphi : \mathcal{F} \to \mathcal{G}$ et
$\psi : \mathcal{K} \to \mathcal{L}$. Supposons le diagramme
$$
\xymatrix{
\mathcal{I} \otimes_{\mathcal{O}_X} \mathcal{F}
\ar[r]_-{c_{\mathcal{F}'}} \ar[d]_{1 \otimes \varphi} &
\mathcal{K} \ar[d]^\psi \\
\mathcal{I} \otimes_{\mathcal{O}_X} \mathcal{G}
\ar[r]^-{c_{\mathcal{G}'}} &
\mathcal{L}
}
$$
commutatif. Il existe alors un élément
$$
o(\varphi, \psi) \in
\Ext^1_{\mathcal{O}_X}(\mathcal{F}, \mathcal{L})
$$
dont l'annulation est une condition nécessaire et suffisante pour qu'il existe
une application $\varphi' : \mathcal{F}' \to \mathcal{G}'$ compatible avec
$\varphi$ et $\psi$.
\end{lemma}

\begin{proof}
On peut construire explicitement une extension
$$
0 \to \mathcal{L} \to \mathcal{H} \to \mathcal{F} \to 0
$$
en prenant pour $\mathcal{H}$ la cohomologie au milieu du complexe
$$
\mathcal{K}
\xrightarrow{1, - \psi}
\mathcal{F}' \oplus \mathcal{G}' \xrightarrow{\varphi, 1}
\mathcal{G}
$$
(avec les notations évidentes). Un calcul sur les sections locales,
utilisant l'hypothèse de commutativité du diagramme du lemme,
montre que $\mathcal{H}$ est annulé par $\mathcal{I}$. Par conséquent,
$\mathcal{H}$ définit une classe dans
$$
\Ext^1_{\mathcal{O}_X}(\mathcal{F}, \mathcal{L})
\subset
\Ext^1_{\mathcal{O}_{X'}}(\mathcal{F}, \mathcal{L})
$$
Enfin, la classe de $\mathcal{H}$ est la différence entre la somme amalgamée
de l'extension $\mathcal{F}'$ suivant $\psi$ et l'image réciproque
de l'extension $\mathcal{G}'$ suivant $\varphi$ (calculs omis).
L'annulation de la classe de $\mathcal{H}$ équivaut donc à
l'existence d'un diagramme commutatif
$$
\xymatrix{
0 \ar[r] &
\mathcal{K} \ar[r] \ar[d]_{\psi} &
\mathcal{F}' \ar[r] \ar[d]_{\varphi'} &
\mathcal{F} \ar[r] \ar[d]_\varphi & 0\\
0 \ar[r] &
\mathcal{L} \ar[r] &
\mathcal{G}' \ar[r] &
\mathcal{G} \ar[r] & 0
}
$$
comme voulu.
\end{proof}

\begin{lemma}
\label{lemma-inf-ext}
Soit $i : (X, \mathcal{O}_X) \to (X', \mathcal{O}_{X'})$ un épaississement
du premier ordre d'espaces annelés.
Supposons donnés des $\mathcal{O}_X$-modules $\mathcal{F}$, $\mathcal{K}$
et une application $\mathcal{O}_X$-linéaire
$c : \mathcal{I} \otimes_{\mathcal{O}_X} \mathcal{F} \to \mathcal{K}$.
S'il existe une suite (\ref{equation-extension}) telle que
$c_{\mathcal{F}'} = c$, alors l'ensemble des classes d'isomorphisme de ces
extensions est un espace principal homogène sous
$\Ext^1_{\mathcal{O}_X}(\mathcal{F}, \mathcal{K})$.
\end{lemma}

\begin{proof}
Supposons données des extensions
$$
0 \to \mathcal{K} \to \mathcal{F}'_1 \to \mathcal{F} \to 0
\quad\text{et}\quad
0 \to \mathcal{K} \to \mathcal{F}'_2 \to \mathcal{F} \to 0
$$
telles que $c_{\mathcal{F}'_1} = c_{\mathcal{F}'_2} = c$. Leur différence
(dans le groupe des extensions ; voir
Homology, section \ref{homology-section-extensions})
est une extension
$$
0 \to \mathcal{K} \to \mathcal{E} \to \mathcal{F} \to 0
$$
où $\mathcal{E}$ est annulé par $\mathcal{I}$ (on omet le calcul local).
La suite est donc une extension de $\mathcal{O}_X$-modules ;
voir Modules, lemme \ref{modules-lemma-i-star-equivalence}.
Réciproquement, étant donnée une telle extension $\mathcal{E}$, on peut ajouter l'extension
$\mathcal{E}$ à l'extension de $\mathcal{O}_{X'}$-modules $\mathcal{F}'$ sans
modifier l'application $c_{\mathcal{F}'}$. Nous omettons certains détails.
\end{proof}

\begin{lemma}
\label{lemma-inf-obs-ext}
Soit $i : (X, \mathcal{O}_X) \to (X', \mathcal{O}_{X'})$
un épaississement du premier ordre d'espaces annelés. Supposons donnés des
$\mathcal{O}_X$-modules $\mathcal{F}$, $\mathcal{K}$
et une application $\mathcal{O}_X$-linéaire
$c : \mathcal{I} \otimes_{\mathcal{O}_X} \mathcal{F} \to \mathcal{K}$.
Il existe alors un élément
$$
o(\mathcal{F}, \mathcal{K}, c) \in
\Ext^2_{\mathcal{O}_X}(\mathcal{F}, \mathcal{K})
$$
dont l'annulation est une condition nécessaire et suffisante pour qu'il existe
une suite (\ref{equation-extension}) telle que $c_{\mathcal{F}'} = c$.
\end{lemma}

\begin{proof}
Montrons d'abord que, si $\mathcal{K}$ est un $\mathcal{O}_X$-module injectif,
il existe bien une suite (\ref{equation-extension}) telle que
$c_{\mathcal{F}'} = c$. Pour cela, choisissons un
$\mathcal{O}_{X'}$-module plat $\mathcal{H}'$ et une surjection
$\mathcal{H}' \to \mathcal{F}$
(Modules, lemme \ref{modules-lemma-module-quotient-flat}).
Soit $\mathcal{J} \subset \mathcal{H}'$ le noyau. Comme $\mathcal{H}'$
est plat, on a
$$
\mathcal{I} \otimes_{\mathcal{O}_{X'}} \mathcal{H}' =
\mathcal{I}\mathcal{H}'
\subset \mathcal{J} \subset \mathcal{H}'
$$
Remarquons que l'application
$$
\mathcal{I}\mathcal{H}' =
\mathcal{I} \otimes_{\mathcal{O}_{X'}} \mathcal{H}'
\longrightarrow
\mathcal{I} \otimes_{\mathcal{O}_{X'}} \mathcal{F} =
\mathcal{I} \otimes_{\mathcal{O}_X} \mathcal{F}
$$
annule $\mathcal{I}\mathcal{J}$. En effet, si $f$ est une section locale
de $\mathcal{I}$ et $s$ une section locale de $\mathcal{H}$, alors
$fs$ s'envoie sur $f \otimes \overline{s}$, où $\overline{s}$ est
l'image de $s$ dans $\mathcal{F}$. On obtient ainsi
$$
\xymatrix{
\mathcal{I}\mathcal{H}'/\mathcal{I}\mathcal{J}
\ar@{^{(}->}[r] \ar[d] &
\mathcal{J}/\mathcal{I}\mathcal{J} \ar@{..>}[d]_\gamma \\
\mathcal{I} \otimes_{\mathcal{O}_X} \mathcal{F} \ar[r]^-c &
\mathcal{K}
}
$$
un diagramme de $\mathcal{O}_X$-modules. Si $\mathcal{K}$ est injectif
comme $\mathcal{O}_X$-module, on obtient la flèche pointillée.
Notons $\gamma' : \mathcal{J} \to \mathcal{K}$ la composée
de $\gamma$ avec $\mathcal{J} \to \mathcal{J}/\mathcal{I}\mathcal{J}$.
Un calcul local montre que la somme amalgamée
$$
\xymatrix{
0 \ar[r] &
\mathcal{J} \ar[r] \ar[d]_{\gamma'} &
\mathcal{H}' \ar[r] \ar[d] &
\mathcal{F} \ar[r] \ar@{=}[d] &
0 \\
0 \ar[r] &
\mathcal{K} \ar[r] &
\mathcal{F}' \ar[r] &
\mathcal{F} \ar[r] &
0
}
$$
est une solution au problème posé par le lemme.

\medskip\noindent
Cas général. Choisissons un plongement $\mathcal{K} \subset \mathcal{K}'$,
où $\mathcal{K}'$ est un $\mathcal{O}_X$-module injectif. Soit $\mathcal{Q}$
le quotient, de sorte que l'on ait une suite exacte
$$
0 \to \mathcal{K} \to \mathcal{K}' \to \mathcal{Q} \to 0
$$
Notons
$c' : \mathcal{I} \otimes_{\mathcal{O}_X} \mathcal{F} \to \mathcal{K}'$
la composée. Le paragraphe précédent fournit une suite
$$
0 \to \mathcal{K}' \to \mathcal{E}' \to \mathcal{F} \to 0
$$
comme dans (\ref{equation-extension}), telle que $c_{\mathcal{E}'} = c'$.
Remarquons que la composée de $c'$ avec l'application $\mathcal{K}' \to \mathcal{Q}$
est nulle ; la somme amalgamée de $\mathcal{E}'$ suivant
$\mathcal{K}' \to \mathcal{Q}$ est donc une extension
$$
0 \to \mathcal{Q} \to \mathcal{D}' \to \mathcal{F} \to 0
$$
comme dans (\ref{equation-extension}), telle que $c_{\mathcal{D}'} = 0$.
Cela signifie exactement que $\mathcal{D}'$ est annulé par
$\mathcal{I}$ ; autrement dit, $\mathcal{D}'$ est une extension
de $\mathcal{O}_X$-modules, c'est-à-dire qu'il définit un élément
$$
o(\mathcal{F}, \mathcal{K}, c) \in
\Ext^1_{\mathcal{O}_X}(\mathcal{F}, \mathcal{Q}) =
\Ext^2_{\mathcal{O}_X}(\mathcal{F}, \mathcal{K})
$$
(l'égalité résulte de la suite exacte longue de cohomologie associée
à la suite exacte précédente et de l'annulation des groupes Ext supérieurs
à valeurs dans le module injectif $\mathcal{K}'$). Si
$o(\mathcal{F}, \mathcal{K}, c) = 0$, on peut choisir un scindage
$s : \mathcal{F} \to \mathcal{D}'$ et poser
$$
\mathcal{F}' = \Ker(\mathcal{E}' \to \mathcal{D}'/s(\mathcal{F}))
$$
ce qui donne le diagramme suivant
$$
\xymatrix{
0 \ar[r] &
\mathcal{K} \ar[r] \ar[d] &
\mathcal{F}' \ar[r] \ar[d] &
\mathcal{F} \ar[r] \ar@{=}[d] &
0 \\
0 \ar[r] &
\mathcal{K}' \ar[r] &
\mathcal{E}' \ar[r] &
\mathcal{F} \ar[r] & 0
}
$$
à lignes exactes, qui montre que $c_{\mathcal{F}'} = c$. Réciproquement, si
$\mathcal{F}'$ existe, alors la somme amalgamée de $\mathcal{F}'$ suivant l'application
$\mathcal{K} \to \mathcal{K}'$ est isomorphe à $\mathcal{E}'$, d'après le
lemme \ref{lemma-inf-ext} et l'annulation des groupes Ext supérieurs
à valeurs dans le module injectif $\mathcal{K}'$. On obtient ainsi un diagramme
comme ci-dessus, ce qui implique que $\mathcal{D}'$ est scindée comme extension, c'est-à-dire que
la classe $o(\mathcal{F}, \mathcal{K}, c)$ est nulle.
\end{proof}

\begin{remark}
\label{remark-trivial-thickening}
Soit $(X, \mathcal{O}_X)$ un espace annelé. Un épaississement du premier ordre
$i : (X, \mathcal{O}_X) \to (X', \mathcal{O}_{X'})$ est dit
{\it trivial} s'il existe un morphisme d'espaces annelés
$\pi : (X', \mathcal{O}_{X'}) \to (X, \mathcal{O}_X)$ qui soit un
inverse à gauche de $i$. Le choix d'un tel morphisme
$\pi$ est appelé une {\it trivialisation} de l'épaississement du premier ordre.
Étant donné $\pi$, on obtient un scindage
\begin{equation}
\label{equation-splitting}
\mathcal{O}_{X'} = \mathcal{O}_X \oplus \mathcal{I}
\end{equation}
comme faisceaux d'algèbres sur $X$, en utilisant $\pi^\sharp$ pour scinder la surjection
$\mathcal{O}_{X'} \to \mathcal{O}_X$. Réciproquement, un tel scindage détermine
un morphisme $\pi$. La catégorie des épaississements du premier ordre trivialisés de
$(X, \mathcal{O}_X)$ est équivalente à la catégorie des
$\mathcal{O}_X$-modules.
\end{remark}

\begin{remark}
\label{remark-trivial-extension}
Soit $i : (X, \mathcal{O}_X) \to (X', \mathcal{O}_{X'})$
un épaississement trivial du premier ordre d'espaces annelés,
et soit $\pi : (X', \mathcal{O}_{X'}) \to (X, \mathcal{O}_X)$
une trivialisation. Pour tout triplet
$(\mathcal{F}, \mathcal{K}, c)$ formé d'une paire de
$\mathcal{O}_X$-modules et d'une application
$c : \mathcal{I} \otimes_{\mathcal{O}_X} \mathcal{F} \to \mathcal{K}$,
on peut poser
$$
\mathcal{F}'_{c, triv} = \mathcal{F} \oplus \mathcal{K}
$$
et utiliser le scindage (\ref{equation-splitting}) associé à $\pi$
ainsi que l'application $c$ pour définir la structure de $\mathcal{O}_{X'}$-module
et obtenir une extension (\ref{equation-extension}). On appellera
$\mathcal{F}'_{c, triv}$ l'{\it extension triviale} de $\mathcal{F}$
par $\mathcal{K}$ correspondant
à $c$ et à la trivialisation $\pi$. Pour toute extension
$\mathcal{F}'$ comme dans (\ref{equation-extension}), on peut utiliser
$\pi^\sharp : \mathcal{O}_X \to \mathcal{O}_{X'}$ pour considérer $\mathcal{F}'$
comme une extension de $\mathcal{O}_X$-modules, et donc comme une classe $\xi_{\mathcal{F}'}$
dans $\Ext^1_{\mathcal{O}_X}(\mathcal{F}, \mathcal{K})$.
Le lemme \ref{lemma-inf-ext} assure que
$\mathcal{F}' \mapsto \xi_{\mathcal{F}'}$
induit une bijection
$$
\left\{
\begin{matrix}
\text{isomorphism classes of extensions}\\
\mathcal{F}'\text{ comme dans (\ref{equation-extension}) avec }c = c_{\mathcal{F}'}
\end{matrix}
\right\}
\longrightarrow
\Ext^1_{\mathcal{O}_X}(\mathcal{F}, \mathcal{K})
$$
De plus, l'extension triviale $\mathcal{F}'_{c, triv}$ s'envoie sur la classe nulle.
\end{remark}

\begin{remark}
\label{remark-extension-functorial}
Soit $(X, \mathcal{O}_X)$ un espace annelé. Soient
$(X, \mathcal{O}_X) \to (X'_i, \mathcal{O}_{X'_i})$, $i = 1, 2$,
des épaississements du premier ordre, de faisceaux d'idéaux $\mathcal{I}_i$.
Soit $h : (X'_1, \mathcal{O}_{X'_1}) \to (X'_2, \mathcal{O}_{X'_2})$
un morphisme d'épaississements du premier ordre de $(X, \mathcal{O}_X)$.
On a le diagramme
$$
\xymatrix{
& (X, \mathcal{O}_X) \ar[ld] \ar[rd] & \\
(X'_1, \mathcal{O}_{X'_1}) \ar[rr]^h & & 
(X'_2, \mathcal{O}_{X'_2})
}
$$
Remarquons que $h^\sharp : \mathcal{O}_{X'_2} \to \mathcal{O}_{X'_1}$
induit en particulier une application de $\mathcal{O}_X$-modules
$\mathcal{I}_2 \to \mathcal{I}_1$.
Soit $\mathcal{F}$ un
$\mathcal{O}_X$-module. Soit $(\mathcal{K}_i, c_i)$, pour $i = 1, 2$, une paire
formée d'un $\mathcal{O}_X$-module $\mathcal{K}_i$ et d'une application
$c_i : \mathcal{I}_i \otimes_{\mathcal{O}_X} \mathcal{F} \to
\mathcal{K}_i$. Supposons en outre donnée une application
de $\mathcal{O}_X$-modules $\mathcal{K}_2 \to \mathcal{K}_1$
telle que le diagramme
$$
\xymatrix{
\mathcal{I}_2 \otimes_{\mathcal{O}_X} \mathcal{F}
\ar[r]_-{c_2} \ar[d] &
\mathcal{K}_2 \ar[d] \\
\mathcal{I}_1 \otimes_{\mathcal{O}_X} \mathcal{F}
\ar[r]^-{c_1} &
\mathcal{K}_1
}
$$
soit commutatif. On dispose alors d'une fonctorialité canonique
$$
\left\{
\begin{matrix}
\mathcal{F}'_2\text{ comme dans (\ref{equation-extension}) avec }\\
c_2 = c_{\mathcal{F}'_2}\text{ et }\mathcal{K} = \mathcal{K}_2
\end{matrix}
\right\}
\longrightarrow
\left\{
\begin{matrix}
\mathcal{F}'_1\text{ comme dans (\ref{equation-extension}) avec }\\
c_1 = c_{\mathcal{F}'_1}\text{ et }\mathcal{K} = \mathcal{K}_1
\end{matrix}
\right\}
$$
En effet, en considérant tous les faisceaux $\mathcal{O}_X$, $\mathcal{O}_{X'_i}$,
$\mathcal{F}$, $\mathcal{K}_i$, etc., comme des faisceaux sur $X$, on associe
à $\mathcal{F}'_2$ le faisceau $\mathcal{F}'_1$ défini comme la
somme amalgamée, c'est-à-dire s'insérant dans le diagramme d'extensions suivant :
$$
\xymatrix{
0 \ar[r] &
\mathcal{K}_2 \ar[r] \ar[d] &
\mathcal{F}'_2 \ar[r] \ar[d] &
\mathcal{F} \ar@{=}[d] \ar[r] & 0 \\
0 \ar[r] &
\mathcal{K}_1 \ar[r] &
\mathcal{F}'_1 \ar[r] &
\mathcal{F} \ar[r] & 0
}
$$
Nous omettons la construction de la structure de $\mathcal{O}_{X'_1}$-module
sur la somme amalgamée (elle utilise la commutativité du diagramme
faisant intervenir $c_1$ et $c_2$).
\end{remark}

\begin{remark}
\label{remark-trivial-extension-functorial}
Soient $(X, \mathcal{O}_X)$, $(X, \mathcal{O}_X) \to (X'_i, \mathcal{O}_{X'_i})$,
$\mathcal{I}_i$, et
$h : (X'_1, \mathcal{O}_{X'_1}) \to (X'_2, \mathcal{O}_{X'_2})$
comme dans la remarque \ref{remark-extension-functorial}. Supposons données des
trivialisations $\pi_i : X'_i \to X$ telles que
$\pi_1 = h \circ \pi_2$. Autrement dit, supposons que $h$ soit un morphisme
d'épaississements du premier ordre trivialisés de $(X, \mathcal{O}_X)$. Soit
$(\mathcal{K}_i, c_i)$, pour $i = 1, 2$, une paire formée d'un
$\mathcal{O}_X$-module $\mathcal{K}_i$ et d'une application
$c_i : \mathcal{I}_i \otimes_{\mathcal{O}_X} \mathcal{F} \to
\mathcal{K}_i$. Supposons en outre donnée une application
de $\mathcal{O}_X$-modules $\mathcal{K}_2 \to \mathcal{K}_1$
telle que le diagramme
$$
\xymatrix{
\mathcal{I}_2 \otimes_{\mathcal{O}_X} \mathcal{F}
\ar[r]_-{c_2} \ar[d] &
\mathcal{K}_2 \ar[d] \\
\mathcal{I}_1 \otimes_{\mathcal{O}_X} \mathcal{F}
\ar[r]^-{c_1} &
\mathcal{K}_1
}
$$
soit commutatif. Dans cette situation, la construction de la
remarque \ref{remark-trivial-extension} induit
un diagramme commutatif
$$
\xymatrix{
\{\mathcal{F}'_2\text{ comme dans (\ref{equation-extension}) avec }
c_2 = c_{\mathcal{F}'_2}\text{ et }\mathcal{K} = \mathcal{K}_2\}
\ar[d] \ar[rr] & &
\Ext^1_{\mathcal{O}_X}(\mathcal{F}, \mathcal{K}_2) \ar[d] \\
\{\mathcal{F}'_1\text{ comme dans (\ref{equation-extension}) avec }
c_1 = c_{\mathcal{F}'_1}\text{ et }\mathcal{K} = \mathcal{K}_1\}
\ar[rr] & &
\Ext^1_{\mathcal{O}_X}(\mathcal{F}, \mathcal{K}_1)
}
$$
où l'application verticale de droite est donnée par la fonctorialité de $\Ext$
et l'application $\mathcal{K}_2 \to \mathcal{K}_1$, tandis que l'application verticale de gauche
est celle de la remarque \ref{remark-extension-functorial}.
\end{remark}

\begin{remark}
\label{remark-short-exact-sequence-thickenings}
Soit $(X, \mathcal{O}_X)$ un espace annelé. On dit qu'une suite de morphismes
d'épaississements du premier ordre
$$
(X'_1, \mathcal{O}_{X'_1}) \to
(X'_2, \mathcal{O}_{X'_2}) \to
(X'_3, \mathcal{O}_{X'_3})
$$
de $(X, \mathcal{O}_X)$ est un {\it complexe}
si les applications correspondantes entre
les faisceaux d'idéaux $\mathcal{I}_i$
forment un complexe de $\mathcal{O}_X$-modules
$\mathcal{I}_3 \to \mathcal{I}_2 \to \mathcal{I}_1$
(c'est-à-dire si la composée est nulle). Dans ce cas, la composée
$(X'_1, \mathcal{O}_{X'_1}) \to (X_3', \mathcal{O}_{X'_3})$ se factorise par
$(X, \mathcal{O}_X) \to (X'_3, \mathcal{O}_{X'_3})$ ; autrement dit,
l'épaississement du premier ordre $(X'_1, \mathcal{O}_{X'_1})$ de
$(X, \mathcal{O}_X)$ est trivial et muni d'une
trivialisation canonique
$\pi : (X'_1, \mathcal{O}_{X'_1}) \to (X, \mathcal{O}_X)$.

\medskip\noindent
On dit qu'une suite de morphismes d'épaississements du premier ordre
$$
(X'_1, \mathcal{O}_{X'_1}) \to
(X'_2, \mathcal{O}_{X'_2}) \to
(X'_3, \mathcal{O}_{X'_3})
$$
de $(X, \mathcal{O}_X)$ est une {\it suite exacte courte} si les
applications correspondantes entre les faisceaux d'idéaux forment une suite exacte courte
$$
0 \to \mathcal{I}_3 \to \mathcal{I}_2 \to \mathcal{I}_1 \to 0
$$
de $\mathcal{O}_X$-modules.
\end{remark}

\begin{remark}
\label{remark-complex-thickenings-and-ses-modules}
Soit $(X, \mathcal{O}_X)$ un espace annelé. Soit $\mathcal{F}$ un
$\mathcal{O}_X$-module. Soit
$$
(X'_1, \mathcal{O}_{X'_1}) \to
(X'_2, \mathcal{O}_{X'_2}) \to
(X'_3, \mathcal{O}_{X'_3})
$$
un complexe d'épaississements du premier ordre de $(X, \mathcal{O}_X)$, voir
Remarque \ref{remark-short-exact-sequence-thickenings}.
Soient $(\mathcal{K}_i, c_i)$, $i = 1, 2, 3$, des couples formés d'un
$\mathcal{O}_X$-module $\mathcal{K}_i$ et d'une application
$c_i : \mathcal{I}_i \otimes_{\mathcal{O}_X} \mathcal{F} \to
\mathcal{K}_i$. Supposons donnée une suite exacte courte
de $\mathcal{O}_X$-modules
$$
0 \to \mathcal{K}_3 \to \mathcal{K}_2 \to \mathcal{K}_1 \to 0
$$
telle que
$$
\vcenter{
\xymatrix{
\mathcal{I}_2 \otimes_{\mathcal{O}_X} \mathcal{F}
\ar[r]_-{c_2} \ar[d] &
\mathcal{K}_2 \ar[d] \\
\mathcal{I}_1 \otimes_{\mathcal{O}_X} \mathcal{F}
\ar[r]^-{c_1} &
\mathcal{K}_1
}
}
\quad\text{et}\quad
\vcenter{
\xymatrix{
\mathcal{I}_3 \otimes_{\mathcal{O}_X} \mathcal{F}
\ar[r]_-{c_3} \ar[d] &
\mathcal{K}_3 \ar[d] \\
\mathcal{I}_2 \otimes_{\mathcal{O}_X} \mathcal{F}
\ar[r]^-{c_2} &
\mathcal{K}_2
}
}
$$
soient commutatifs. Supposons enfin donnée une extension
$$
0 \to \mathcal{K}_2 \to \mathcal{F}'_2 \to \mathcal{F} \to 0
$$
comme dans (\ref{equation-extension}), avec $\mathcal{K} = \mathcal{K}_2$,
de $\mathcal{O}_{X'_2}$-modules telle que $c_{\mathcal{F}'_2} = c_2$.
Dans cette situation, nous pouvons appliquer la fonctorialité de la
Remarque \ref{remark-extension-functorial} pour obtenir une extension
$\mathcal{F}'_1$ sur $X'_1$ (nous décrirons ci-dessous $\mathcal{F}'_1$
dans ce cas particulier). Par la
Remarque \ref{remark-trivial-extension},
en utilisant le scindage canonique
$\pi : (X'_1, \mathcal{O}_{X'_1}) \to (X, \mathcal{O}_X)$ de la
Remarque \ref{remark-short-exact-sequence-thickenings},
nous obtenons
$\xi_{\mathcal{F}'_1} \in
\Ext^1_{\mathcal{O}_X}(\mathcal{F}, \mathcal{K}_1)$.
Enfin, nous avons l'obstruction
$$
o(\mathcal{F}, \mathcal{K}_3, c_3) \in
\Ext^2_{\mathcal{O}_X}(\mathcal{F}, \mathcal{K}_3)
$$
du Lemme \ref{lemma-inf-obs-ext}.
Dans cette situation, nous {\bf affirmons} que l'application canonique
$$
\partial :
\Ext^1_{\mathcal{O}_X}(\mathcal{F}, \mathcal{K}_1)
\longrightarrow
\Ext^2_{\mathcal{O}_X}(\mathcal{F}, \mathcal{K}_3)
$$
provenant de la suite exacte courte
$0 \to \mathcal{K}_3 \to \mathcal{K}_2 \to \mathcal{K}_1 \to 0$
envoie $\xi_{\mathcal{F}'_1}$
sur la classe d'obstruction $o(\mathcal{F}, \mathcal{K}_3, c_3)$.

\medskip\noindent
Pour démontrer cette affirmation, choisissons un plongement
$j : \mathcal{K}_3 \to \mathcal{K}$, où $\mathcal{K}$ est un
$\mathcal{O}_X$-module injectif.
Nous pouvons relever $j$ en une application $j' : \mathcal{K}_2 \to \mathcal{K}$.
Posons $\mathcal{E}'_2 = j'_*\mathcal{F}'_2$, égal à la somme amalgamée
de $\mathcal{F}'_2$ par $j'$, de sorte que $c_{\mathcal{E}'_2} = j' \circ c_2$.
Voici le diagramme :
$$
\xymatrix{
0 \ar[r] &
\mathcal{K}_2 \ar[r] \ar[d]_{j'} &
\mathcal{F}'_2 \ar[r] \ar[d] &
\mathcal{F} \ar[r] \ar[d] & 0 \\
0 \ar[r] &
\mathcal{K} \ar[r] &
\mathcal{E}'_2 \ar[r] &
\mathcal{F} \ar[r] & 0
}
$$
Posons $\mathcal{E}'_3 = \mathcal{E}'_2$, mais considérons-le comme un
$\mathcal{O}_{X'_3}$-module au moyen de $\mathcal{O}_{X'_3} \to \mathcal{O}_{X'_2}$.
Alors $c_{\mathcal{E}'_3} = j \circ c_3$.
La démonstration du Lemme \ref{lemma-inf-obs-ext} construit
$o(\mathcal{F}, \mathcal{K}_3, c_3)$
comme le cobord de la classe de l'extension de $\mathcal{O}_X$-modules
$$
0 \to
\mathcal{K}/\mathcal{K}_3 \to
\mathcal{E}'_3/\mathcal{K}_3 \to
\mathcal{F} \to 0
$$
D'autre part, remarquons que
$\mathcal{F}'_1 = \mathcal{F}'_2/\mathcal{K}_3$ ; ainsi, la classe
$\xi_{\mathcal{F}'_1}$ est la classe de l'extension
$$
0 \to \mathcal{K}_2/\mathcal{K}_3 \to \mathcal{F}'_2/\mathcal{K}_3
\to \mathcal{F} \to 0
$$
considérée comme une suite de $\mathcal{O}_X$-modules au moyen de
$\pi^\sharp$, où
$\pi : (X'_1, \mathcal{O}_{X'_1}) \to (X, \mathcal{O}_X)$
est le scindage canonique.
Ainsi, l'affirmation résulte finalement de l'existence du diagramme
commutatif
$$
\xymatrix{
0 \ar[r] &
\mathcal{K}_2/\mathcal{K}_3 \ar[r] \ar[d] &
\mathcal{F}'_2/\mathcal{K}_3 \ar[r] \ar[d] &
\mathcal{F} \ar[r] \ar[d] & 0 \\
0 \ar[r] &
\mathcal{K}/\mathcal{K}_3 \ar[r] &
\mathcal{E}'_3/\mathcal{K}_3 \ar[r] &
\mathcal{F} \ar[r] & 0
}
$$
qui est $\mathcal{O}_X$-linéaire (pour les structures de
$\mathcal{O}_X$-modules données ci-dessus).
\end{remark}








\section{Déformations infinitésimales de modules sur des espaces annelés}
\sectionmark{Déformations de modules sur des espaces annelés}
\label{section-deformation-modules}

\noindent
Soit $i : (X, \mathcal{O}_X) \to (X', \mathcal{O}_{X'})$ un épaississement
du premier ordre d'espaces annelés. Nous employons librement les notations introduites à la
section \ref{section-thickenings-spaces}.
Soit $\mathcal{F}'$ un $\mathcal{O}_{X'}$-module
et posons $\mathcal{F} = i^*\mathcal{F}'$.
Dans cette situation, nous avons une suite exacte courte
$$
0 \to \mathcal{I}\mathcal{F}' \to \mathcal{F}' \to \mathcal{F} \to 0
$$
de $\mathcal{O}_{X'}$-modules. Puisque $\mathcal{I}^2 = 0$, la structure de
$\mathcal{O}_{X'}$-module sur $\mathcal{I}\mathcal{F}'$
provient d'une unique structure de $\mathcal{O}_X$-module.
Ainsi, la suite ci-dessus est une extension comme dans (\ref{equation-extension}).
Dans le cas particulier où $\mathcal{F}' = \mathcal{O}_{X'}$, nous avons
$i^*\mathcal{O}_{X'} = \mathcal{O}_X$ et
$\mathcal{I}\mathcal{O}_{X'} = \mathcal{I}$, et nous retrouvons la
suite des faisceaux structuraux
$$
0 \to \mathcal{I} \to \mathcal{O}_{X'} \to \mathcal{O}_X \to 0
$$

\begin{lemma}
\label{lemma-inf-map-special}
Soit $i : (X, \mathcal{O}_X) \to (X', \mathcal{O}_{X'})$
un épaississement du premier ordre d'espaces annelés.
Soient $\mathcal{F}'$, $\mathcal{G}'$ des $\mathcal{O}_{X'}$-modules.
Posons $\mathcal{F} = i^*\mathcal{F}'$ et $\mathcal{G} = i^*\mathcal{G}'$.
Soit $\varphi : \mathcal{F} \to \mathcal{G}$ une application
$\mathcal{O}_X$-linéaire.
L'ensemble des relèvements de $\varphi$ en une application
$\mathcal{O}_{X'}$-linéaire
$\varphi' : \mathcal{F}' \to \mathcal{G}'$ est, s'il est non vide, un espace
principal homogène sous
$\Hom_{\mathcal{O}_X}(\mathcal{F}, \mathcal{I}\mathcal{G}')$.
\end{lemma}

\begin{proof}
C'est un cas particulier du Lemme \ref{lemma-inf-map}, mais nous en donnons
aussi une démonstration directe. Nous avons des suites exactes courtes de modules
$$
0 \to \mathcal{I} \to \mathcal{O}_{X'} \to \mathcal{O}_X \to 0
\quad\text{et}\quad
0 \to \mathcal{I}\mathcal{G}' \to \mathcal{G}' \to \mathcal{G} \to 0
$$
et de même pour $\mathcal{F}'$.
Comme $\mathcal{I}$ est de carré nul, les structures de
$\mathcal{O}_{X'}$-modules sur $\mathcal{I}$ et
$\mathcal{I}\mathcal{G}'$ proviennent d'uniques structures de
$\mathcal{O}_X$-modules. Il s'ensuit que
$$
\Hom_{\mathcal{O}_{X'}}(\mathcal{F}', \mathcal{I}\mathcal{G}') =
\Hom_{\mathcal{O}_X}(\mathcal{F}, \mathcal{I}\mathcal{G}')
\quad\text{et}\quad
\Hom_{\mathcal{O}_{X'}}(\mathcal{F}', \mathcal{G}) =
\Hom_{\mathcal{O}_X}(\mathcal{F}, \mathcal{G})
$$
Le lemme résulte alors de la suite exacte
$$
0 \to \Hom_{\mathcal{O}_{X'}}(\mathcal{F}', \mathcal{I}\mathcal{G}') \to
\Hom_{\mathcal{O}_{X'}}(\mathcal{F}', \mathcal{G}') \to
\Hom_{\mathcal{O}_{X'}}(\mathcal{F}', \mathcal{G})
$$
voir Homology, lemme \ref{homology-lemma-check-exactness}.
\end{proof}

\begin{lemma}
\label{lemma-deform-module}
Soit $(f, f')$ un morphisme d'épaississements du premier ordre d'espaces annelés
comme dans la Situation \ref{situation-morphism-thickenings}.
Soit $\mathcal{F}'$ un $\mathcal{O}_{X'}$-module
et posons $\mathcal{F} = i^*\mathcal{F}'$.
Supposons que $\mathcal{F}$ soit plat sur $S$
et que $(f, f')$ soit un morphisme strict d'épaississements
(Définition \ref{definition-strict-morphism-thickenings}).
Alors les assertions suivantes sont équivalentes :
\begin{enumerate}
\item $\mathcal{F}'$ est plat sur $S'$ ;
\item l'application canonique
$f^*\mathcal{J} \otimes_{\mathcal{O}_X} \mathcal{F} \to
\mathcal{I}\mathcal{F}'$
est un isomorphisme.
\end{enumerate}
De plus, dans ce cas, les applications
$$
f^*\mathcal{J} \otimes_{\mathcal{O}_X} \mathcal{F} \to
\mathcal{I} \otimes_{\mathcal{O}_X} \mathcal{F} \to
\mathcal{I}\mathcal{F}'
$$
sont des isomorphismes.
\end{lemma}

\begin{proof}
L'application $f^*\mathcal{J} \to \mathcal{I}$ est surjective,
car $(f, f')$ est un morphisme strict d'épaississements.
La dernière assertion résulte donc de (2).

\medskip\noindent
Démontrons l'équivalence de (1) et (2). Nous pouvons vérifier ces conditions
sur les germes. Soit $x \in X \subset X'$
un point d'image $s = f(x) \in S \subset S'$.
Posons $A' = \mathcal{O}_{S', s}$, $B' = \mathcal{O}_{X', x}$,
$A = \mathcal{O}_{S, s}$ et $B = \mathcal{O}_{X, x}$.
Alors $A = A'/J$ et $B = B'/I$ pour certains idéaux de carré nul.
Puisque $(f, f')$ est un morphisme strict d'épaississements, nous avons
$I = JB'$.
Soient $M' = \mathcal{F}'_x$ et $M = \mathcal{F}_x$.
Alors $M'$ est un $B'$-module et $M$ est un $B$-module.
Comme $\mathcal{F} = i^*\mathcal{F}'$, le noyau de la
surjection $M' \to M$ est $IM' = JM'$. Nous avons donc une suite exacte
courte
$$
0 \to JM' \to M' \to M \to 0
$$
En utilisant
Faisceaux, lemme \ref{sheaves-lemma-stalk-pullback-modules}
et
Modules, lemme \ref{modules-lemma-stalk-tensor-product}
pour identifier les germes des images réciproques et des produits tensoriels, nous voyons
que le germe en $x$ de l'application canonique de la lemme est l'application
$$
(J \otimes_A B) \otimes_B M = J \otimes_A M = J \otimes_{A'} M'
\longrightarrow JM'
$$
L'hypothèse que $\mathcal{F}$ est plat sur $S$ signifie que
$M$ est un $A$-module plat.

\medskip\noindent
Supposons (1). La platitude implique
$\text{Tor}_1^{A'}(M', A) = 0$ d'après
Algèbre, lemme \ref{algebra-lemma-characterize-flat}.
Cela signifie que $J \otimes_{A'} M' \to M'$ est injective d'après
Algèbre, remarque \ref{algebra-remark-Tor-ring-mod-ideal}.
Ainsi, $J \otimes_A M \to JM'$ est un isomorphisme.

\medskip\noindent
Supposons (2). Alors $J \otimes_{A'} M' \to M'$ est injective. Donc
$\text{Tor}_1^{A'}(M', A) = 0$ d'après
Algèbre, remarque \ref{algebra-remark-Tor-ring-mod-ideal}.
Ainsi, $M'$ est plat sur $A'$ d'après
Algèbre, lemme \ref{algebra-lemma-what-does-it-mean}.
\end{proof}

\begin{lemma}
\label{lemma-inf-map-rel}
Soit $(f, f')$ un morphisme d'épaississements du premier ordre comme dans la
Situation \ref{situation-morphism-thickenings}.
Soient $\mathcal{F}'$, $\mathcal{G}'$ des $\mathcal{O}_{X'}$-modules et posons
$\mathcal{F} = i^*\mathcal{F}'$ et $\mathcal{G} = i^*\mathcal{G}'$.
Soit $\varphi : \mathcal{F} \to \mathcal{G}$ une application
$\mathcal{O}_X$-linéaire.
Supposons que $\mathcal{G}'$ soit plat sur $S'$ et que
$(f, f')$ soit un morphisme strict d'épaississements.
L'ensemble des relèvements de $\varphi$ en une application
$\mathcal{O}_{X'}$-linéaire
$\varphi' : \mathcal{F}' \to \mathcal{G}'$ est, s'il est non vide, un espace
principal homogène sous
$$
\Hom_{\mathcal{O}_X}(\mathcal{F},
\mathcal{G} \otimes_{\mathcal{O}_X} f^*\mathcal{J})
$$
\end{lemma}

\begin{proof}
Il suffit de combiner les Lemmes \ref{lemma-inf-map-special} et
\ref{lemma-deform-module}.
\end{proof}

\begin{lemma}
\label{lemma-inf-obs-map-special}
Soit $i : (X, \mathcal{O}_X) \to (X', \mathcal{O}_{X'})$
un épaississement du premier ordre d'espaces annelés.
Soient $\mathcal{F}'$, $\mathcal{G}'$ des $\mathcal{O}_{X'}$-modules et posons
$\mathcal{F} = i^*\mathcal{F}'$ et $\mathcal{G} = i^*\mathcal{G}'$.
Soit $\varphi : \mathcal{F} \to \mathcal{G}$ une application
$\mathcal{O}_X$-linéaire.
Il existe un élément
$$
o(\varphi) \in
\Ext^1_{\mathcal{O}_X}(Li^*\mathcal{F}',
\mathcal{I}\mathcal{G}')
$$
dont l'annulation est une condition nécessaire et suffisante pour
l'existence d'un relèvement de $\varphi$ en une application
$\mathcal{O}_{X'}$-linéaire
$\varphi' : \mathcal{F}' \to \mathcal{G}'$.
\end{lemma}

\begin{proof}
Il ressort de la démonstration du Lemme \ref{lemma-inf-map-special} que
l'annulation du cobord de $\varphi$ par l'application
$$
\Hom_{\mathcal{O}_X}(\mathcal{F}, \mathcal{G}) =
\Hom_{\mathcal{O}_{X'}}(\mathcal{F}', \mathcal{G}) \longrightarrow
\Ext^1_{\mathcal{O}_{X'}}(\mathcal{F}', \mathcal{I}\mathcal{G}')
$$
est une condition nécessaire et suffisante pour l'existence d'un relèvement. Nous
concluons puisque
$$
\Ext^1_{\mathcal{O}_{X'}}(\mathcal{F}', \mathcal{I}\mathcal{G}') =
\Ext^1_{\mathcal{O}_X}(Li^*\mathcal{F}', \mathcal{I}\mathcal{G}')
$$
par l'adjonction de $i_* = Ri_*$ et $Li^*$ dans la catégorie dérivée
(Cohomologie, lemme \ref{cohomology-lemma-adjoint}).
\end{proof}

\begin{lemma}
\label{lemma-inf-obs-map-rel}
Soit $(f, f')$ un morphisme d'épaississements du
premier ordre comme dans la Situation \ref{situation-morphism-thickenings}.
Soient $\mathcal{F}'$, $\mathcal{G}'$ des $\mathcal{O}_{X'}$-modules et posons
$\mathcal{F} = i^*\mathcal{F}'$ et $\mathcal{G} = i^*\mathcal{G}'$.
Soit $\varphi : \mathcal{F} \to \mathcal{G}$ une application
$\mathcal{O}_X$-linéaire.
Supposons que $\mathcal{F}'$ et $\mathcal{G}'$ soient plats sur $S'$ et
que $(f, f')$ soit un morphisme strict d'épaississements. Il existe un élément
$$
o(\varphi) \in  \Ext^1_{\mathcal{O}_X}(\mathcal{F},
\mathcal{G} \otimes_{\mathcal{O}_X} f^*\mathcal{J})
$$
dont l'annulation est une condition nécessaire et suffisante pour
l'existence d'un relèvement de $\varphi$ en une application
$\mathcal{O}_{X'}$-linéaire
$\varphi' : \mathcal{F}' \to \mathcal{G}'$.
\end{lemma}

\begin{proof}[Première démonstration]
Cela résulte du Lemme \ref{lemma-inf-obs-map-special},
car nous affirmons que, sous les hypothèses de la lemme, nous avons
$$
\Ext^1_{\mathcal{O}_X}(Li^*\mathcal{F}',
\mathcal{I}\mathcal{G}') =
\Ext^1_{\mathcal{O}_X}(\mathcal{F},
\mathcal{G} \otimes_{\mathcal{O}_X} f^*\mathcal{J})
$$
En effet, nous avons
$\mathcal{I}\mathcal{G}' =
\mathcal{G} \otimes_{\mathcal{O}_X} f^*\mathcal{J}$
d'après le Lemme \ref{lemma-deform-module}.
D'autre part, remarquons que
$$
H^{-1}(Li^*\mathcal{F}') =
\text{Tor}_1^{\mathcal{O}_{X'}}(\mathcal{F}', \mathcal{O}_X)
$$
(nous omettons le calcul local). En utilisant la suite exacte courte
$$
0 \to \mathcal{I} \to \mathcal{O}_{X'} \to \mathcal{O}_X \to 0
$$
nous voyons que ce $\text{Tor}_1$ est calculé par le noyau de l'application
$\mathcal{I} \otimes_{\mathcal{O}_X} \mathcal{F} \to \mathcal{I}\mathcal{F}'$,
qui est nul d'après la dernière assertion du Lemme \ref{lemma-deform-module}.
Ainsi, $\tau_{\geq -1}Li^*\mathcal{F}' = \mathcal{F}$.
D'autre part, nous avons
$$
\Ext^1_{\mathcal{O}_X}(Li^*\mathcal{F}',
\mathcal{I}\mathcal{G}') =
\Ext^1_{\mathcal{O}_X}(\tau_{\geq -1}Li^*\mathcal{F}',
\mathcal{I}\mathcal{G}')
$$
par le dual de
Catégories dérivées, lemme \ref{derived-lemma-negative-vanishing}.
\end{proof}

\begin{proof}[Seconde démonstration]
Nous pouvons appliquer le Lemme \ref{lemma-inf-obs-map} comme suit. Remarquons que
$\mathcal{K} = \mathcal{I} \otimes_{\mathcal{O}_X} \mathcal{F}$ et
$\mathcal{L} = \mathcal{I} \otimes_{\mathcal{O}_X} \mathcal{G}$
d'après le Lemme \ref{lemma-deform-module}, que
$c_{\mathcal{F}'} = 1 \otimes 1$ et $c_{\mathcal{G}'} = 1 \otimes 1$,
et qu'en prenant $\psi = 1 \otimes \varphi$, le diagramme de la lemme
est commutatif. Ainsi, $o(\varphi) = o(\varphi, 1 \otimes \varphi)$
convient.
\end{proof}

\begin{lemma}
\label{lemma-inf-ext-rel}
Soit $(f, f')$ un morphisme d'épaississements du premier ordre comme dans la
Situation \ref{situation-morphism-thickenings}.
Soit $\mathcal{F}$ un $\mathcal{O}_X$-module.
Supposons que $(f, f')$ soit un morphisme strict d'épaississements et que
$\mathcal{F}$ soit plat sur $S$. S'il existe un couple
$(\mathcal{F}', \alpha)$ formé d'un
$\mathcal{O}_{X'}$-module $\mathcal{F}'$ plat sur $S'$ et d'un isomorphisme
$\alpha : i^*\mathcal{F}' \to \mathcal{F}$, alors l'ensemble des
classes d'isomorphisme de tels couples est un espace principal homogène
sous
$\Ext^1_{\mathcal{O}_X}(
\mathcal{F}, \mathcal{I} \otimes_{\mathcal{O}_X} \mathcal{F})$.
\end{lemma}

\begin{proof}
Si nous supposons qu'il existe un tel module, alors l'application canonique
$$
f^*\mathcal{J} \otimes_{\mathcal{O}_X} \mathcal{F} \to
\mathcal{I} \otimes_{\mathcal{O}_X} \mathcal{F}
$$
est un isomorphisme d'après le Lemme \ref{lemma-deform-module}. Appliquons le
Lemme \ref{lemma-inf-ext} avec $\mathcal{K} = 
\mathcal{I} \otimes_{\mathcal{O}_X} \mathcal{F}$
et $c = 1$. D'après le Lemme \ref{lemma-deform-module}, les extensions
$\mathcal{F}'$ correspondantes sont toutes plates sur $S'$.
\end{proof}

\begin{lemma}
\label{lemma-inf-obs-ext-rel}
Soit $(f, f')$ un morphisme d'épaississements du premier ordre comme dans la
Situation \ref{situation-morphism-thickenings}.
Soit $\mathcal{F}$ un $\mathcal{O}_X$-module. Supposons que
$(f, f')$ soit un morphisme strict d'épaississements
et que $\mathcal{F}$ soit plat sur $S$. Il existe un
$\mathcal{O}_{X'}$-module $\mathcal{F}'$ plat sur $S'$ tel que
$i^*\mathcal{F}' \cong \mathcal{F}$, si et seulement si
\begin{enumerate}
\item l'application canonique $
f^*\mathcal{J} \otimes_{\mathcal{O}_X} \mathcal{F} \to
\mathcal{I} \otimes_{\mathcal{O}_X} \mathcal{F}$
est un isomorphisme ;
\item la classe
$o(\mathcal{F}, \mathcal{I} \otimes_{\mathcal{O}_X} \mathcal{F}, 1)
\in \Ext^2_{\mathcal{O}_X}(
\mathcal{F}, \mathcal{I} \otimes_{\mathcal{O}_X} \mathcal{F})$
du Lemme \ref{lemma-inf-obs-ext} est nulle.
\end{enumerate}
\end{lemma}

\begin{proof}
Cela résulte immédiatement de la caractérisation des
$\mathcal{O}_{X'}$-modules plats sur $S'$ donnée par la
Lemme \ref{lemma-deform-module} et de la
Lemme \ref{lemma-inf-obs-ext}.
\end{proof}






\section{Application aux modules plats sur les épaississements plats d'espaces annelés}
\label{section-flat}

\noindent
Considérons un diagramme commutatif
$$
\xymatrix{
(X, \mathcal{O}_X) \ar[r]_i \ar[d]_f & (X', \mathcal{O}_{X'}) \ar[d]^{f'} \\
(S, \mathcal{O}_S) \ar[r]^t & (S', \mathcal{O}_{S'})
}
$$
d'espaces annelés dont les flèches horizontales sont des épaississements du premier ordre,
comme dans la Situation \ref{situation-morphism-thickenings}. Posons
$\mathcal{I} = \Ker(i^\sharp) \subset \mathcal{O}_{X'}$ et
$\mathcal{J} = \Ker(t^\sharp) \subset \mathcal{O}_{S'}$.
Soit $\mathcal{F}$ un $\mathcal{O}_X$-module. Supposons que
\begin{enumerate}
\item $(f, f')$ soit un morphisme strict d'épaississements ;
\item $f'$ soit plat ;
\item $\mathcal{F}$ soit plat sur $S$.
\end{enumerate}
Remarquons que (1) $+$ (2) impliquent que $\mathcal{I} = f^*\mathcal{J}$
(appliquer le Lemme \ref{lemma-deform-module} à $\mathcal{O}_{X'}$).
La théorie de la section précédente prend une forme particulièrement simple
sous ces hypothèses. Nous résumons dans la lemme suivante les résultats déjà obtenus.

\begin{lemma}
\label{lemma-flat}
Dans la situation ci-dessus, on a les assertions suivantes.
\begin{enumerate}
\item Il existe un $\mathcal{O}_{X'}$-module $\mathcal{F}'$ plat sur
$S'$ tel que $i^*\mathcal{F}' \cong \mathcal{F}$, si et seulement si
la classe
$o(\mathcal{F}, f^*\mathcal{J} \otimes_{\mathcal{O}_X} \mathcal{F}, 1)
\in \Ext^2_{\mathcal{O}_X}(
\mathcal{F}, f^*\mathcal{J} \otimes_{\mathcal{O}_X} \mathcal{F})$
du Lemme \ref{lemma-inf-obs-ext} est nulle.
\item S'il existe un tel module, alors l'ensemble des classes d'isomorphisme
des relèvements est un espace principal homogène sous
$\Ext^1_{\mathcal{O}_X}(
\mathcal{F}, f^*\mathcal{J} \otimes_{\mathcal{O}_X} \mathcal{F})$.
\item Étant donné un relèvement $\mathcal{F}'$, l'ensemble des automorphismes de
$\mathcal{F}'$ dont l'image réciproque est $\text{id}_\mathcal{F}$ est canoniquement
isomorphe à $\Ext^0_{\mathcal{O}_X}(
\mathcal{F}, f^*\mathcal{J} \otimes_{\mathcal{O}_X} \mathcal{F})$.
\end{enumerate}
\end{lemma}

\begin{proof}
L'assertion (1) résulte du Lemme \ref{lemma-inf-obs-ext-rel},
car nous avons vu ci-dessus que $\mathcal{I} = f^*\mathcal{J}$.
L'assertion (2) résulte du Lemme \ref{lemma-inf-ext-rel}.
L'assertion (3) résulte du Lemme \ref{lemma-inf-map-rel}.
\end{proof}

\begin{situation}
\label{situation-ses-flat-thickenings}
Soit $f : (X, \mathcal{O}_X) \to (S, \mathcal{O}_S)$ un morphisme
d'espaces annelés. Considérons un diagramme commutatif
$$
\xymatrix{
(X'_1, \mathcal{O}'_1) \ar[r]_h \ar[d]_{f'_1} &
(X'_2, \mathcal{O}'_2) \ar[r] \ar[d]_{f'_2} &
(X'_3, \mathcal{O}'_3) \ar[d]_{f'_3} \\
(S'_1, \mathcal{O}_{S'_1}) \ar[r] &
(S'_2, \mathcal{O}_{S'_2}) \ar[r] &
(S'_3, \mathcal{O}_{S'_3})
}
$$
où (a) la ligne supérieure est une suite exacte courte d'épaississements du premier ordre
de $X$, (b) la ligne inférieure est une suite exacte courte d'épaississements du premier ordre
de $S$, (c) chaque $f'_i$ induit $f$ par restriction, (d) chaque couple
$(f, f_i')$ est un morphisme strict d'épaississements et (e) chaque $f'_i$ est plat.
Enfin, soit $\mathcal{F}'_2$ un $\mathcal{O}'_2$-module plat sur
$S'_2$, et posons $\mathcal{F} = \mathcal{F}'_2|_X$. Soit $\pi : X'_1 \to X$
le scindage canonique
(Remarque \ref{remark-short-exact-sequence-thickenings}).
\end{situation}

\begin{lemma}
\label{lemma-verify-iv}
Dans la Situation \ref{situation-ses-flat-thickenings}, les modules
$\pi^*\mathcal{F}$ et $h^*\mathcal{F}'_2$ sont des $\mathcal{O}'_1$-modules
plats sur $S'_1$ et dont la restriction est $\mathcal{F}$ sur $X$.
Leur différence (Lemme \ref{lemma-flat}) est un élément
$\theta$ de $\Ext^1_{\mathcal{O}_X}(
\mathcal{F}, f^*\mathcal{J}_1 \otimes_{\mathcal{O}_X} \mathcal{F})$
dont le cobord dans
$\Ext^2_{\mathcal{O}_X}(
\mathcal{F}, f^*\mathcal{J}_3 \otimes_{\mathcal{O}_X} \mathcal{F})$
est égal à l'obstruction (Lemme \ref{lemma-flat})
au relèvement de $\mathcal{F}$ en un $\mathcal{O}'_3$-module plat sur $S'_3$.
\end{lemma}

\begin{proof}
Remarquons que $\pi^*\mathcal{F}$ et $h^*\mathcal{F}'_2$
ont tous deux pour restriction $\mathcal{F}$ à $X$ et que les noyaux de
$\pi^*\mathcal{F} \to \mathcal{F}$ et $h^*\mathcal{F}'_2 \to \mathcal{F}$
sont donnés par $f^*\mathcal{J}_1 \otimes_{\mathcal{O}_X} \mathcal{F}$.
Ils sont donc plats d'après le Lemme \ref{lemma-deform-module}.
Prendre le cobord a un sens, car la suite de modules
$$
0 \to f^*\mathcal{J}_3 \otimes_{\mathcal{O}_X} \mathcal{F} \to
f^*\mathcal{J}_2 \otimes_{\mathcal{O}_X} \mathcal{F} \to
f^*\mathcal{J}_1 \otimes_{\mathcal{O}_X} \mathcal{F} \to 0
$$
est exacte courte par les hypothèses de la
Situation \ref{situation-ses-flat-thickenings}
et le fait que $\mathcal{F}$ est plat sur $S$.
L'assertion sur la classe d'obstruction est une traduction directe du résultat de la
Remarque \ref{remark-complex-thickenings-and-ses-modules}
dans cette situation particulière.
\end{proof}





\section{Déformations d'espaces annelés et complexe cotangent naïf}
\label{section-deformations-ringed-spaces}

\noindent
Dans cette section, nous utilisons le complexe cotangent naïf pour faire un peu
de théorie des déformations. Nous partons d'un épaississement du premier ordre
$t : (S, \mathcal{O}_S) \to (S', \mathcal{O}_{S'})$ d'espaces annelés.
Nous notons $\mathcal{J} = \Ker(t^\sharp)$ et identifions les espaces
topologiques sous-jacents à $S$ et $S'$.
Nous nous donnons en outre un morphisme d'espaces annelés
$f : (X, \mathcal{O}_X) \to (S, \mathcal{O}_S)$, un $\mathcal{O}_X$-module
$\mathcal{G}$ et une $f$-application $c : \mathcal{J} \to \mathcal{G}$
de faisceaux de modules (Faisceaux, Définition
\ref{sheaves-definition-f-map} et section
\ref{sheaves-section-ringed-spaces-functoriality-modules}).
Dans cette section, nous cherchons à savoir si l'on peut trouver
le point d'interrogation qui complète le diagramme suivant
\begin{equation}
\label{equation-to-solve-ringed-spaces}
\vcenter{
\xymatrix{
0 \ar[r] & \mathcal{G} \ar[r] & {?} \ar[r] & \mathcal{O}_X \ar[r] & 0 \\
0 \ar[r] & \mathcal{J} \ar[u]^c \ar[r] & \mathcal{O}_{S'} \ar[u] \ar[r] &
\mathcal{O}_S \ar[u] \ar[r] & 0
}
}
\end{equation}
(où les flèches verticales sont des $f$-applications), et aussi dans quelle
mesure la solution est unique (lorsqu'elle existe). Plus précisément,
nous cherchons un épaississement du premier ordre
$i : (X, \mathcal{O}_X) \to (X', \mathcal{O}_{X'})$
et un morphisme d'épaississements $(f, f')$ comme dans
(\ref{equation-morphism-thickenings}), où $\Ker(i^\sharp)$ est identifié
à $\mathcal{G}$, tels que $(f')^\sharp$ induise l'application donnée $c$.
Nous dirons que $X'$ est une {\it solution} de
(\ref{equation-to-solve-ringed-spaces}).

\begin{lemma}
\label{lemma-huge-diagram-ringed-spaces}
Supposons donné un diagramme commutatif de morphismes d'espaces annelés
\begin{equation}
\label{equation-huge-1}
\vcenter{
\xymatrix{
& (X_2, \mathcal{O}_{X_2}) \ar[r]_{i_2} \ar[d]_{f_2} \ar[ddl]_g &
(X'_2, \mathcal{O}_{X'_2}) \ar[d]^{f'_2} \\
& (S_2, \mathcal{O}_{S_2}) \ar[r]^{t_2} \ar[ddl]|\hole &
(S'_2, \mathcal{O}_{S'_2}) \ar[ddl] \\
(X_1, \mathcal{O}_{X_1}) \ar[r]_{i_1} \ar[d]_{f_1} &
(X'_1, \mathcal{O}_{X'_1}) \ar[d]^{f'_1} \\
(S_1, \mathcal{O}_{S_1}) \ar[r]^{t_1} &
(S'_1, \mathcal{O}_{S'_1})
}
}
\end{equation}
dont les flèches horizontales sont des épaississements du premier ordre. Posons
$\mathcal{G}_j = \Ker(i_j^\sharp)$ et supposons donnée
une $g$-application $\nu : \mathcal{G}_1 \to \mathcal{G}_2$ de modules
donnant lieu au diagramme commutatif
\begin{equation}
\label{equation-huge-2}
\vcenter{
\xymatrix{
& 0 \ar[r] & \mathcal{G}_2 \ar[r] &
\mathcal{O}_{X'_2} \ar[r] &
\mathcal{O}_{X_2} \ar[r] & 0 \\
& 0 \ar[r]|\hole &
\mathcal{J}_2 \ar[u]_{c_2} \ar[r] &
\mathcal{O}_{S'_2} \ar[u] \ar[r]|\hole &
\mathcal{O}_{S_2} \ar[u] \ar[r] & 0 \\
0 \ar[r] & \mathcal{G}_1 \ar[ruu] \ar[r] &
\mathcal{O}_{X'_1} \ar[r] &
\mathcal{O}_{X_1} \ar[ruu] \ar[r] & 0 \\
0 \ar[r] & \mathcal{J}_1 \ar[ruu]|\hole \ar[u]^{c_1} \ar[r] &
\mathcal{O}_{S'_1} \ar[ruu]|\hole \ar[u] \ar[r] &
\mathcal{O}_{S_1} \ar[ruu]|\hole \ar[u] \ar[r] & 0
}
}
\end{equation}
dont les faces avant et arrière sont des solutions de
(\ref{equation-to-solve-ringed-spaces}).
\begin{enumerate}
\item Il existe un élément canonique de
$\Ext^1_{\mathcal{O}_{X_2}}(Lg^*\NL_{X_1/S_1}, \mathcal{G}_2)$
dont l'annulation est une condition nécessaire et suffisante pour qu'il existe
un morphisme d'espaces annelés $X'_2 \to X'_1$ complétant
(\ref{equation-huge-1}) de façon compatible à $\nu$.
\item S'il existe un morphisme $X'_2 \to X'_1$ complétant
(\ref{equation-huge-1}) de façon compatible à $\nu$, l'ensemble de tous ces
morphismes est un espace principal homogène sous
$$
\Hom_{\mathcal{O}_{X_1}}(\Omega_{X_1/S_1}, g_*\mathcal{G}_2) =
\Hom_{\mathcal{O}_{X_2}}(g^*\Omega_{X_1/S_1}, \mathcal{G}_2) =
\Ext^0_{\mathcal{O}_{X_2}}(Lg^*\NL_{X_1/S_1}, \mathcal{G}_2).
$$
\end{enumerate}
\end{lemma}

\begin{proof}
Le complexe cotangent naïf $\NL_{X_1/S_1}$ est défini dans Modules,
Définition \ref{modules-definition-cotangent-complex-morphism-ringed-topoi}.
Les égalités de la dernière assertion de la lemme résultent de ce que
$g^*$ est adjoint à $g_*$, de l'égalité
$H^0(\NL_{X_1/S_1}) = \Omega_{X_1/S_1}$ (par construction du complexe
cotangent naïf) et du fait que $Lg^*$ est le foncteur dérivé à gauche de
$g^*$. Nous travaillerons donc, dans la suite de la démonstration, avec les
groupes
$\Ext^k_{\mathcal{O}_{X_2}}(Lg^*\NL_{X_1/S_1}, \mathcal{G}_2)$,
$k = 0, 1$. Montrons d'abord que l'on peut se ramener au cas où les espaces
topologiques sous-jacents à tous les espaces annelés de la lemme sont les mêmes.

\medskip\noindent
Pour cela, remarquons que $g^{-1}\NL_{X_1/S_1}$ est égal au complexe
cotangent naïf de l'homomorphisme de faisceaux d'anneaux
$g^{-1}f_1^{-1}\mathcal{O}_{S_1} \to g^{-1}\mathcal{O}_{X_1}$ ; voir
Modules, Lemme \ref{modules-lemma-pullback-NL}.
De plus, le terme de degré $0$ de $\NL_{X_1/S_1}$ est un
$\mathcal{O}_{X_1}$-module plat ; l'application canonique
$$
Lg^*\NL_{X_1/S_1}
\longrightarrow
g^{-1}\NL_{X_1/S_1} \otimes_{g^{-1}\mathcal{O}_{X_1}} \mathcal{O}_{X_2}
$$
induit donc un isomorphisme sur les faisceaux de cohomologie en degrés $0$ et
$-1$. On peut ainsi remplacer les groupes Ext de la lemme par
$$
\Ext^k_{g^{-1}\mathcal{O}_{X_1}}(g^{-1}\NL_{X_1/S_1}, \mathcal{G}_2) =
\Ext^k_{g^{-1}\mathcal{O}_{X_1}}(
\NL_{g^{-1}\mathcal{O}_{X_1}/g^{-1}f_1^{-1}\mathcal{O}_{S_1}}, \mathcal{G}_2)
$$
L'ensemble des morphismes d'espaces annelés $X'_2 \to X'_1$ complétant
(\ref{equation-huge-1}) de façon compatible à $\nu$ est en bijection avec
l'ensemble des homomorphismes de $g^{-1}f_1^{-1}\mathcal{O}_{S'_1}$-algèbres
$g^{-1}\mathcal{O}_{X'_1} \to \mathcal{O}_{X'_2}$ compatibles avec
$f^\sharp$ et $\nu$. Nous voyons ainsi que nous pouvons supposer donné un
diagramme (\ref{equation-huge-2}) de faisceaux sur $X$ et chercher un
homomorphisme de faisceaux d'anneaux
$\mathcal{O}_{X'_1} \to \mathcal{O}_{X'_2}$ qui le complète.

\medskip\noindent
Dans la suite de la démonstration, nous supposons que tous les espaces
topologiques sous-jacents sont identiques, c'est-à-dire que nous disposons
d'un diagramme (\ref{equation-huge-2}) de faisceaux sur un espace $X$ et
cherchons les homomorphismes de faisceaux d'anneaux
$\mathcal{O}_{X'_1} \to \mathcal{O}_{X'_2}$ qui le complètent.
Comme groupes Ext, nous utiliserons
$\Ext^k_{\mathcal{O}_{X_1}}(
\NL_{\mathcal{O}_{X_1}/\mathcal{O}_{S_1}}, \mathcal{G}_2)$, $k = 0, 1$.

\medskip\noindent
Étape 1. Construction de la classe d'obstruction. Considérons le faisceau
d'ensembles
$$
\mathcal{E} = \mathcal{O}_{X'_1} \times_{\mathcal{O}_{X_2}} \mathcal{O}_{X'_2}
$$
Il est muni d'une application surjective
$\alpha : \mathcal{E} \to \mathcal{O}_{X_1}$ ; nous pouvons donc employer
$\NL(\alpha)$ à la place de
$\NL_{\mathcal{O}_{X_1}/\mathcal{O}_{S_1}}$ ; voir Modules, Lemme
\ref{modules-lemma-NL-up-to-qis}.
Posons
$$
\mathcal{I}' =
\Ker(\mathcal{O}_{S'_1}[\mathcal{E}] \to \mathcal{O}_{X_1})
\quad\text{et}\quad
\mathcal{I} =
\Ker(\mathcal{O}_{S_1}[\mathcal{E}] \to \mathcal{O}_{X_1})
$$
Il existe une surjection $\mathcal{I}' \to \mathcal{I}$ dont le noyau est
$\mathcal{J}_1\mathcal{O}_{S'_1}[\mathcal{E}]$.
Nous obtenons deux homomorphismes de $\mathcal{O}_{S'_1}$-algèbres
$$
a : \mathcal{O}_{S'_1}[\mathcal{E}] \to \mathcal{O}_{X'_1}
\quad\text{et}\quad
b : \mathcal{O}_{S'_1}[\mathcal{E}] \to \mathcal{O}_{X'_2}
$$
qui induisent des applications
$a|_{\mathcal{I}'} : \mathcal{I}' \to \mathcal{G}_1$ et
$b|_{\mathcal{I}'} : \mathcal{I}' \to \mathcal{G}_2$.
Les deux applications $a$ et $b$ annulent $(\mathcal{I}')^2$.
De plus, $a$ et $b$ coïncident sur
$\mathcal{J}_1\mathcal{O}_{S'_1}[\mathcal{E}]$ en tant qu'applications à
valeurs dans $\mathcal{G}_2$, car le carré de gauche de
(\ref{equation-huge-2}) est commutatif. Par conséquent, la différence
$b|_{\mathcal{I}'} - \nu \circ a|_{\mathcal{I}'}$
induit une application $\mathcal{O}_{X_1}$-linéaire bien définie
$$
\xi : \mathcal{I}/\mathcal{I}^2 \longrightarrow \mathcal{G}_2
$$
qui envoie la classe d'une section locale $f$ de $\mathcal{I}$ sur
$\nu(a(f')) - b(f')$, où $f'$ est un relèvement de $f$ en une section locale
de $\mathcal{I}'$. Nous notons
$[\xi] \in \Ext^1_{\mathcal{O}_{X_1}}(\NL(\alpha), \mathcal{G}_2)$
son image (voir ci-dessous).

\medskip\noindent
Étape 2. L'annulation de $[\xi]$ est nécessaire. Écrivons
$\Omega = \Omega_{\mathcal{O}_{S_1}[\mathcal{E}]/\mathcal{O}_{S_1}}
\otimes_{\mathcal{O}_{S_1}[\mathcal{E}]} \mathcal{O}_{X_1}$.
Remarquons que $\NL(\alpha) = (\mathcal{I}/\mathcal{I}^2 \to \Omega)$
s'insère dans un triangle distingué
$$
\Omega[0] \to
\NL(\alpha) \to
\mathcal{I}/\mathcal{I}^2[1] \to
\Omega[1]
$$
On voit donc que $[\xi]$ est nul si et seulement si $\xi$ est un composé
$\mathcal{I}/\mathcal{I}^2 \to \Omega \to \mathcal{G}_2$
pour une certaine application $\Omega \to \mathcal{G}_2$. Supposons qu'il
existe un homomorphisme de faisceaux d'anneaux
$\varphi : \mathcal{O}_{X'_1} \to \mathcal{O}_{X'_2}$ complétant
(\ref{equation-huge-2}). Considérons alors l'application
$\mathcal{O}_{S'_1}[\mathcal{E}] \to \mathcal{G}_2$,
$f' \mapsto b(f') - \varphi(a(f'))$. Un calcul montre qu'elle annule
$\mathcal{J}_1\mathcal{O}_{S'_1}[\mathcal{E}]$ et induit une dérivation
$\mathcal{O}_{S_1}[\mathcal{E}] \to \mathcal{G}_2$.
L'application linéaire $\Omega \to \mathcal{G}_2$ ainsi obtenue atteste que
$[\xi] = 0$ dans ce cas.

\medskip\noindent
Étape 3. L'annulation de $[\xi]$ est suffisante. Soit
$\theta : \Omega \to \mathcal{G}_2$ une application
$\mathcal{O}_{X_1}$-linéaire telle que $\xi$ soit égal à
$\theta \circ (\mathcal{I}/\mathcal{I}^2 \to \Omega)$.
Un calcul montre alors que
$$
b + \theta \circ d : \mathcal{O}_{S'_1}[\mathcal{E}] \to \mathcal{O}_{X'_2}
$$
restreint à $\mathcal{I}'$ coïncide avec
$\nu \circ a : \mathcal{I}' \to \mathcal{G}_2$. Comme
$\mathcal{O}_{X'_1}$ est la somme amalgamée de
$\mathcal{I}' \to \mathcal{O}_{S'_1}[\mathcal{E}]$ et
$\mathcal{I}' \to \mathcal{G}_1$, les applications
$b + \theta \circ d$ et $a$ définissent une application
$\mathcal{O}_{X'_1} \to \mathcal{O}_{X'_2}$ complétant
(\ref{equation-huge-2}).

\medskip\noindent
Démonstration de (2) dans le cas particulier ci-dessus. Omise. Indication :
c'est exactement la même que la démonstration de (2) de la
Lemme \ref{lemma-huge-diagram}.
\end{proof}

\begin{lemma}
\label{lemma-NL-represent-ext-class}
Soit $X$ un espace topologique. Soit $\mathcal{A} \to \mathcal{B}$ un
homomorphisme de faisceaux d'anneaux. Soit $\mathcal{G}$ un
$\mathcal{B}$-module. Soit
$\xi \in \Ext^1_\mathcal{B}(\NL_{\mathcal{B}/\mathcal{A}}, \mathcal{G})$.
Il existe une application de faisceaux d'ensembles
$\alpha : \mathcal{E} \to \mathcal{B}$ telle que
$\xi \in \Ext^1_\mathcal{B}(\NL(\alpha), \mathcal{G})$
soit la classe d'une application
$\mathcal{I}/\mathcal{I}^2 \to \mathcal{G}$
(voir la démonstration pour les notations).
\end{lemma}

\begin{proof}
Rappelons que, si $\alpha : \mathcal{E} \to \mathcal{B}$ est telle que
$\mathcal{A}[\mathcal{E}] \to \mathcal{B}$ soit surjective, de noyau
$\mathcal{I}$, le complexe
$\NL(\alpha) = (\mathcal{I}/\mathcal{I}^2 \to 
\Omega_{\mathcal{A}[\mathcal{E}]/\mathcal{A}}
\otimes_{\mathcal{A}[\mathcal{E}]} \mathcal{B})$ est canoniquement
isomorphe à $\NL_{\mathcal{B}/\mathcal{A}}$ ; voir Modules, Lemme
\ref{modules-lemma-NL-up-to-qis}.
Remarquons de plus que
$\Omega = \Omega_{\mathcal{A}[\mathcal{E}]/\mathcal{A}}
\otimes_{\mathcal{A}[\mathcal{E}]} \mathcal{B}$ est le faisceau associé
au préfaisceau
$U \mapsto \bigoplus_{e \in \mathcal{E}(U)} \mathcal{B}(U)$.
Autrement dit, $\Omega$ est le $\mathcal{B}$-module libre sur le faisceau
d'ensembles $\mathcal{E}$ et il existe en particulier une application
canonique $\mathcal{E} \to \Omega$.

\medskip\noindent
Cela étant, choisissons un $\mathcal{E}$ (par exemple
$\mathcal{E} = \mathcal{B}$ comme dans la définition du complexe cotangent
naïf). L'obstruction à l'écriture de $\xi$ comme classe d'une application
$\mathcal{I}/\mathcal{I}^2 \to \mathcal{G}$ est un élément de
$\Ext^1_\mathcal{B}(\Omega, \mathcal{G})$. Supposons qu'il soit représenté
par l'extension
$0 \to \mathcal{G} \to \mathcal{H} \to \Omega \to 0$
de $\mathcal{B}$-modules. Considérons le faisceau d'ensembles
$\mathcal{E}' = \mathcal{E} \times_\Omega \mathcal{H}$,
muni de l'application induite $\alpha' : \mathcal{E}' \to \mathcal{B}$.
Posons $\mathcal{I}' = \Ker(\mathcal{A}[\mathcal{E}'] \to \mathcal{B})$
et $\Omega' = \Omega_{\mathcal{A}[\mathcal{E}']/\mathcal{A}}
\otimes_{\mathcal{A}[\mathcal{E}']} \mathcal{B}$.
L'image réciproque de $\xi$ par le quasi-isomorphisme
$\NL(\alpha') \to \NL(\alpha)$ s'envoie sur zéro dans
$\Ext^1_\mathcal{B}(\Omega', \mathcal{G})$, car l'image réciproque de
l'extension $\mathcal{H}$ par l'application $\Omega' \to \Omega$ est
scindée : en effet, $\Omega'$ est le $\mathcal{B}$-module libre sur le
faisceau d'ensembles $\mathcal{E}'$ et, par construction, il existe un
diagramme commutatif
$$
\xymatrix{
\mathcal{E}' \ar[r] \ar[d] & \mathcal{E} \ar[d] \\
\mathcal{H} \ar[r] & \Omega
}
$$
Cela termine la démonstration.
\end{proof}

\begin{lemma}
\label{lemma-choices-ringed-spaces}
S'il existe une solution de (\ref{equation-to-solve-ringed-spaces}),
l'ensemble des classes d'isomorphisme de solutions est un espace principal
homogène sous $\Ext^1_{\mathcal{O}_X}(\NL_{X/S}, \mathcal{G})$.
\end{lemma}

\begin{proof}
Remarquons tout de suite qu'étant données deux solutions $X'_1$ et $X'_2$
de (\ref{equation-to-solve-ringed-spaces}), la
Lemme \ref{lemma-huge-diagram-ringed-spaces} fournit un élément d'obstruction
$o(X'_1, X'_2) \in \Ext^1_{\mathcal{O}_X}(\NL_{X/S}, \mathcal{G})$
à l'existence d'une application $X'_1 \to X'_2$. Manifestement, cet élément
est l'obstruction à l'existence d'un isomorphisme, et sépare donc les classes
d'isomorphisme. Pour achever la démonstration, il suffit par conséquent de
montrer qu'étant donnés une solution $X'$ et un élément
$\xi \in \Ext^1_{\mathcal{O}_X}(\NL_{X/S}, \mathcal{G})$,
on peut trouver une seconde solution $X'_\xi$ telle que
$o(X', X'_\xi) = \xi$.

\medskip\noindent
Choisissons $\alpha : \mathcal{E} \to \mathcal{O}_X$ comme dans la
Lemme \ref{lemma-NL-represent-ext-class} pour la classe $\xi$.
Considérons la surjection
$f^{-1}\mathcal{O}_S[\mathcal{E}] \to \mathcal{O}_X$
de noyau $\mathcal{I}$ et le complexe cotangent naïf correspondant
$\NL(\alpha) = (\mathcal{I}/\mathcal{I}^2 \to
\Omega_{f^{-1}\mathcal{O}_S[\mathcal{E}]/f^{-1}\mathcal{O}_S}
\otimes_{f^{-1}\mathcal{O}_S[\mathcal{E}]} \mathcal{O}_X)$.
D'après la lemme, $\xi$ est la classe d'un morphisme
$\delta : \mathcal{I}/\mathcal{I}^2 \to \mathcal{G}$.
Après avoir remplacé $\mathcal{E}$ par
$\mathcal{E} \times_{\mathcal{O}_X} \mathcal{O}_{X'}$, nous pouvons
également supposer que $\alpha$ se factorise par une application
$\alpha' : \mathcal{E} \to \mathcal{O}_{X'}$.

\medskip\noindent
Ces choix déterminent un homomorphisme de
$f^{-1}\mathcal{O}_{S'}$-algèbres
$\varphi : f^{-1}\mathcal{O}_{S'}[\mathcal{E}] \to \mathcal{O}_{X'}$.
Posons
$\mathcal{I}' = \Ker(f^{-1}\mathcal{O}_{S'}[\mathcal{E}] \to \mathcal{O}_X)$.
Remarquons que $\varphi$ induit une application
$\varphi|_{\mathcal{I}'} : \mathcal{I}' \to \mathcal{G}$
et que $\mathcal{O}_{X'}$ est la somme amalgamée, comme dans le diagramme
suivant
$$
\xymatrix{
0 \ar[r] & \mathcal{G} \ar[r] & \mathcal{O}_{X'} \ar[r] &
\mathcal{O}_X \ar[r] & 0 \\
0 \ar[r] & \mathcal{I}' \ar[u]^{\varphi|_{\mathcal{I}'}} \ar[r] &
f^{-1}\mathcal{O}_{S'}[\mathcal{E}] \ar[u] \ar[r] &
\mathcal{O}_X \ar[u]_{=} \ar[r] & 0
}
$$
Soit $\psi : \mathcal{I}' \to \mathcal{G}$ la somme de l'application
$\varphi|_{\mathcal{I}'}$ et du composé
$$
\mathcal{I}' \to \mathcal{I}'/(\mathcal{I}')^2 \to
\mathcal{I}/\mathcal{I}^2 \xrightarrow{\delta} \mathcal{G}.
$$
La somme amalgamée suivant $\psi$ est alors une autre extension d'anneaux
$\mathcal{O}_{X'_\xi}$ qui s'insère dans un diagramme comme ci-dessus.
Un calcul (omis) montre que $o(X', X'_\xi) = \xi$, comme voulu.
\end{proof}

\begin{lemma}
\label{lemma-extensions-of-relative-ringed-spaces}
Soit $f : (X, \mathcal{O}_X) \to (S, \mathcal{O}_S)$ un morphisme
d'espaces annelés. Soit $\mathcal{G}$ un $\mathcal{O}_X$-module.
L'ensemble des classes d'isomorphisme des extensions de
$f^{-1}\mathcal{O}_S$-algèbres
$$
0 \to \mathcal{G} \to \mathcal{O}_{X'} \to \mathcal{O}_X \to 0
$$
où $\mathcal{G}$ est un idéal de carré nul\footnote{Autrement dit,
l'ensemble des classes d'isomorphisme des épaississements du premier ordre
$i : X \to X'$ au-dessus de $S$, munis d'un isomorphisme
$\mathcal{G} \to \Ker(i^\sharp)$ de $\mathcal{O}_X$-modules.}
est canoniquement en bijection avec
$\Ext^1_{\mathcal{O}_X}(\NL_{X/S}, \mathcal{G})$.
\end{lemma}

\begin{proof}
Pour le démontrer, appliquons les résultats précédents au cas où
(\ref{equation-to-solve-ringed-spaces}) est donné par le diagramme
$$
\xymatrix{
0 \ar[r] &
\mathcal{G} \ar[r] &
{?} \ar[r] &
\mathcal{O}_X \ar[r] & 0 \\
0 \ar[r] &
0 \ar[u] \ar[r] &
\mathcal{O}_S \ar[u] \ar[r]^{\text{id}} &
\mathcal{O}_S \ar[u] \ar[r] & 0
}
$$
Le lemme résulte donc du Lemme \ref{lemma-choices-ringed-spaces}
et de l'existence d'une solution, à savoir
$\mathcal{G} \oplus \mathcal{O}_X$.
(Voir la remarque ci-dessous pour une construction directe de la bijection.)
\end{proof}

\begin{remark}
\label{remark-extensions-of-relative-ringed-spaces}
Soient $f : (X, \mathcal{O}_X) \to (S, \mathcal{O}_S)$ et
$\mathcal{G}$ comme dans la
Lemme \ref{lemma-extensions-of-relative-ringed-spaces}.
Considérons une extension
$0 \to \mathcal{G} \to \mathcal{O}_{X'} \to \mathcal{O}_X \to 0$
comme dans la lemme. On peut choisir un faisceau d'ensembles $\mathcal{E}$
et un diagramme commutatif
$$
\xymatrix{
\mathcal{E} \ar[d]_{\alpha'} \ar[rd]^\alpha \\
\mathcal{O}_{X'} \ar[r] & \mathcal{O}_X
}
$$
tels que $f^{-1}\mathcal{O}_S[\mathcal{E}] \to \mathcal{O}_X$
soit surjective, de noyau $\mathcal{J}$.
(On peut par exemple prendre n'importe quel faisceau d'ensembles se
surjectant sur $\mathcal{O}_{X'}$.) Alors
$$
\NL_{X/S} \cong \NL(\alpha) = 
\left(
\mathcal{J}/\mathcal{J}^2
\longrightarrow
\Omega_{f^{-1}\mathcal{O}_S[\mathcal{E}]/f^{-1}\mathcal{O}_S}
\otimes_{f^{-1}\mathcal{O}_S[\mathcal{E}]} \mathcal{O}_X\right)
$$
Voir Modules, section \ref{modules-section-netherlander}, et en particulier
Lemme \ref{modules-lemma-NL-up-to-qis}. Bien entendu, $\alpha'$ détermine une
application $f^{-1}\mathcal{O}_S[\mathcal{E}] \to \mathcal{O}_{X'}$,
qui détermine à son tour une application
$$
\mathcal{J}/\mathcal{J}^2 \longrightarrow \mathcal{G}
$$
laquelle détermine à son tour l'élément de
$\Ext^1_{\mathcal{O}_X}(\NL(\alpha), \mathcal{G}) =
\Ext^1_{\mathcal{O}_X}(\NL_{X/S}, \mathcal{G})$
correspondant à $\mathcal{O}_{X'}$ par la bijection de la lemme.
\end{remark}

\begin{lemma}
\label{lemma-extensions-of-relative-ringed-spaces-functorial}
Soient $f : (X, \mathcal{O}_X) \to (S, \mathcal{O}_S)$ et
$g : (Y, \mathcal{O}_Y) \to (X, \mathcal{O}_X)$ des morphismes
d'espaces annelés. Soit $\mathcal{F}$ un $\mathcal{O}_X$-module.
Soit $\mathcal{G}$ un $\mathcal{O}_Y$-module. Soit
$c : \mathcal{F} \to \mathcal{G}$ une $g$-application. Enfin, considérons
\begin{enumerate}
\item[(a)] $0 \to \mathcal{F} \to \mathcal{O}_{X'} \to \mathcal{O}_X \to 0$,
une extension de $f^{-1}\mathcal{O}_S$-algèbres,
correspondant à
$\xi \in \Ext^1_{\mathcal{O}_X}(\NL_{X/S}, \mathcal{F})$, et
\item[(b)] $0 \to \mathcal{G} \to \mathcal{O}_{Y'} \to \mathcal{O}_Y \to 0$,
une extension de $g^{-1}f^{-1}\mathcal{O}_S$-algèbres,
correspondant à
$\zeta \in \Ext^1_{\mathcal{O}_Y}(\NL_{Y/S}, \mathcal{G})$.
\end{enumerate}
Voir le Lemme \ref{lemma-extensions-of-relative-ringed-spaces}.
Il existe alors un $S$-morphisme $g' : Y' \to X'$ compatible avec $g$ et
$c$ si et seulement si $\xi$ et $\zeta$ ont la même image dans
$\Ext^1_{\mathcal{O}_Y}(Lg^*\NL_{X/S}, \mathcal{G})$.
\end{lemma}

\begin{proof}
L'énoncé a bien un sens, car nous disposons des applications
$$
\Ext^1_{\mathcal{O}_X}(\NL_{X/S}, \mathcal{F}) \to
\Ext^1_{\mathcal{O}_Y}(Lg^*\NL_{X/S}, Lg^*\mathcal{F}) \to
\Ext^1_{\mathcal{O}_Y}(Lg^*\NL_{X/S}, \mathcal{G})
$$
obtenues à l'aide de l'application
$Lg^*\mathcal{F} \to g^*\mathcal{F} \xrightarrow{c} \mathcal{G}$,
ainsi que de l'application
$$
\Ext^1_{\mathcal{O}_Y}(\NL_{Y/S}, \mathcal{G}) \to
\Ext^1_{\mathcal{O}_Y}(Lg^*\NL_{X/S}, \mathcal{G})
$$
obtenue à l'aide de $Lg^*\NL_{X/S} \to \NL_{Y/S}$.
L'énoncé de la lemme se déduit de la
Lemme \ref{lemma-huge-diagram-ringed-spaces} appliquée au diagramme
$$
\xymatrix{
& 0 \ar[r] &
\mathcal{G} \ar[r] &
\mathcal{O}_{Y'} \ar[r] &
\mathcal{O}_Y \ar[r] & 0 \\
& 0 \ar[r]|\hole & 0 \ar[u] \ar[r] &
\mathcal{O}_S \ar[u] \ar[r]|\hole &
\mathcal{O}_S \ar[u] \ar[r] & 0 \\
0 \ar[r] &
\mathcal{F} \ar[ruu] \ar[r] &
\mathcal{O}_{X'} \ar[r] &
\mathcal{O}_X \ar[ruu] \ar[r] & 0 \\
0 \ar[r] & 0 \ar[ruu]|\hole \ar[u] \ar[r] &
\mathcal{O}_S \ar[ruu]|\hole \ar[u] \ar[r] &
\mathcal{O}_S \ar[ruu]|\hole \ar[u] \ar[r] & 0
}
$$
et d'une compatibilité entre les constructions des démonstrations des
Lemmes \ref{lemma-extensions-of-relative-ringed-spaces} et
\ref{lemma-huge-diagram-ringed-spaces}, dont nous omettons l'énoncé et la
démonstration. (Voir la remarque ci-dessous pour un argument direct.)
\end{proof}

\begin{remark}
\label{remark-extensions-of-relative-ringed-spaces-functorial}
Soient $f : (X, \mathcal{O}_X) \to (S, \mathcal{O}_S)$,
$g : (Y, \mathcal{O}_Y) \to (X, \mathcal{O}_X)$,
$\mathcal{F}$,
$\mathcal{G}$,
$c : \mathcal{F} \to \mathcal{G}$,
$0 \to \mathcal{F} \to \mathcal{O}_{X'} \to \mathcal{O}_X \to 0$,
$\xi \in \Ext^1_{\mathcal{O}_X}(\NL_{X/S}, \mathcal{F})$,
$0 \to \mathcal{G} \to \mathcal{O}_{Y'} \to \mathcal{O}_Y \to 0$
et $\zeta \in \Ext^1_{\mathcal{O}_Y}(\NL_{Y/S}, \mathcal{G})$
comme dans la
Lemme \ref{lemma-extensions-of-relative-ringed-spaces-functorial}.
En prenant la somme amalgamée suivant
$c : g^{-1}\mathcal{F} \to \mathcal{G}$, on construit une extension
$$
\xymatrix{
0 \ar[r] &
\mathcal{G} \ar[r] &
\mathcal{O}'_1 \ar[r] &
g^{-1}\mathcal{O}_X \ar[r] & 0 \\
0 \ar[r] &
g^{-1}\mathcal{F} \ar[u]^c \ar[r] &
g^{-1}\mathcal{O}_{X'} \ar[u] \ar[r] &
g^{-1}\mathcal{O}_X \ar@{=}[u] \ar[r] & 0
}
$$
En prenant l'image réciproque suivant
$g^\sharp : g^{-1}\mathcal{O}_X \to \mathcal{O}_Y$,
on construit une extension
$$
\xymatrix{
0 \ar[r] &
\mathcal{G} \ar[r] &
\mathcal{O}_{Y'} \ar[r] &
\mathcal{O}_Y \ar[r] & 0 \\
0 \ar[r] &
\mathcal{G} \ar@{=}[u] \ar[r] &
\mathcal{O}'_2 \ar[u] \ar[r] &
g^{-1}\mathcal{O}_X \ar[u] \ar[r] & 0
}
$$
Une chasse au diagramme montre qu'il existe un $S$-morphisme
$Y' \to X'$ compatible avec $g$ et $c$ si et seulement si
$\mathcal{O}'_1$ est isomorphe à $\mathcal{O}'_2$ comme extension de
$g^{-1}f^{-1}\mathcal{O}_S$-algèbres de $g^{-1}\mathcal{O}_X$ par
$\mathcal{G}$. D'après la
Lemme \ref{lemma-extensions-of-relative-ringed-spaces}, ces extensions sont
classifiées par le membre de gauche de
$$
\Ext^1_{g^{-1}\mathcal{O}_X}(
\NL_{g^{-1}\mathcal{O}_X/g^{-1}f^{-1}\mathcal{O}_S}, \mathcal{G}) =
\Ext^1_{\mathcal{O}_Y}(Lg^*\NL_{X/S}, \mathcal{G})
$$
Ici, l'égalité provient de l'adjonction tensor-hom et des égalités
$$
\NL_{g^{-1}\mathcal{O}_X/g^{-1}f^{-1}\mathcal{O}_S} = g^{-1}\NL_{X/S}
\quad\text{et}\quad
Lg^*\NL_{X/S} =
g^{-1}\NL_{X/S} \otimes_{g^{-1}\mathcal{O}_X}^\mathbf{L} \mathcal{O}_Y
$$
Pour la première, voir Modules, Lemme \ref{modules-lemma-pullback-NL} ;
la seconde résulte de la définition de l'image réciproque dérivée.
Ainsi, pour voir que la
Lemme \ref{lemma-extensions-of-relative-ringed-spaces-functorial} est vraie,
il suffit de montrer que $\mathcal{O}'_1$ correspond à l'image de $\xi$ et
que $\mathcal{O}'_2$ correspond à l'image de $\zeta$.
La correspondance entre $\xi$ et $\mathcal{O}'_1$ résulte immédiatement de
la construction de la classe $\xi$ dans la
Remarque \ref{remark-extensions-of-relative-ringed-spaces}.
Pour la correspondance entre $\zeta$ et $\mathcal{O}'_2$, choisissons d'abord
un diagramme commutatif
$$
\xymatrix{
\mathcal{E} \ar[d]_{\beta'} \ar[rd]^\beta \\
\mathcal{O}_{Y'} \ar[r] & \mathcal{O}_Y
}
$$
tel que $g^{-1}f^{-1}\mathcal{O}_S[\mathcal{E}] \to \mathcal{O}_Y$
soit surjective, de noyau $\mathcal{K}$. Choisissons ensuite un diagramme
commutatif
$$
\xymatrix{
\mathcal{E} \ar[d]_{\beta'} &
\mathcal{E}' \ar[l]^\varphi \ar[d]_{\alpha'} \ar[rd]^\alpha \\
\mathcal{O}_{Y'} &
\mathcal{O}'_2 \ar[l] \ar[r] &
g^{-1}\mathcal{O}_X
}
$$
tel que
$g^{-1}f^{-1}\mathcal{O}_S[\mathcal{E}'] \to g^{-1}\mathcal{O}_X$
soit surjective, de noyau $\mathcal{J}$. (Il suffit par exemple de prendre
$\mathcal{E}' = \mathcal{E} \amalg \mathcal{O}'_2$ comme faisceau
d'ensembles.) L'application $\varphi$ induit une application de complexes
$\NL(\alpha) \to \NL(\beta)$
(notations de Modules, section \ref{modules-section-netherlander})
et en particulier
$\bar\varphi : \mathcal{J}/\mathcal{J}^2 \to \mathcal{K}/\mathcal{K}^2$.
On a alors $\NL(\alpha) \cong \NL_{Y/S}$ et
$\NL(\beta) \cong \NL_{g^{-1}\mathcal{O}_X/g^{-1}f^{-1}\mathcal{O}_S}$,
et l'application de complexes $\NL(\alpha) \to \NL(\beta)$ représente
l'application $Lg^*\NL_{X/S} \to \NL_{Y/S}$ employée dans l'énoncé de la
Lemme \ref{lemma-extensions-of-relative-ringed-spaces-functorial}
(voir la première partie de sa démonstration). Or $\zeta$ correspond à la
classe de l'application $\mathcal{K}/\mathcal{K}^2 \to \mathcal{G}$
induite par $\beta'$ ; voir la
Remarque \ref{remark-extensions-of-relative-ringed-spaces}.
De même, l'extension $\mathcal{O}'_2$ correspond à l'application
$\mathcal{J}/\mathcal{J}^2 \to \mathcal{G}$ induite par $\alpha'$.
Le diagramme commutatif ci-dessus montre que cette application est le composé
de l'application $\mathcal{K}/\mathcal{K}^2 \to \mathcal{G}$ induite par
$\beta'$ et de l'application
$\bar\varphi : \mathcal{J}/\mathcal{J}^2 \to \mathcal{K}/\mathcal{K}^2$.
C'est la compatibilité recherchée.
\end{remark}

\begin{lemma}
\label{lemma-parametrize-solutions-ringed-spaces}
Soient $t : (S, \mathcal{O}_S) \to (S', \mathcal{O}_{S'})$,
$\mathcal{J} = \Ker(t^\sharp)$,
$f : (X, \mathcal{O}_X) \to (S, \mathcal{O}_S)$, $\mathcal{G}$ et
$c : \mathcal{J} \to \mathcal{G}$ comme dans
(\ref{equation-to-solve-ringed-spaces}).
Notons $\xi \in \Ext^1_{\mathcal{O}_S}(\NL_{S/S'}, \mathcal{J})$
l'élément correspondant, par la
Lemme \ref{lemma-extensions-of-relative-ringed-spaces}, à l'extension
$\mathcal{O}_{S'}$ de $\mathcal{O}_S$ par $\mathcal{J}$.
L'ensemble des classes d'isomorphisme de solutions est canoniquement en
bijection avec la fibre de
$$
\Ext^1_{\mathcal{O}_X}(\NL_{X/S'}, \mathcal{G}) \to
\Ext^1_{\mathcal{O}_X}(Lf^*\NL_{S/S'}, \mathcal{G})
$$
au-dessus de l'image de $\xi$.
\end{lemma}

\begin{proof}
D'après le Lemme \ref{lemma-extensions-of-relative-ringed-spaces}, appliqué
à $X \to S'$ et au $\mathcal{O}_X$-module $\mathcal{G}$, les éléments
$\zeta$ de
$\Ext^1_{\mathcal{O}_X}(\NL_{X/S'}, \mathcal{G})$
paramètrent les extensions
$0 \to \mathcal{G} \to \mathcal{O}_{X'} \to \mathcal{O}_X \to 0$
de $f^{-1}\mathcal{O}_{S'}$-algèbres.
D'après la
Lemme \ref{lemma-extensions-of-relative-ringed-spaces-functorial}, appliquée
à $X \to S \to S'$ et à $c : \mathcal{J} \to \mathcal{G}$, il existe un
$S'$-morphisme $X' \to S'$ compatible avec $c$ et
$f : X \to S$ si et seulement si $\zeta$ s'envoie sur $\xi$.
Bien entendu, cela revient à dire que $\mathcal{O}_{X'}$ est une solution de
(\ref{equation-to-solve-ringed-spaces}).
\end{proof}

\begin{remark}
\label{remark-parametrize-solutions-ringed-spaces}
Dans la situation de la
Lemme \ref{lemma-parametrize-solutions-ringed-spaces}, nous avons des
applications de complexes
$$
Lf^*\NL_{S'/S} \to \NL_{X/S'} \to \NL_{X/S}
$$
Ces applications sont proches de former un triangle distingué ; voir
Modules, Lemme \ref{modules-lemma-exact-sequence-NL-ringed-topoi}.
Si elles formaient un triangle distingué, on en conclurait que l'image de
$\xi$ dans $\Ext^2_{\mathcal{O}_X}(\NL_{X/S}, \mathcal{G})$
serait l'obstruction à l'existence d'une solution de
(\ref{equation-to-solve-ringed-spaces}).
\end{remark}

\section{Déformations de schémas}
\label{section-deformations-schemes}

\noindent
Dans cette section, nous explicitons la signification des résultats de la
section \ref{section-deformations-ringed-spaces} pour les déformations de
schémas.

\begin{lemma}
\label{lemma-deform}
Soit $S \subset S'$ un épaississement du premier ordre de schémas.
Soit $f : X \to S$ un morphisme plat de schémas.
S'il existe un morphisme plat $f' : X' \to S'$ de schémas et un isomorphisme
$a : X \to X' \times_{S'} S$ au-dessus de $S$, alors
\begin{enumerate}
\item l'ensemble des classes d'isomorphisme de couples
$(f' : X' \to S', a)$ est un ensemble principal homogène sous
$\Ext^1_{\mathcal{O}_X}(\NL_{X/S}, f^*\mathcal{C}_{S/S'})$, et
\item l'ensemble des automorphismes $\varphi : X' \to X'$ au-dessus de $S'$
qui se réduisent à l'identité sur $X' \times_{S'} S$ est
$\Ext^0_{\mathcal{O}_X}(\NL_{X/S}, f^*\mathcal{C}_{S/S'})$.
\end{enumerate}
\end{lemma}

\begin{proof}
Observons d'abord que les épaississements de schémas tels qu'ils sont définis
dans Morphismes, section \ref{more-morphisms-section-thickenings}, ne sont
autres que les morphismes de schémas qui sont des épaississements au sens de la
section \ref{section-thickenings-spaces}.
Nous pouvons considérer $X$ comme un sous-schéma fermé de $X'$ de sorte que
$(f, f') : (X \subset X') \to (S \subset S')$ soit un morphisme
d'épaississements du premier ordre. Il résulte alors de Morphismes,
Lemme \ref{more-morphisms-lemma-deform} (ou du
Lemme \ref{lemma-deform-module}, plus général) que le faisceau d'idéaux de $X$
dans $X'$ est égal à $f^*\mathcal{C}_{S/S'}$.
Nous avons donc un diagramme commutatif
$$
\xymatrix{
0 \ar[r] & f^*\mathcal{C}_{S/S'} \ar[r] &
\mathcal{O}_{X'} \ar[r] &
\mathcal{O}_X \ar[r] & 0 \\
0 \ar[r] & \mathcal{C}_{S/S'} \ar[u] \ar[r] &
\mathcal{O}_{S'} \ar[u] \ar[r] &
\mathcal{O}_S \ar[u] \ar[r] & 0
}
$$
où les flèches verticales sont des $f$-applications ; comparer avec
(\ref{equation-to-solve-ringed-spaces}).
La partie (1) découle donc du Lemme \ref{lemma-choices-ringed-spaces}, et la
partie (2), de la partie (2) du
Lemme \ref{lemma-huge-diagram-ringed-spaces}.
(Notons que $\NL_{X/S}$, tel qu'il est défini pour un morphisme de schémas dans
Morphismes, section \ref{more-morphisms-section-netherlander}, coïncide avec
$\NL_{X/S}$ tel qu'il est employé dans la
section \ref{section-deformations-ringed-spaces}.)
\end{proof}

\section{Épaississements de topos annelés}
\label{section-thickenings-ringed-topoi}

\noindent
Cette section est l'analogue de la section \ref{section-thickenings-spaces}
pour les topos annelés.
Dans les quelques sections qui suivent, nous utiliserons les notions suivantes :
\begin{enumerate}
\item Un faisceau d'idéaux $\mathcal{I} \subset \mathcal{O}'$ sur un topos
annelé $(\Sh(\mathcal{D}), \mathcal{O}')$ est {\it localement nilpotent} si
toute section locale de $\mathcal{I}$ est localement nilpotente.
\item Un {\it épaississement} de topos annelés est un morphisme
$i : (\Sh(\mathcal{C}), \mathcal{O}) \to (\Sh(\mathcal{D}), \mathcal{O}')$
de topos annelés tel que
\begin{enumerate}
\item $i_*$ soit une équivalence $\Sh(\mathcal{C}) \to \Sh(\mathcal{D})$,
\item l'application $i^\sharp : \mathcal{O}' \to i_*\mathcal{O}$ soit
surjective, et
\item le noyau de $i^\sharp$ soit un faisceau d'idéaux localement nilpotent.
\end{enumerate}
\item Un {\it épaississement du premier ordre} de topos annelés est un
épaississement
$i : (\Sh(\mathcal{C}), \mathcal{O}) \to (\Sh(\mathcal{D}), \mathcal{O}')$
de topos annelés tel que $\Ker(i^\sharp)$ soit de carré nul.
\item On voit clairement comment définir les
{\it morphismes d'épaississements de topos annelés}, les
{\it morphismes d'épaississements de topos annelés au-dessus d'un topos annelé
de base}, etc.
\end{enumerate}
Si
$i : (\Sh(\mathcal{C}), \mathcal{O}) \to (\Sh(\mathcal{D}), \mathcal{O}')$
est un épaississement de topos annelés, nous identifions les topos sous-jacents
et considérons $\mathcal{O}$, $\mathcal{O}'$ et
$\mathcal{I} = \Ker(i^\sharp)$ comme des faisceaux sur $\mathcal{C}$.
Nous obtenons une suite exacte courte
$$
0 \to \mathcal{I} \to \mathcal{O}' \to \mathcal{O} \to 0
$$
de $\mathcal{O}'$-modules. D'après Modules sur les sites,
Lemme \ref{sites-modules-lemma-i-star-equivalence}, la catégorie des
$\mathcal{O}$-modules est équivalente à la catégorie des
$\mathcal{O}'$-modules annulés par $\mathcal{I}$. En particulier, si $i$ est un
épaississement du premier ordre, alors $\mathcal{I}$ est un
$\mathcal{O}$-module.

\begin{situation}
\label{situation-morphism-thickenings-ringed-topoi}
Un morphisme d'épaississements de topos annelés $(f, f')$ est donné par un
diagramme commutatif
\begin{equation}
\label{equation-morphism-thickenings-ringed-topoi}
\vcenter{
\xymatrix{
(\Sh(\mathcal{C}), \mathcal{O}) \ar[r]_i \ar[d]_f &
(\Sh(\mathcal{D}), \mathcal{O}') \ar[d]^{f'} \\
(\Sh(\mathcal{B}), \mathcal{O}_\mathcal{B}) \ar[r]^t &
(\Sh(\mathcal{B}'), \mathcal{O}_{\mathcal{B}'})
}
}
\end{equation}
de topos annelés dont les flèches horizontales sont des épaississements. Dans
cette situation, nous posons
$\mathcal{I} = \Ker(i^\sharp) \subset \mathcal{O}'$ et
$\mathcal{J} = \Ker(t^\sharp) \subset \mathcal{O}_{\mathcal{B}'}$.
Comme $f = f'$ sur les topos sous-jacents, nous identifierons les foncteurs
image réciproque $f^{-1}$ et $(f')^{-1}$.
Observons que
$(f')^\sharp : f^{-1}\mathcal{O}_{\mathcal{B}'} \to \mathcal{O}'$
induit en particulier une application $f^{-1}\mathcal{J} \to \mathcal{I}$,
et donc une application de $\mathcal{O}'$-modules
$$
(f')^*\mathcal{J} \longrightarrow \mathcal{I}
$$
Si $i$ et $t$ sont des épaississements du premier ordre, alors
$(f')^*\mathcal{J} = f^*\mathcal{J}$ et l'application ci-dessus devient une
application $f^*\mathcal{J} \to \mathcal{I}$.
\end{situation}

\begin{definition}
\label{definition-strict-morphism-thickenings-ringed-topoi}
Dans la Situation \ref{situation-morphism-thickenings-ringed-topoi}, nous
disons que $(f, f')$ est un {\it morphisme strict d'épaississements} si
l'application $(f')^*\mathcal{J} \longrightarrow \mathcal{I}$ est surjective.
\end{definition}

\section{Modules sur les épaississements du premier ordre de topos annelés}
\label{section-modules-thickenings-ringed-topoi}

\noindent
Dans cette section, nous présentons quelques préliminaires à la théorie des
déformations des modules. Soit
$i : (\Sh(\mathcal{C}, \mathcal{O}) \to (\Sh(\mathcal{D}), \mathcal{O}')$
un épaississement du premier ordre de topos annelés. Nous utiliserons librement
les notations introduites dans la
section \ref{section-thickenings-ringed-topoi} ; en particulier, nous
identifierons les topos sous-jacents.
Dans cette section, nous considérons des suites exactes courtes
\begin{equation}
\label{equation-extension-ringed-topoi}
0 \to \mathcal{K} \to \mathcal{F}' \to \mathcal{F} \to 0
\end{equation}
de $\mathcal{O}'$-modules, où $\mathcal{F}$ et $\mathcal{K}$ sont des
$\mathcal{O}$-modules, et $\mathcal{F}'$ est un $\mathcal{O}'$-module.
Dans cette situation, nous avons une application canonique de
$\mathcal{O}$-modules
$$
c_{\mathcal{F}'} :
\mathcal{I} \otimes_\mathcal{O} \mathcal{F}
\longrightarrow
\mathcal{K}
$$
où $\mathcal{I} = \Ker(i^\sharp)$.
En effet, étant données des sections locales $f$ de $\mathcal{I}$ et $s$ de
$\mathcal{F}$, nous posons $c_{\mathcal{F}'}(f \otimes s) = fs'$, où $s'$ est
une section locale de $\mathcal{F}'$ relevant $s$.

\begin{lemma}
\label{lemma-inf-map-ringed-topoi}
Soit $i : (\Sh(\mathcal{C}), \mathcal{O}) \to (\Sh(\mathcal{D}), \mathcal{O}')$
un épaississement du premier ordre de topos annelés. Supposons données des
extensions
$$
0 \to \mathcal{K} \to \mathcal{F}' \to \mathcal{F} \to 0
\quad\text{et}\quad
0 \to \mathcal{L} \to \mathcal{G}' \to \mathcal{G} \to 0
$$
comme dans (\ref{equation-extension-ringed-topoi}), ainsi que des applications
$\varphi : \mathcal{F} \to \mathcal{G}$ et
$\psi : \mathcal{K} \to \mathcal{L}$.
\begin{enumerate}
\item S'il existe une application de $\mathcal{O}'$-modules
$\varphi' : \mathcal{F}' \to \mathcal{G}'$ compatible avec $\varphi$ et
$\psi$, alors le diagramme
$$
\xymatrix{
\mathcal{I} \otimes_\mathcal{O} \mathcal{F}
\ar[r]_-{c_{\mathcal{F}'}} \ar[d]_{1 \otimes \varphi} &
\mathcal{K} \ar[d]^\psi \\
\mathcal{I} \otimes_\mathcal{O} \mathcal{G}
\ar[r]^-{c_{\mathcal{G}'}} &
\mathcal{L}
}
$$
est commutatif.
\item L'ensemble des applications de $\mathcal{O}'$-modules
$\varphi' : \mathcal{F}' \to \mathcal{G}'$ compatibles avec $\varphi$ et
$\psi$ est, s'il est non vide, un espace principal homogène sous
$\Hom_\mathcal{O}(\mathcal{F}, \mathcal{L})$.
\end{enumerate}
\end{lemma}

\begin{proof}
La partie (1) résulte immédiatement de la description des applications.
Pour (2), si $\varphi'$ et $\varphi''$ sont deux applications
$\mathcal{F}' \to \mathcal{G}'$ compatibles avec $\varphi$ et $\psi$, alors
$\varphi' - \varphi''$ se factorise sous la forme
$$
\mathcal{F}' \to \mathcal{F} \to \mathcal{L} \to \mathcal{G}'
$$
L'application du milieu provient d'un unique élément de
$\Hom_\mathcal{O}(\mathcal{F}, \mathcal{L})$, d'après Modules sur les sites,
Lemme \ref{sites-modules-lemma-i-star-equivalence}.
Réciproquement, étant donné un élément $\alpha$ de ce groupe, nous pouvons
ajouter la composée affichée ci-dessus, avec $\alpha$ au milieu, à $\varphi'$.
Nous omettons certains détails.
\end{proof}

\begin{lemma}
\label{lemma-inf-obs-map-ringed-topoi}
Soit $i : (\Sh(\mathcal{C}), \mathcal{O}) \to (\Sh(\mathcal{D}), \mathcal{O}')$
un épaississement du premier ordre de topos annelés. Supposons données des
extensions
$$
0 \to \mathcal{K} \to \mathcal{F}' \to \mathcal{F} \to 0
\quad\text{et}\quad
0 \to \mathcal{L} \to \mathcal{G}' \to \mathcal{G} \to 0
$$
comme dans (\ref{equation-extension-ringed-topoi}), ainsi que des applications
$\varphi : \mathcal{F} \to \mathcal{G}$ et
$\psi : \mathcal{K} \to \mathcal{L}$. Supposons le diagramme
$$
\xymatrix{
\mathcal{I} \otimes_\mathcal{O} \mathcal{F}
\ar[r]_-{c_{\mathcal{F}'}} \ar[d]_{1 \otimes \varphi} &
\mathcal{K} \ar[d]^\psi \\
\mathcal{I} \otimes_\mathcal{O} \mathcal{G}
\ar[r]^-{c_{\mathcal{G}'}} &
\mathcal{L}
}
$$
commutatif. Il existe alors un élément
$$
o(\varphi, \psi) \in
\Ext^1_\mathcal{O}(\mathcal{F}, \mathcal{L})
$$
dont l'annulation est une condition nécessaire et suffisante pour qu'il existe
une application $\varphi' : \mathcal{F}' \to \mathcal{G}'$ compatible avec
$\varphi$ et $\psi$.
\end{lemma}

\begin{proof}
Nous pouvons construire explicitement une extension
$$
0 \to \mathcal{L} \to \mathcal{H} \to \mathcal{F} \to 0
$$
en prenant pour $\mathcal{H}$ la cohomologie au milieu du complexe
$$
\mathcal{K}
\xrightarrow{1, - \psi}
\mathcal{F}' \oplus \mathcal{G}' \xrightarrow{\varphi, 1}
\mathcal{G}
$$
(avec les notations évidentes). Un calcul sur les sections locales utilisant
l'hypothèse de commutativité du diagramme du lemme montre que $\mathcal{H}$ est
annulé par $\mathcal{I}$. Par conséquent, $\mathcal{H}$ définit une classe dans
$$
\Ext^1_\mathcal{O}(\mathcal{F}, \mathcal{L})
\subset
\Ext^1_{\mathcal{O}'}(\mathcal{F}, \mathcal{L})
$$
Enfin, la classe de $\mathcal{H}$ est la différence entre la somme amalgamée de
l'extension $\mathcal{F}'$ suivant $\psi$ et l'image réciproque de l'extension
$\mathcal{G}'$ suivant $\varphi$ (calculs omis).
L'annulation de la classe de $\mathcal{H}$ équivaut donc à l'existence d'un
diagramme commutatif
$$
\xymatrix{
0 \ar[r] &
\mathcal{K} \ar[r] \ar[d]_{\psi} &
\mathcal{F}' \ar[r] \ar[d]_{\varphi'} &
\mathcal{F} \ar[r] \ar[d]_\varphi & 0\\
0 \ar[r] &
\mathcal{L} \ar[r] &
\mathcal{G}' \ar[r] &
\mathcal{G} \ar[r] & 0
}
$$
comme voulu.
\end{proof}

\begin{lemma}
\label{lemma-inf-ext-ringed-topoi}
Soit $i : (\Sh(\mathcal{C}), \mathcal{O}) \to (\Sh(\mathcal{D}), \mathcal{O}')$
un épaississement du premier ordre de topos annelés. Supposons donnés des
$\mathcal{O}$-modules $\mathcal{F}$, $\mathcal{K}$ et une application
$\mathcal{O}$-linéaire
$c : \mathcal{I} \otimes_\mathcal{O} \mathcal{F} \to \mathcal{K}$.
S'il existe une suite (\ref{equation-extension-ringed-topoi}) telle que
$c_{\mathcal{F}'} = c$, alors l'ensemble des classes d'isomorphisme de ces
extensions est un espace principal homogène sous
$\Ext^1_\mathcal{O}(\mathcal{F}, \mathcal{K})$.
\end{lemma}

\begin{proof}
Supposons données des extensions
$$
0 \to \mathcal{K} \to \mathcal{F}'_1 \to \mathcal{F} \to 0
\quad\text{et}\quad
0 \to \mathcal{K} \to \mathcal{F}'_2 \to \mathcal{F} \to 0
$$
telles que $c_{\mathcal{F}'_1} = c_{\mathcal{F}'_2} = c$. Leur différence
(dans le groupe des extensions ; voir Homologie,
section \ref{homology-section-extensions}) est une extension
$$
0 \to \mathcal{K} \to \mathcal{E} \to \mathcal{F} \to 0
$$
où $\mathcal{E}$ est annulé par $\mathcal{I}$ (on omet le calcul local).
La suite est donc une extension de $\mathcal{O}$-modules ; voir Modules sur les
sites, Lemme \ref{sites-modules-lemma-i-star-equivalence}.
Réciproquement, étant donnée une telle extension $\mathcal{E}$, nous pouvons
ajouter l'extension $\mathcal{E}$ à l'extension de $\mathcal{O}'$-modules
$\mathcal{F}'$ sans modifier l'application $c_{\mathcal{F}'}$. Nous omettons
certains détails.
\end{proof}

\begin{lemma}
\label{lemma-inf-obs-ext-ringed-topoi}
Soit $i : (\Sh(\mathcal{C}), \mathcal{O}) \to (\Sh(\mathcal{D}), \mathcal{O}')$
un épaississement du premier ordre de topos annelés. Supposons donnés des
$\mathcal{O}$-modules $\mathcal{F}$, $\mathcal{K}$ et une application
$\mathcal{O}$-linéaire
$c : \mathcal{I} \otimes_\mathcal{O} \mathcal{F} \to \mathcal{K}$.
Il existe alors un élément
$$
o(\mathcal{F}, \mathcal{K}, c) \in
\Ext^2_\mathcal{O}(\mathcal{F}, \mathcal{K})
$$
dont l'annulation est une condition nécessaire et suffisante pour qu'il existe
une suite (\ref{equation-extension-ringed-topoi}) telle que
$c_{\mathcal{F}'} = c$.
\end{lemma}

\begin{proof}
Montrons d'abord que, si $\mathcal{K}$ est un $\mathcal{O}$-module injectif,
il existe bien une suite (\ref{equation-extension-ringed-topoi}) telle que
$c_{\mathcal{F}'} = c$. Pour cela, choisissons un $\mathcal{O}'$-module plat
$\mathcal{H}'$ et une surjection $\mathcal{H}' \to \mathcal{F}$
(Modules sur les sites,
Lemme \ref{sites-modules-lemma-module-quotient-flat}).
Soit $\mathcal{J} \subset \mathcal{H}'$ le noyau. Comme $\mathcal{H}'$ est
plat, nous avons
$$
\mathcal{I} \otimes_{\mathcal{O}'} \mathcal{H}' =
\mathcal{I}\mathcal{H}'
\subset \mathcal{J} \subset \mathcal{H}'
$$
Observons que l'application
$$
\mathcal{I}\mathcal{H}' =
\mathcal{I} \otimes_{\mathcal{O}'} \mathcal{H}'
\longrightarrow
\mathcal{I} \otimes_{\mathcal{O}'} \mathcal{F} =
\mathcal{I} \otimes_\mathcal{O} \mathcal{F}
$$
annule $\mathcal{I}\mathcal{J}$. En effet, si $f$ est une section locale de
$\mathcal{I}$ et $s$ une section locale de $\mathcal{H}$, alors $fs$ s'envoie
sur $f \otimes \overline{s}$, où $\overline{s}$ est l'image de $s$ dans
$\mathcal{F}$. Nous obtenons ainsi
$$
\xymatrix{
\mathcal{I}\mathcal{H}'/\mathcal{I}\mathcal{J}
\ar@{^{(}->}[r] \ar[d] &
\mathcal{J}/\mathcal{I}\mathcal{J} \ar@{..>}[d]_\gamma \\
\mathcal{I} \otimes_\mathcal{O} \mathcal{F} \ar[r]^-c &
\mathcal{K}
}
$$
un diagramme de $\mathcal{O}$-modules. Si $\mathcal{K}$ est injectif comme
$\mathcal{O}$-module, nous obtenons la flèche pointillée.
Notons $\gamma' : \mathcal{J} \to \mathcal{K}$ la composée de $\gamma$ avec
$\mathcal{J} \to \mathcal{J}/\mathcal{I}\mathcal{J}$.
Un calcul local montre que la somme amalgamée
$$
\xymatrix{
0 \ar[r] &
\mathcal{J} \ar[r] \ar[d]_{\gamma'} &
\mathcal{H}' \ar[r] \ar[d] &
\mathcal{F} \ar[r] \ar@{=}[d] &
0 \\
0 \ar[r] &
\mathcal{K} \ar[r] &
\mathcal{F}' \ar[r] &
\mathcal{F} \ar[r] &
0
}
$$
est une solution au problème posé par le lemme.

\medskip\noindent
Cas général. Choisissons un plongement $\mathcal{K} \subset \mathcal{K}'$,
où $\mathcal{K}'$ est un $\mathcal{O}$-module injectif. Soit $\mathcal{Q}$ le
quotient, de sorte que nous ayons une suite exacte
$$
0 \to \mathcal{K} \to \mathcal{K}' \to \mathcal{Q} \to 0
$$
Notons
$c' : \mathcal{I} \otimes_\mathcal{O} \mathcal{F} \to \mathcal{K}'$
la composée. Le paragraphe précédent fournit une suite
$$
0 \to \mathcal{K}' \to \mathcal{E}' \to \mathcal{F} \to 0
$$
comme dans (\ref{equation-extension-ringed-topoi}), telle que
$c_{\mathcal{E}'} = c'$.
Remarquons que la composée de $c'$ avec l'application
$\mathcal{K}' \to \mathcal{Q}$ est nulle ; la somme amalgamée de
$\mathcal{E}'$ suivant $\mathcal{K}' \to \mathcal{Q}$ est donc une extension
$$
0 \to \mathcal{Q} \to \mathcal{D}' \to \mathcal{F} \to 0
$$
comme dans (\ref{equation-extension-ringed-topoi}), telle que
$c_{\mathcal{D}'} = 0$.
Cela signifie exactement que $\mathcal{D}'$ est annulé par $\mathcal{I}$ ;
autrement dit, $\mathcal{D}'$ est une extension de $\mathcal{O}$-modules,
c'est-à-dire qu'il définit un élément
$$
o(\mathcal{F}, \mathcal{K}, c) \in
\Ext^1_\mathcal{O}(\mathcal{F}, \mathcal{Q}) =
\Ext^2_\mathcal{O}(\mathcal{F}, \mathcal{K})
$$
(l'égalité résulte de la suite exacte longue de cohomologie associée à la suite
exacte précédente et de l'annulation des groupes Ext supérieurs à valeurs dans
le module injectif $\mathcal{K}'$). Si
$o(\mathcal{F}, \mathcal{K}, c) = 0$, nous pouvons choisir un scindage
$s : \mathcal{F} \to \mathcal{D}'$ et poser
$$
\mathcal{F}' = \Ker(\mathcal{E}' \to \mathcal{D}'/s(\mathcal{F}))
$$
de sorte que nous obtenions le diagramme suivant
$$
\xymatrix{
0 \ar[r] &
\mathcal{K} \ar[r] \ar[d] &
\mathcal{F}' \ar[r] \ar[d] &
\mathcal{F} \ar[r] \ar@{=}[d] &
0 \\
0 \ar[r] &
\mathcal{K}' \ar[r] &
\mathcal{E}' \ar[r] &
\mathcal{F} \ar[r] & 0
}
$$
à lignes exactes, qui montre que $c_{\mathcal{F}'} = c$. Réciproquement, si
$\mathcal{F}'$ existe, alors la somme amalgamée de $\mathcal{F}'$ suivant
l'application $\mathcal{K} \to \mathcal{K}'$ est isomorphe à $\mathcal{E}'$,
d'après le Lemme \ref{lemma-inf-ext-ringed-topoi} et l'annulation des groupes
Ext supérieurs à valeurs dans le module injectif $\mathcal{K}'$.
Nous obtenons ainsi un diagramme comme ci-dessus, ce qui implique que
$\mathcal{D}'$ est scindée comme extension, c'est-à-dire que la classe
$o(\mathcal{F}, \mathcal{K}, c)$ est nulle.
\end{proof}

\begin{remark}
\label{remark-trivial-thickening-ringed-topoi}
Soit $(\Sh(\mathcal{C}), \mathcal{O})$ un topos annelé. Un épaississement du
premier ordre $i : (\Sh(\mathcal{C}), \mathcal{O}) \to
(\Sh(\mathcal{D}), \mathcal{O}')$ est dit {\it trivial} s'il existe un
morphisme de topos annelés
$\pi : (\Sh(\mathcal{D}), \mathcal{O}') \to (\Sh(\mathcal{C}), \mathcal{O})$
qui soit un inverse à gauche de $i$. Le choix d'un tel morphisme $\pi$ est
appelé une {\it trivialisation} de l'épaississement du premier ordre.
Étant donné $\pi$, nous obtenons un scindage
\begin{equation}
\label{equation-splitting-ringed-topoi}
\mathcal{O}' = \mathcal{O} \oplus \mathcal{I}
\end{equation}
comme faisceaux d'algèbres sur $\mathcal{C}$, en utilisant $\pi^\sharp$ pour
scinder la surjection $\mathcal{O}' \to \mathcal{O}$.
Réciproquement, un tel scindage détermine un morphisme $\pi$. La catégorie des
épaississements du premier ordre trivialisés de
$(\Sh(\mathcal{C}), \mathcal{O})$ est équivalente à la catégorie des
$\mathcal{O}$-modules.
\end{remark}

\begin{remark}
\label{remark-trivial-extension-ringed-topoi}
Soit $i : (\Sh(\mathcal{C}), \mathcal{O}) \to (\Sh(\mathcal{D}), \mathcal{O}')$
un épaississement trivial du premier ordre de topos annelés, et soit
$\pi : (\Sh(\mathcal{D}), \mathcal{O}') \to
(\Sh(\mathcal{C}), \mathcal{O})$ une trivialisation. Pour tout triplet
$(\mathcal{F}, \mathcal{K}, c)$ formé d'une paire de $\mathcal{O}$-modules et
d'une application
$c : \mathcal{I} \otimes_\mathcal{O} \mathcal{F} \to \mathcal{K}$,
nous pouvons poser
$$
\mathcal{F}'_{c, triv} = \mathcal{F} \oplus \mathcal{K}
$$
et utiliser le scindage (\ref{equation-splitting-ringed-topoi}) associé à
$\pi$, ainsi que l'application $c$, pour définir la structure de
$\mathcal{O}'$-module et obtenir une extension
(\ref{equation-extension-ringed-topoi}).
Nous appellerons $\mathcal{F}'_{c, triv}$ l'{\it extension triviale} de
$\mathcal{F}$ par $\mathcal{K}$ correspondant à $c$ et à la trivialisation
$\pi$. Pour toute extension $\mathcal{F}'$ comme dans
(\ref{equation-extension-ringed-topoi}), nous pouvons utiliser
$\pi^\sharp : \mathcal{O} \to \mathcal{O}'$ pour considérer $\mathcal{F}'$
comme une extension de $\mathcal{O}$-modules, et donc comme une classe
$\xi_{\mathcal{F}'}$ dans
$\Ext^1_\mathcal{O}(\mathcal{F}, \mathcal{K})$.
Le Lemme \ref{lemma-inf-ext-ringed-topoi} assure que
$\mathcal{F}' \mapsto \xi_{\mathcal{F}'}$ induit une bijection
$$
\left\{
\begin{matrix}
\text{isomorphism classes of extensions}\\
\mathcal{F}'\text{ comme dans (\ref{equation-extension-ringed-topoi}) avec }
c = c_{\mathcal{F}'}
\end{matrix}
\right\}
\longrightarrow
\Ext^1_\mathcal{O}(\mathcal{F}, \mathcal{K})
$$
De plus, l'extension triviale $\mathcal{F}'_{c, triv}$ s'envoie sur la classe
nulle.
\end{remark}

\begin{remark}
\label{remark-extension-functorial-ringed-topoi}
Soit $(\Sh(\mathcal{C}), \mathcal{O})$ un topos annelé. Soient
$(\Sh(\mathcal{C}), \mathcal{O}) \to (\Sh(\mathcal{D}_i), \mathcal{O}'_i)$,
$i = 1, 2$, des épaississements du
premier ordre de faisceaux d'idéaux $\mathcal{I}_i$.
Soit $h : (\Sh(\mathcal{D}_1), \mathcal{O}'_1) \to
(\Sh(\mathcal{D}_2), \mathcal{O}'_2)$ un morphisme d'épaississements du
premier ordre de $(\Sh(\mathcal{C}), \mathcal{O})$.
Voici le diagramme
$$
\xymatrix{
& (\Sh(\mathcal{C}), \mathcal{O}) \ar[ld] \ar[rd] & \\
(\Sh(\mathcal{D}_1), \mathcal{O}'_1) \ar[rr]^h & & 
(\Sh(\mathcal{D}_2), \mathcal{O}'_2)
}
$$
Remarquons que $h^\sharp : \mathcal{O}'_2 \to \mathcal{O}'_1$ induit en
particulier une application de $\mathcal{O}$-modules
$\mathcal{I}_2 \to \mathcal{I}_1$.
Soit $\mathcal{F}$ un $\mathcal{O}$-module.
Soit $(\mathcal{K}_i, c_i)$, $i = 1, 2$, une paire formée d'un
$\mathcal{O}$-module $\mathcal{K}_i$ et d'une application
$c_i : \mathcal{I}_i \otimes_\mathcal{O} \mathcal{F} \to
\mathcal{K}_i$. Supposons en outre donnée une application de
$\mathcal{O}$-modules $\mathcal{K}_2 \to \mathcal{K}_1$ telle que le diagramme
$$
\xymatrix{
\mathcal{I}_2 \otimes_\mathcal{O} \mathcal{F}
\ar[r]_-{c_2} \ar[d] &
\mathcal{K}_2 \ar[d] \\
\mathcal{I}_1 \otimes_\mathcal{O} \mathcal{F}
\ar[r]^-{c_1} &
\mathcal{K}_1
}
$$
soit commutatif. Nous disposons alors d'une fonctorialité canonique
$$
\left\{
\begin{matrix}
\mathcal{F}'_2\text{ comme dans (\ref{equation-extension-ringed-topoi}) avec }\\
c_2 = c_{\mathcal{F}'_2}\text{ et }\mathcal{K} = \mathcal{K}_2
\end{matrix}
\right\}
\longrightarrow
\left\{
\begin{matrix}
\mathcal{F}'_1\text{ comme dans (\ref{equation-extension-ringed-topoi}) avec }\\
c_1 = c_{\mathcal{F}'_1}\text{ et }\mathcal{K} = \mathcal{K}_1
\end{matrix}
\right\}
$$
En effet, en considérant tous les faisceaux $\mathcal{O}$,
$\mathcal{O}'_i$, $\mathcal{F}$, $\mathcal{K}_i$, etc., comme des faisceaux
sur $\mathcal{C}$, nous associons à $\mathcal{F}'_2$ le faisceau
$\mathcal{F}'_1$ défini comme la somme amalgamée, c'est-à-dire s'insérant dans
le diagramme d'extensions suivant
$$
\xymatrix{
0 \ar[r] &
\mathcal{K}_2 \ar[r] \ar[d] &
\mathcal{F}'_2 \ar[r] \ar[d] &
\mathcal{F} \ar@{=}[d] \ar[r] & 0 \\
0 \ar[r] &
\mathcal{K}_1 \ar[r] &
\mathcal{F}'_1 \ar[r] &
\mathcal{F} \ar[r] & 0
}
$$
Nous omettons la construction de la structure de $\mathcal{O}'_1$-module sur
la somme amalgamée (elle utilise la commutativité du diagramme faisant
intervenir $c_1$ et $c_2$).
\end{remark}

\begin{remark}
\label{remark-trivial-extension-functorial-ringed-topoi}
Soient $(\Sh(\mathcal{C}), \mathcal{O})$,
$(\Sh(\mathcal{C}), \mathcal{O}) \to (\Sh(\mathcal{D}_i), \mathcal{O}'_i)$,
$\mathcal{I}_i$, et
$h : (\Sh(\mathcal{D}_1), \mathcal{O}'_1) \to
(\Sh(\mathcal{D}_2), \mathcal{O}'_2)$ comme dans la
Remarque \ref{remark-extension-functorial-ringed-topoi}.
Supposons données des trivialisations
$\pi_i : (\Sh(\mathcal{D}_i), \mathcal{O}'_i) \to
(\Sh(\mathcal{C}), \mathcal{O})$ telles que
$\pi_1 = h \circ \pi_2$. Autrement dit, supposons que $h$ soit un morphisme
d'épaississements du premier ordre trivialisés de
$(\Sh(\mathcal{C}), \mathcal{O})$.
Soit $(\mathcal{K}_i, c_i)$, $i = 1, 2$, une paire formée d'un
$\mathcal{O}$-module $\mathcal{K}_i$ et d'une application
$c_i : \mathcal{I}_i \otimes_\mathcal{O} \mathcal{F} \to
\mathcal{K}_i$. Supposons en outre donnée une application de
$\mathcal{O}$-modules $\mathcal{K}_2 \to \mathcal{K}_1$ telle que le diagramme
$$
\xymatrix{
\mathcal{I}_2 \otimes_\mathcal{O} \mathcal{F}
\ar[r]_-{c_2} \ar[d] &
\mathcal{K}_2 \ar[d] \\
\mathcal{I}_1 \otimes_\mathcal{O} \mathcal{F}
\ar[r]^-{c_1} &
\mathcal{K}_1
}
$$
soit commutatif. Dans cette situation, la construction de la
Remarque \ref{remark-trivial-extension-ringed-topoi} induit un diagramme
commutatif
$$
\xymatrix{
\{\mathcal{F}'_2\text{ comme dans (\ref{equation-extension-ringed-topoi}) avec }
c_2 = c_{\mathcal{F}'_2}\text{ et }\mathcal{K} = \mathcal{K}_2\}
\ar[d] \ar[rr] & &
\Ext^1_\mathcal{O}(\mathcal{F}, \mathcal{K}_2) \ar[d] \\
\{\mathcal{F}'_1\text{ comme dans (\ref{equation-extension-ringed-topoi}) avec }
c_1 = c_{\mathcal{F}'_1}\text{ et }\mathcal{K} = \mathcal{K}_1\}
\ar[rr] & &
\Ext^1_\mathcal{O}(\mathcal{F}, \mathcal{K}_1)
}
$$
où l'application verticale de droite est donnée par la fonctorialité de $\Ext$
et l'application $\mathcal{K}_2 \to \mathcal{K}_1$, tandis que l'application
verticale de gauche est celle de la
Remarque \ref{remark-extension-functorial-ringed-topoi}.
\end{remark}

\begin{remark}
\label{remark-obstruction-extension-functorial-ringed-topoi}
Soient $(\Sh(\mathcal{C}), \mathcal{O})$,
$(\Sh(\mathcal{C}), \mathcal{O}) \to (\Sh(\mathcal{D}_i), \mathcal{O}'_i)$,
$\mathcal{I}_i$, et
$h : (\Sh(\mathcal{D}_1), \mathcal{O}'_1) \to
(\Sh(\mathcal{D}_2), \mathcal{O}'_2)$ comme dans la
Remarque \ref{remark-extension-functorial-ringed-topoi}.
Observons que $h^\sharp : \mathcal{O}'_2 \to \mathcal{O}'_1$ induit en
particulier une application de $\mathcal{O}$-modules
$\mathcal{I}_2 \to \mathcal{I}_1$.
Soit $\mathcal{F}$ un $\mathcal{O}$-module.
Soit $(\mathcal{K}_i, c_i)$, $i = 1, 2$, une paire formée d'un
$\mathcal{O}$-module $\mathcal{K}_i$ et d'une application
$c_i : \mathcal{I}_i \otimes_\mathcal{O} \mathcal{F} \to
\mathcal{K}_i$. Supposons en outre donnée une application de
$\mathcal{O}$-modules $\mathcal{K}_2 \to \mathcal{K}_1$ telle que le diagramme
$$
\xymatrix{
\mathcal{I}_2 \otimes_\mathcal{O} \mathcal{F}
\ar[r]_-{c_2} \ar[d] &
\mathcal{K}_2 \ar[d] \\
\mathcal{I}_1 \otimes_\mathcal{O} \mathcal{F}
\ar[r]^-{c_1} &
\mathcal{K}_1
}
$$
soit commutatif. Nous {\bf affirmons} alors que l'application
$$
\Ext^2_\mathcal{O}(\mathcal{F}, \mathcal{K}_2)
\longrightarrow
\Ext^2_\mathcal{O}(\mathcal{F}, \mathcal{K}_1)
$$
envoie $o(\mathcal{F}, \mathcal{K}_2, c_2)$ sur
$o(\mathcal{F}, \mathcal{K}_1, c_1)$.

\medskip\noindent
Pour démontrer cette affirmation, choisissons un plongement
$j_2 : \mathcal{K}_2 \to \mathcal{K}_2'$, où $\mathcal{K}_2'$ est un
$\mathcal{O}$-module injectif.
Comme dans la démonstration du
Lemme \ref{lemma-inf-obs-ext-ringed-topoi}, nous pouvons choisir une extension
de $\mathcal{O}_2$-modules
$$
0 \to \mathcal{K}_2' \to \mathcal{E}_2 \to \mathcal{F} \to 0
$$
telle que $c_{\mathcal{E}_2} = j_2 \circ c_2$.
La démonstration du Lemme \ref{lemma-inf-obs-ext-ringed-topoi} construit
$o(\mathcal{F}, \mathcal{K}_2, c_2)$ comme la classe de l'extension de Yoneda
(au sens de Catégories dérivées, section \ref{derived-section-ext}) associée à
la suite exacte de $\mathcal{O}$-modules
$$
0 \to
\mathcal{K}_2 \to \mathcal{K}_2' \to
\mathcal{E}_2/\mathcal{K}_2 \to
\mathcal{F} \to 0
$$
Soit $\mathcal{K}_1'$ le conoyau de
$\mathcal{K}_2 \to \mathcal{K}_1 \oplus \mathcal{K}_2'$.
Il existe une injection $j_1 : \mathcal{K}_1 \to \mathcal{K}_1'$ et une
application $\mathcal{K}_2' \to \mathcal{K}_1'$ formant un carré commutatif.
Nous formons la somme amalgamée
$$
\xymatrix{
0 \ar[r] &
\mathcal{K}_2' \ar[r] \ar[d] &
\mathcal{E}_2 \ar[r] \ar[d] &
\mathcal{F} \ar[r] \ar[d] & 0 \\
0 \ar[r] &
\mathcal{K}_1' \ar[r] &
\mathcal{E}_1 \ar[r] &
\mathcal{F} \ar[r] & 0
}
$$
Il existe une structure canonique de $\mathcal{O}_1$-module sur
$\mathcal{E}_1$, et, pour cette structure, nous avons
$c_{\mathcal{E}_1} = j_1 \circ c_1$ (cela utilise la commutativité du diagramme
ci-dessus faisant intervenir $c_1$ et $c_2$).
Le procédé du Lemme \ref{lemma-inf-obs-ext-ringed-topoi} nous dit que
$o(\mathcal{F}, \mathcal{K}_1, c_1)$ est la classe de l'extension de Yoneda
associée à la suite exacte de $\mathcal{O}$-modules
$$
0 \to
\mathcal{K}_1 \to
\mathcal{K}_1' \to
\mathcal{E}_1/\mathcal{K}_1 \to
\mathcal{F} \to 0
$$
Puisque nous avons des applications de suites exactes
$$
\xymatrix{
0 \ar[r] &
\mathcal{K}_2 \ar[d] \ar[r] &
\mathcal{K}_2' \ar[d] \ar[r] &
\mathcal{E}_2/\mathcal{K}_2 \ar[r] \ar[d] &
\mathcal{F} \ar[r] \ar@{=}[d] &
0 \\
0 \ar[r] &
\mathcal{K}_2 \ar[r] &
\mathcal{K}_2' \ar[r] &
\mathcal{E}_2/\mathcal{K}_2 \ar[r] &
\mathcal{F} \ar[r] &
0
}
$$
nous concluons que l'affirmation est vraie.
\end{remark}

\begin{remark}
\label{remark-short-exact-sequence-thickenings-ringed-topoi}
Soit $(\Sh(\mathcal{C}), \mathcal{O})$ un topos annelé.
Nous disons qu'une suite de morphismes d'épaississements du premier ordre
$$
(\Sh(\mathcal{D}_1), \mathcal{O}'_1) \to
(\Sh(\mathcal{D}_2), \mathcal{O}'_2) \to
(\Sh(\mathcal{D}_3), \mathcal{O}'_3)
$$
de $(\Sh(\mathcal{C}), \mathcal{O})$ est un {\it complexe} si les applications
correspondantes entre les faisceaux d'idéaux $\mathcal{I}_i$ forment un
complexe de $\mathcal{O}$-modules
$\mathcal{I}_3 \to \mathcal{I}_2 \to \mathcal{I}_1$ (c'est-à-dire si la
composée est nulle). Dans ce cas, la composée
$(\Sh(\mathcal{D}_1), \mathcal{O}'_1) \to
(\Sh(\mathcal{D}_3), \mathcal{O}'_3)$ se factorise par
$(\Sh(\mathcal{C}), \mathcal{O}) \to
(\Sh(\mathcal{D}_3), \mathcal{O}'_3)$ ; autrement dit, l'épaississement du
premier ordre $(\Sh(\mathcal{D}_1), \mathcal{O}'_1)$ de
$(\Sh(\mathcal{C}), \mathcal{O})$ est trivial et muni d'une trivialisation
canonique
$\pi : (\Sh(\mathcal{D}_1), \mathcal{O}'_1) \to
(\Sh(\mathcal{C}), \mathcal{O})$.

\medskip\noindent
Nous disons qu'une suite de morphismes d'épaississements du premier ordre
$$
(\Sh(\mathcal{D}_1), \mathcal{O}'_1) \to
(\Sh(\mathcal{D}_2), \mathcal{O}'_2) \to
(\Sh(\mathcal{D}_3), \mathcal{O}'_3)
$$
de $(\Sh(\mathcal{C}), \mathcal{O})$ est une {\it suite exacte courte} si les
applications correspondantes entre les faisceaux d'idéaux forment une suite
exacte courte
$$
0 \to \mathcal{I}_3 \to \mathcal{I}_2 \to \mathcal{I}_1 \to 0
$$
de $\mathcal{O}$-modules.
\end{remark}

\begin{remark}
\label{remark-complex-thickenings-and-ses-modules-ringed-topoi}
Soit $(\Sh(\mathcal{C}), \mathcal{O})$ un topos annelé. Soit $\mathcal{F}$ un
$\mathcal{O}$-module. Soit
$$
(\Sh(\mathcal{D}_1), \mathcal{O}'_1) \to
(\Sh(\mathcal{D}_2), \mathcal{O}'_2) \to
(\Sh(\mathcal{D}_3), \mathcal{O}'_3)
$$
un complexe d'épaississements du premier ordre de
$(\Sh(\mathcal{C}), \mathcal{O})$ ; voir la
Remarque \ref{remark-short-exact-sequence-thickenings-ringed-topoi}.
Soient $(\mathcal{K}_i, c_i)$, $i = 1, 2, 3$, des couples formés d'un
$\mathcal{O}$-module $\mathcal{K}_i$ et d'une application
$c_i : \mathcal{I}_i \otimes_\mathcal{O} \mathcal{F} \to
\mathcal{K}_i$. Supposons donnée une suite exacte courte de
$\mathcal{O}$-modules
$$
0 \to \mathcal{K}_3 \to \mathcal{K}_2 \to \mathcal{K}_1 \to 0
$$
telle que
$$
\vcenter{
\xymatrix{
\mathcal{I}_2 \otimes_\mathcal{O} \mathcal{F}
\ar[r]_-{c_2} \ar[d] &
\mathcal{K}_2 \ar[d] \\
\mathcal{I}_1 \otimes_\mathcal{O} \mathcal{F}
\ar[r]^-{c_1} &
\mathcal{K}_1
}
}
\quad\text{et}\quad
\vcenter{
\xymatrix{
\mathcal{I}_3 \otimes_\mathcal{O} \mathcal{F}
\ar[r]_-{c_3} \ar[d] &
\mathcal{K}_3 \ar[d] \\
\mathcal{I}_2 \otimes_\mathcal{O} \mathcal{F}
\ar[r]^-{c_2} &
\mathcal{K}_2
}
}
$$
soient commutatifs. Supposons enfin donnée une extension
$$
0 \to \mathcal{K}_2 \to \mathcal{F}'_2 \to \mathcal{F} \to 0
$$
comme dans (\ref{equation-extension-ringed-topoi}), avec
$\mathcal{K} = \mathcal{K}_2$, de $\mathcal{O}'_2$-modules telle que
$c_{\mathcal{F}'_2} = c_2$.
Dans cette situation, nous pouvons appliquer la fonctorialité de la
Remarque \ref{remark-extension-functorial-ringed-topoi} pour obtenir une
extension $\mathcal{F}'_1$ de $\mathcal{O}'_1$-modules (nous décrirons
ci-dessous $\mathcal{F}'_1$ dans ce cas particulier). D'après la
Remarque \ref{remark-trivial-extension-ringed-topoi}, en utilisant le scindage
canonique
$\pi : (\Sh(\mathcal{D}_1), \mathcal{O}'_1) \to
(\Sh(\mathcal{C}), \mathcal{O})$ de la
Remarque \ref{remark-short-exact-sequence-thickenings-ringed-topoi}, nous
obtenons
$\xi_{\mathcal{F}'_1} \in
\Ext^1_\mathcal{O}(\mathcal{F}, \mathcal{K}_1)$.
Enfin, nous avons l'obstruction
$$
o(\mathcal{F}, \mathcal{K}_3, c_3) \in
\Ext^2_\mathcal{O}(\mathcal{F}, \mathcal{K}_3)
$$
du Lemme \ref{lemma-inf-obs-ext-ringed-topoi}.
Dans cette situation, nous {\bf affirmons} que l'application canonique
$$
\partial :
\Ext^1_\mathcal{O}(\mathcal{F}, \mathcal{K}_1)
\longrightarrow
\Ext^2_\mathcal{O}(\mathcal{F}, \mathcal{K}_3)
$$
provenant de la suite exacte courte
$0 \to \mathcal{K}_3 \to \mathcal{K}_2 \to \mathcal{K}_1 \to 0$ envoie
$\xi_{\mathcal{F}'_1}$ sur la classe d'obstruction
$o(\mathcal{F}, \mathcal{K}_3, c_3)$.

\medskip\noindent
Pour démontrer cette affirmation, choisissons un plongement
$j : \mathcal{K}_3 \to \mathcal{K}$, où $\mathcal{K}$ est un
$\mathcal{O}$-module injectif.
Nous pouvons relever $j$ en une application
$j' : \mathcal{K}_2 \to \mathcal{K}$.
Posons $\mathcal{E}'_2 = j'_*\mathcal{F}'_2$, égal à la somme amalgamée de
$\mathcal{F}'_2$ par $j'$, de sorte que
$c_{\mathcal{E}'_2} = j' \circ c_2$. Voici le diagramme
$$
\xymatrix{
0 \ar[r] &
\mathcal{K}_2 \ar[r] \ar[d]_{j'} &
\mathcal{F}'_2 \ar[r] \ar[d] &
\mathcal{F} \ar[r] \ar[d] & 0 \\
0 \ar[r] &
\mathcal{K} \ar[r] &
\mathcal{E}'_2 \ar[r] &
\mathcal{F} \ar[r] & 0
}
$$
Posons $\mathcal{E}'_3 = \mathcal{E}'_2$, mais considérons-le comme un
$\mathcal{O}'_3$-module au moyen de $\mathcal{O}'_3 \to \mathcal{O}'_2$.
Alors $c_{\mathcal{E}'_3} = j \circ c_3$.
La démonstration du Lemme \ref{lemma-inf-obs-ext-ringed-topoi} construit
$o(\mathcal{F}, \mathcal{K}_3, c_3)$ comme le cobord de la classe de
l'extension de $\mathcal{O}$-modules
$$
0 \to
\mathcal{K}/\mathcal{K}_3 \to
\mathcal{E}'_3/\mathcal{K}_3 \to
\mathcal{F} \to 0
$$
D'autre part, remarquons que
$\mathcal{F}'_1 = \mathcal{F}'_2/\mathcal{K}_3$ ; ainsi, la classe
$\xi_{\mathcal{F}'_1}$ est la classe de l'extension
$$
0 \to \mathcal{K}_2/\mathcal{K}_3 \to \mathcal{F}'_2/\mathcal{K}_3
\to \mathcal{F} \to 0
$$
considérée comme une suite de $\mathcal{O}$-modules au moyen de $\pi^\sharp$,
où
$\pi : (\Sh(\mathcal{D}_1), \mathcal{O}'_1) \to
(\Sh(\mathcal{C}), \mathcal{O})$ est le scindage canonique.
Ainsi, l'affirmation résulte finalement de l'existence du diagramme commutatif
$$
\xymatrix{
0 \ar[r] &
\mathcal{K}_2/\mathcal{K}_3 \ar[r] \ar[d] &
\mathcal{F}'_2/\mathcal{K}_3 \ar[r] \ar[d] &
\mathcal{F} \ar[r] \ar[d] & 0 \\
0 \ar[r] &
\mathcal{K}/\mathcal{K}_3 \ar[r] &
\mathcal{E}'_3/\mathcal{K}_3 \ar[r] &
\mathcal{F} \ar[r] & 0
}
$$
qui est $\mathcal{O}$-linéaire (pour les structures de $\mathcal{O}$-modules
données ci-dessus).
\end{remark}

\section{Déformations infinitésimales de modules sur des topos annelés}
\sectionmark{Déformations de modules sur des topos annelés}
\label{section-deformation-modules-ringed-topoi}

\noindent
Soit $i : (\Sh(\mathcal{C}), \mathcal{O}) \to (\Sh(\mathcal{D}), \mathcal{O}')$
un épaississement du premier ordre de topos annelés. Nous employons librement
les notations introduites dans la
section \ref{section-thickenings-ringed-topoi}.
Soit $\mathcal{F}'$ un $\mathcal{O}'$-module et posons
$\mathcal{F} = i^*\mathcal{F}'$.
Dans cette situation, nous avons une suite exacte courte
$$
0 \to \mathcal{I}\mathcal{F}' \to \mathcal{F}' \to \mathcal{F} \to 0
$$
de $\mathcal{O}'$-modules. Puisque $\mathcal{I}^2 = 0$, la structure de
$\mathcal{O}'$-module sur $\mathcal{I}\mathcal{F}'$ provient d'une unique
structure de $\mathcal{O}$-module.
Ainsi, la suite ci-dessus est une extension comme dans
(\ref{equation-extension-ringed-topoi}).
Dans le cas particulier où $\mathcal{F}' = \mathcal{O}'$, nous avons
$i^*\mathcal{O}' = \mathcal{O}$ et
$\mathcal{I}\mathcal{O}' = \mathcal{I}$, et nous retrouvons la suite des
faisceaux structuraux
$$
0 \to \mathcal{I} \to \mathcal{O}' \to \mathcal{O} \to 0
$$

\begin{lemma}
\label{lemma-inf-map-special-ringed-topoi}
Soit $i : (\Sh(\mathcal{C}), \mathcal{O}) \to (\Sh(\mathcal{D}), \mathcal{O}')$
un épaississement du premier ordre de topos annelés.
Soient $\mathcal{F}'$, $\mathcal{G}'$ des $\mathcal{O}'$-modules.
Posons $\mathcal{F} = i^*\mathcal{F}'$ et
$\mathcal{G} = i^*\mathcal{G}'$.
Soit $\varphi : \mathcal{F} \to \mathcal{G}$ une application
$\mathcal{O}$-linéaire.
L'ensemble des relèvements de $\varphi$ en une application
$\mathcal{O}'$-linéaire $\varphi' : \mathcal{F}' \to \mathcal{G}'$ est, s'il
est non vide, un espace principal homogène sous
$\Hom_\mathcal{O}(\mathcal{F}, \mathcal{I}\mathcal{G}')$.
\end{lemma}

\begin{proof}
C'est un cas particulier du Lemme \ref{lemma-inf-map-ringed-topoi}, mais nous
en donnons aussi une démonstration directe. Nous avons des suites exactes
courtes de modules
$$
0 \to \mathcal{I} \to \mathcal{O}' \to \mathcal{O} \to 0
\quad\text{et}\quad
0 \to \mathcal{I}\mathcal{G}' \to \mathcal{G}' \to \mathcal{G} \to 0
$$
et de même pour $\mathcal{F}'$.
Comme $\mathcal{I}$ est de carré nul, les structures de $\mathcal{O}'$-modules
sur $\mathcal{I}$ et $\mathcal{I}\mathcal{G}'$ proviennent d'uniques structures
de $\mathcal{O}$-modules. Il s'ensuit que
$$
\Hom_{\mathcal{O}'}(\mathcal{F}', \mathcal{I}\mathcal{G}') =
\Hom_\mathcal{O}(\mathcal{F}, \mathcal{I}\mathcal{G}')
\quad\text{et}\quad
\Hom_{\mathcal{O}'}(\mathcal{F}', \mathcal{G}) =
\Hom_\mathcal{O}(\mathcal{F}, \mathcal{G})
$$
Le lemme résulte alors de la suite exacte
$$
0 \to \Hom_{\mathcal{O}'}(\mathcal{F}', \mathcal{I}\mathcal{G}') \to
\Hom_{\mathcal{O}'}(\mathcal{F}', \mathcal{G}') \to
\Hom_{\mathcal{O}'}(\mathcal{F}', \mathcal{G})
$$
voir Homologie, Lemme \ref{homology-lemma-check-exactness}.
\end{proof}

\begin{lemma}
\label{lemma-deform-module-ringed-topoi}
Soit $(f, f')$ un morphisme d'épaississements du premier ordre de topos annelés
comme dans la
Situation \ref{situation-morphism-thickenings-ringed-topoi}.
Soit $\mathcal{F}'$ un $\mathcal{O}'$-module et posons
$\mathcal{F} = i^*\mathcal{F}'$.
Supposons que $\mathcal{F}$ soit plat sur $\mathcal{O}_\mathcal{B}$ et que
$(f, f')$ soit un morphisme strict d'épaississements
(Définition \ref{definition-strict-morphism-thickenings-ringed-topoi}).
Alors les assertions suivantes sont équivalentes :
\begin{enumerate}
\item $\mathcal{F}'$ est plat sur $\mathcal{O}_{\mathcal{B}'}$, et
\item l'application canonique
$f^*\mathcal{J} \otimes_\mathcal{O} \mathcal{F} \to
\mathcal{I}\mathcal{F}'$
est un isomorphisme.
\end{enumerate}
De plus, dans ce cas, les applications
$$
f^*\mathcal{J} \otimes_\mathcal{O} \mathcal{F} \to
\mathcal{I} \otimes_\mathcal{O} \mathcal{F} \to
\mathcal{I}\mathcal{F}'
$$
sont des isomorphismes.
\end{lemma}

\begin{proof}
L'application $f^*\mathcal{J} \to \mathcal{I}$ est surjective, car
$(f, f')$ est un morphisme strict d'épaississements.
La dernière assertion résulte donc de (2).

\medskip\noindent
Démontrons l'équivalence de (1) et (2). Par définition, la platitude sur
$\mathcal{O}_\mathcal{B}$ signifie la platitude sur
$f^{-1}\mathcal{O}_\mathcal{B}$. De même pour la platitude sur
$f^{-1}\mathcal{O}_{\mathcal{B}'}$.
Remarquons que la stricte compatibilité de $(f, f')$ et l'hypothèse
$\mathcal{F} = i^*\mathcal{F}'$ impliquent que
$$
\mathcal{F} = \mathcal{F}'/(f^{-1}\mathcal{J})\mathcal{F}'
$$
comme faisceaux sur $\mathcal{C}$. En outre, observons que
$f^*\mathcal{J} \otimes_\mathcal{O} \mathcal{F} =
f^{-1}\mathcal{J} \otimes_{f^{-1}\mathcal{O}_\mathcal{B}} \mathcal{F}$.
L'équivalence de (1) et (2) résulte donc de Modules sur les sites,
Lemme \ref{sites-modules-lemma-flat-over-thickening}.
\end{proof}

\begin{lemma}
\label{lemma-deform-fp-module-ringed-topoi}
Soit $(f, f')$ un morphisme d'épaississements du premier ordre de topos annelés
comme dans la
Situation \ref{situation-morphism-thickenings-ringed-topoi}.
Soit $\mathcal{F}'$ un $\mathcal{O}'$-module et posons
$\mathcal{F} = i^*\mathcal{F}'$.
Supposons que $\mathcal{F}'$ soit plat sur
$\mathcal{O}_{\mathcal{B}'}$ et que $(f, f')$ soit un morphisme strict
d'épaississements. Alors les assertions suivantes sont équivalentes :
\begin{enumerate}
\item $\mathcal{F}'$ est un $\mathcal{O}'$-module de présentation finie, et
\item $\mathcal{F}$ est un $\mathcal{O}$-module de présentation finie.
\end{enumerate}
\end{lemma}

\begin{proof}
L'implication (1) $\Rightarrow$ (2) résulte de Modules sur les sites,
Lemme \ref{sites-modules-lemma-local-pullback}.
Pour la réciproque, supposons $\mathcal{F}$ de présentation finie.
Nous pouvons et allons supposer que $\mathcal{C} = \mathcal{C}'$.
D'après le Lemme \ref{lemma-deform-module-ringed-topoi}, nous avons une suite
exacte courte
$$
0 \to \mathcal{I} \otimes_{\mathcal{O}_X} \mathcal{F} \to
\mathcal{F}' \to \mathcal{F} \to 0
$$
Soit $U$ un objet de $\mathcal{C}$ tel que $\mathcal{F}|_U$ admette une
présentation
$$
\mathcal{O}_U^{\oplus m} \to \mathcal{O}_U^{\oplus n} \to \mathcal{F}|_U \to 0
$$
Quitte à remplacer $U$ par les membres d'un recouvrement, nous pouvons
supposer que l'application $\mathcal{O}_U^{\oplus n} \to \mathcal{F}|_U$ se
relève en une application
$(\mathcal{O}'_U)^{\oplus n} \to \mathcal{F}'|_U$. L'application induite
$\mathcal{I}^{\oplus n} \to \mathcal{I} \otimes \mathcal{F}$ est surjective
par exactitude à droite de $\otimes$. Ainsi, quitte à remplacer encore $U$ par
les membres d'un recouvrement, nous pouvons trouver un relèvement
$(\mathcal{O}'|_U)^{\oplus m} \to (\mathcal{O}'|_U)^{\oplus n}$ de
l'application donnée
$\mathcal{O}_U^{\oplus m} \to \mathcal{O}_U^{\oplus n}$ tel que
$$
(\mathcal{O}'_U)^{\oplus m} \to (\mathcal{O}'_U)^{\oplus n} \to
\mathcal{F}'|_U \to 0
$$
soit un complexe. Une nouvelle application de l'exactitude à droite de
$\otimes$ montre que ce complexe est exact.
\end{proof}

\begin{lemma}
\label{lemma-inf-map-rel-ringed-topoi}
Soit $(f, f')$ un morphisme d'épaississements du premier ordre comme dans la
Situation \ref{situation-morphism-thickenings-ringed-topoi}.
Soient $\mathcal{F}'$, $\mathcal{G}'$ des $\mathcal{O}'$-modules et posons
$\mathcal{F} = i^*\mathcal{F}'$ et $\mathcal{G} = i^*\mathcal{G}'$.
Soit $\varphi : \mathcal{F} \to \mathcal{G}$ une application
$\mathcal{O}$-linéaire.
Supposons que $\mathcal{G}'$ soit plat sur
$\mathcal{O}_{\mathcal{B}'}$ et que $(f, f')$ soit un morphisme strict
d'épaississements.
L'ensemble des relèvements de $\varphi$ en une application
$\mathcal{O}'$-linéaire $\varphi' : \mathcal{F}' \to \mathcal{G}'$ est, s'il
est non vide, un espace principal homogène sous
$$
\Hom_\mathcal{O}(\mathcal{F},
\mathcal{G} \otimes_\mathcal{O} f^*\mathcal{J})
$$
\end{lemma}

\begin{proof}
Combiner les Lemmes \ref{lemma-inf-map-special-ringed-topoi} et
\ref{lemma-deform-module-ringed-topoi}.
\end{proof}

\begin{lemma}
\label{lemma-inf-obs-map-special-ringed-topoi}
Soit $i : (\Sh(\mathcal{C}), \mathcal{O}) \to (\Sh(\mathcal{D}), \mathcal{O}')$
un épaississement du premier ordre de topos annelés.
Soient $\mathcal{F}'$, $\mathcal{G}'$ des $\mathcal{O}'$-modules et posons
$\mathcal{F} = i^*\mathcal{F}'$ et $\mathcal{G} = i^*\mathcal{G}'$.
Soit $\varphi : \mathcal{F} \to \mathcal{G}$ une application
$\mathcal{O}$-linéaire.
Il existe un élément
$$
o(\varphi) \in
\Ext^1_\mathcal{O}(Li^*\mathcal{F}', \mathcal{I}\mathcal{G}')
$$
dont l'annulation est une condition nécessaire et suffisante pour l'existence
d'un relèvement de $\varphi$ en une application $\mathcal{O}'$-linéaire
$\varphi' : \mathcal{F}' \to \mathcal{G}'$.
\end{lemma}

\begin{proof}
Il ressort de la démonstration du
Lemme \ref{lemma-inf-map-special-ringed-topoi} que l'annulation du cobord de
$\varphi$ par l'application
$$
\Hom_\mathcal{O}(\mathcal{F}, \mathcal{G}) =
\Hom_{\mathcal{O}'}(\mathcal{F}', \mathcal{G}) \longrightarrow
\Ext^1_{\mathcal{O}'}(\mathcal{F}', \mathcal{I}\mathcal{G}')
$$
est une condition nécessaire et suffisante pour l'existence d'un relèvement.
Nous concluons puisque
$$
\Ext^1_{\mathcal{O}'}(\mathcal{F}', \mathcal{I}\mathcal{G}') =
\Ext^1_\mathcal{O}(Li^*\mathcal{F}', \mathcal{I}\mathcal{G}')
$$
par l'adjonction de $i_* = Ri_*$ et $Li^*$ dans la catégorie dérivée
(Cohomologie sur les sites, Lemme \ref{sites-cohomology-lemma-adjoint}).
\end{proof}

\begin{lemma}
\label{lemma-inf-obs-map-rel-ringed-topoi}
Soit $(f, f')$ un morphisme d'épaississements du premier ordre comme dans la
Situation \ref{situation-morphism-thickenings-ringed-topoi}.
Soient $\mathcal{F}'$, $\mathcal{G}'$ des $\mathcal{O}'$-modules et posons
$\mathcal{F} = i^*\mathcal{F}'$ et $\mathcal{G} = i^*\mathcal{G}'$.
Soit $\varphi : \mathcal{F} \to \mathcal{G}$ une application
$\mathcal{O}$-linéaire.
Supposons que $\mathcal{F}'$ et $\mathcal{G}'$ soient plats sur
$\mathcal{O}_{\mathcal{B}'}$ et que $(f, f')$ soit un morphisme strict
d'épaississements. Il existe un élément
$$
o(\varphi) \in
\Ext^1_\mathcal{O}(\mathcal{F},
\mathcal{G} \otimes_\mathcal{O} f^*\mathcal{J})
$$
dont l'annulation est une condition nécessaire et suffisante pour l'existence
d'un relèvement de $\varphi$ en une application $\mathcal{O}'$-linéaire
$\varphi' : \mathcal{F}' \to \mathcal{G}'$.
\end{lemma}

\begin{proof}[Première démonstration]
Cela résulte du Lemme \ref{lemma-inf-obs-map-special-ringed-topoi}, car nous
affirmons que, sous les hypothèses du lemme, nous avons
$$
\Ext^1_\mathcal{O}(Li^*\mathcal{F}', \mathcal{I}\mathcal{G}') =
\Ext^1_\mathcal{O}(\mathcal{F},
\mathcal{G} \otimes_\mathcal{O} f^*\mathcal{J})
$$
En effet, nous avons
$\mathcal{I}\mathcal{G}' =
\mathcal{G} \otimes_\mathcal{O} f^*\mathcal{J}$ d'après le
Lemme \ref{lemma-deform-module-ringed-topoi}.
D'autre part, observons que
$$
H^{-1}(Li^*\mathcal{F}') =
\text{Tor}_1^{\mathcal{O}'}(\mathcal{F}', \mathcal{O})
$$
(nous omettons le calcul local). En utilisant la suite exacte courte
$$
0 \to \mathcal{I} \to \mathcal{O}' \to \mathcal{O} \to 0
$$
nous voyons que ce $\text{Tor}_1$ est calculé par le noyau de l'application
$\mathcal{I} \otimes_\mathcal{O} \mathcal{F} \to \mathcal{I}\mathcal{F}'$,
qui est nul d'après la dernière assertion du
Lemme \ref{lemma-deform-module-ringed-topoi}.
Ainsi, $\tau_{\geq -1}Li^*\mathcal{F}' = \mathcal{F}$.
D'autre part, nous avons
$$
\Ext^1_\mathcal{O}(Li^*\mathcal{F}',
\mathcal{I}\mathcal{G}') =
\Ext^1_\mathcal{O}(\tau_{\geq -1}Li^*\mathcal{F}',
\mathcal{I}\mathcal{G}')
$$
par le dual de Catégories dérivées,
Lemme \ref{derived-lemma-negative-vanishing}.
\end{proof}

\begin{proof}[Seconde démonstration]
Nous pouvons appliquer le Lemme \ref{lemma-inf-obs-map-ringed-topoi} comme
suit. Remarquons que
$\mathcal{K} = \mathcal{I} \otimes_\mathcal{O} \mathcal{F}$ et
$\mathcal{L} = \mathcal{I} \otimes_\mathcal{O} \mathcal{G}$ d'après le
Lemme \ref{lemma-deform-module-ringed-topoi}, que
$c_{\mathcal{F}'} = 1 \otimes 1$ et $c_{\mathcal{G}'} = 1 \otimes 1$, et
qu'en prenant $\psi = 1 \otimes \varphi$, le diagramme du lemme est commutatif.
Ainsi, $o(\varphi) = o(\varphi, 1 \otimes \varphi)$ convient.
\end{proof}

\begin{lemma}
\label{lemma-inf-ext-rel-ringed-topoi}
Soit $(f, f')$ un morphisme d'épaississements du premier ordre comme dans la
Situation \ref{situation-morphism-thickenings-ringed-topoi}.
Soit $\mathcal{F}$ un $\mathcal{O}$-module.
Supposons que $(f, f')$ soit un morphisme strict d'épaississements et que
$\mathcal{F}$ soit plat sur $\mathcal{O}_\mathcal{B}$. S'il existe un couple
$(\mathcal{F}', \alpha)$ formé d'un $\mathcal{O}'$-module $\mathcal{F}'$ plat
sur $\mathcal{O}_{\mathcal{B}'}$ et d'un isomorphisme
$\alpha : i^*\mathcal{F}' \to \mathcal{F}$, alors l'ensemble des classes
d'isomorphisme de tels couples est un espace principal homogène sous
$\Ext^1_\mathcal{O}(
\mathcal{F}, \mathcal{I} \otimes_\mathcal{O} \mathcal{F})$.
\end{lemma}

\begin{proof}
Si nous supposons qu'il existe un tel module, alors l'application canonique
$$
f^*\mathcal{J} \otimes_\mathcal{O} \mathcal{F} \to
\mathcal{I} \otimes_\mathcal{O} \mathcal{F}
$$
est un isomorphisme d'après le
Lemme \ref{lemma-deform-module-ringed-topoi}. Appliquons le
Lemme \ref{lemma-inf-ext-ringed-topoi} avec $\mathcal{K} = 
\mathcal{I} \otimes_\mathcal{O} \mathcal{F}$ et $c = 1$.
D'après le Lemme \ref{lemma-deform-module-ringed-topoi}, les extensions
$\mathcal{F}'$ correspondantes sont toutes plates sur
$\mathcal{O}_{\mathcal{B}'}$.
\end{proof}

\begin{lemma}
\label{lemma-inf-obs-ext-rel-ringed-topoi}
Soit $(f, f')$ un morphisme d'épaississements du premier ordre comme dans la
Situation \ref{situation-morphism-thickenings-ringed-topoi}.
Soit $\mathcal{F}$ un $\mathcal{O}$-module. Supposons que $(f, f')$ soit un
morphisme strict d'épaississements et que $\mathcal{F}$ soit plat sur
$\mathcal{O}_\mathcal{B}$. Il existe un $\mathcal{O}'$-module $\mathcal{F}'$
plat sur $\mathcal{O}_{\mathcal{B}'}$ tel que
$i^*\mathcal{F}' \cong \mathcal{F}$ si et seulement si
\begin{enumerate}
\item l'application canonique
$f^*\mathcal{J} \otimes_\mathcal{O} \mathcal{F} \to
\mathcal{I} \otimes_\mathcal{O} \mathcal{F}$
est un isomorphisme, et
\item la classe
$o(\mathcal{F}, \mathcal{I} \otimes_\mathcal{O} \mathcal{F}, 1)
\in \Ext^2_\mathcal{O}(
\mathcal{F}, \mathcal{I} \otimes_\mathcal{O} \mathcal{F})$
du Lemme \ref{lemma-inf-obs-ext-ringed-topoi} est nulle.
\end{enumerate}
\end{lemma}

\begin{proof}
Cela résulte immédiatement de la caractérisation des $\mathcal{O}'$-modules
plats sur $\mathcal{O}_{\mathcal{B}'}$ donnée par le
Lemme \ref{lemma-deform-module-ringed-topoi} et du
Lemme \ref{lemma-inf-obs-ext-ringed-topoi}.
\end{proof}

\section{Application aux modules plats sur les épaississements plats de topos annelés}
\label{section-flat-ringed-topoi}

\noindent
Considérons un diagramme commutatif
$$
\xymatrix{
(\Sh(\mathcal{C}), \mathcal{O}) \ar[r]_i \ar[d]_f &
(\Sh(\mathcal{D}), \mathcal{O}') \ar[d]^{f'} \\
(\Sh(\mathcal{B}), \mathcal{O}_\mathcal{B}) \ar[r]^t &
(\Sh(\mathcal{B}'), \mathcal{O}_{\mathcal{B}'})
}
$$
de topos annelés dont les flèches horizontales sont des épaississements du
premier ordre, comme dans la
Situation \ref{situation-morphism-thickenings-ringed-topoi}. Posons
$\mathcal{I} = \Ker(i^\sharp) \subset \mathcal{O}'$ et
$\mathcal{J} = \Ker(t^\sharp) \subset \mathcal{O}_{\mathcal{B}'}$.
Soit $\mathcal{F}$ un $\mathcal{O}$-module. Supposons que
\begin{enumerate}
\item $(f, f')$ soit un morphisme strict d'épaississements,
\item $f'$ soit plat, et
\item $\mathcal{F}$ soit plat sur $\mathcal{O}_\mathcal{B}$.
\end{enumerate}
Remarquons que (1) $+$ (2) impliquent que
$\mathcal{I} = f^*\mathcal{J}$ (appliquer le
Lemme \ref{lemma-deform-module-ringed-topoi} à $\mathcal{O}'$).
La théorie de la section précédente prend une forme particulièrement simple
sous ces hypothèses. Nous résumons les résultats déjà obtenus dans le lemme
suivant.

\begin{lemma}
\label{lemma-flat-ringed-topoi}
Dans la situation ci-dessus, les assertions suivantes valent.
\begin{enumerate}
\item Il existe un $\mathcal{O}'$-module $\mathcal{F}'$ plat sur
$\mathcal{O}_{\mathcal{B}'}$ tel que
$i^*\mathcal{F}' \cong \mathcal{F}$ si et seulement si la classe
$o(\mathcal{F}, f^*\mathcal{J} \otimes_\mathcal{O} \mathcal{F}, 1)
\in \Ext^2_\mathcal{O}(
\mathcal{F}, f^*\mathcal{J} \otimes_\mathcal{O} \mathcal{F})$
du Lemme \ref{lemma-inf-obs-ext-ringed-topoi} est nulle.
\item S'il existe un tel module, alors l'ensemble des classes d'isomorphisme
des relèvements est un espace principal homogène sous
$\Ext^1_\mathcal{O}(
\mathcal{F}, f^*\mathcal{J} \otimes_\mathcal{O} \mathcal{F})$.
\item Étant donné un relèvement $\mathcal{F}'$, l'ensemble des automorphismes de
$\mathcal{F}'$ dont l'image inverse est $\text{id}_\mathcal{F}$ est
canoniquement isomorphe à
$\Ext^0_\mathcal{O}(
\mathcal{F}, f^*\mathcal{J} \otimes_\mathcal{O} \mathcal{F})$.
\end{enumerate}
\end{lemma}

\begin{proof}
La partie (1) résulte du
Lemme \ref{lemma-inf-obs-ext-rel-ringed-topoi}, car nous avons vu ci-dessus que
$\mathcal{I} = f^*\mathcal{J}$.
La partie (2) résulte du Lemme \ref{lemma-inf-ext-rel-ringed-topoi}.
La partie (3) résulte du Lemme \ref{lemma-inf-map-rel-ringed-topoi}.
\end{proof}

\begin{situation}
\label{situation-morphism-flat-thickenings-ringed-topoi}
Soit $f : (\Sh(\mathcal{C}), \mathcal{O}) \to
(\Sh(\mathcal{B}), \mathcal{O}_\mathcal{B})$ un morphisme de topos annelés.
Considérons un diagramme commutatif
$$
\xymatrix{
(\Sh(\mathcal{C}'_1), \mathcal{O}'_1) \ar[r]_h \ar[d]_{f'_1} &
(\Sh(\mathcal{C}'_2), \mathcal{O}'_2) \ar[d]_{f'_2} \\
(\Sh(\mathcal{B}'_1), \mathcal{O}_{\mathcal{B}'_1}) \ar[r] &
(\Sh(\mathcal{B}'_2), \mathcal{O}_{\mathcal{B}'_2})
}
$$
où $h$ est un morphisme d'épaississements du premier ordre de
$(\Sh(\mathcal{C}), \mathcal{O})$, la flèche horizontale inférieure est un
morphisme d'épaississements du premier ordre de
$(\Sh(\mathcal{B}), \mathcal{O}_\mathcal{B})$, chaque $f'_i$ se restreint à
$f$, les deux couples $(f, f_i')$ sont des morphismes stricts
d'épaississements, et les deux $f'_i$ sont plats. Enfin, soit $\mathcal{F}$ un
$\mathcal{O}$-module plat sur $\mathcal{O}_\mathcal{B}$.
\end{situation}

\begin{lemma}
\label{lemma-functorial-ringed-topoi}
Dans la Situation \ref{situation-morphism-flat-thickenings-ringed-topoi}, la
classe d'obstruction
$o(\mathcal{F}, f^*\mathcal{J}_2 \otimes_\mathcal{O} \mathcal{F}, 1)$
s'envoie sur la classe d'obstruction
$o(\mathcal{F}, f^*\mathcal{J}_1 \otimes_\mathcal{O} \mathcal{F}, 1)$
par l'application canonique
$$
\Ext^2_\mathcal{O}(
\mathcal{F}, f^*\mathcal{J}_2 \otimes_\mathcal{O} \mathcal{F})
\to \Ext^2_\mathcal{O}(
\mathcal{F}, f^*\mathcal{J}_1 \otimes_\mathcal{O} \mathcal{F})
$$
\end{lemma}

\begin{proof}
Cela résulte de la
Remarque \ref{remark-obstruction-extension-functorial-ringed-topoi}.
\end{proof}

\begin{situation}
\label{situation-ses-flat-thickenings-ringed-topoi}
Soit $f : (\Sh(\mathcal{C}), \mathcal{O}) \to
(\Sh(\mathcal{B}), \mathcal{O}_\mathcal{B})$ un morphisme de topos annelés.
Considérons un diagramme commutatif
$$
\xymatrix{
(\Sh(\mathcal{C}'_1), \mathcal{O}'_1) \ar[r]_h \ar[d]_{f'_1} &
(\Sh(\mathcal{C}'_2), \mathcal{O}'_2) \ar[r] \ar[d]_{f'_2} &
(\Sh(\mathcal{C}'_3), \mathcal{O}'_3) \ar[d]_{f'_3} \\
(\Sh(\mathcal{B}'_1), \mathcal{O}_{\mathcal{B}'_1}) \ar[r] &
(\Sh(\mathcal{B}'_2), \mathcal{O}_{\mathcal{B}'_2}) \ar[r] &
(\Sh(\mathcal{B}'_3), \mathcal{O}_{\mathcal{B}'_3})
}
$$
où (a) la ligne supérieure est une suite exacte courte d'épaississements du
premier ordre de $(\Sh(\mathcal{C}), \mathcal{O})$, (b) la ligne inférieure
est une suite exacte courte d'épaississements du premier ordre de
$(\Sh(\mathcal{B}), \mathcal{O}_\mathcal{B})$, (c) chaque $f'_i$ se restreint
à $f$, (d) chaque couple $(f, f_i')$ est un morphisme strict
d'épaississements et (e) chaque $f'_i$ est plat. Enfin, soit
$\mathcal{F}'_2$ un $\mathcal{O}'_2$-module plat sur
$\mathcal{O}_{\mathcal{B}'_2}$ et posons
$\mathcal{F} = \mathcal{F}'_2 \otimes \mathcal{O}$. Soit
$\pi : (\Sh(\mathcal{C}'_1), \mathcal{O}'_1) \to
(\Sh(\mathcal{C}), \mathcal{O})$ le scindage canonique
(Remarque \ref{remark-short-exact-sequence-thickenings-ringed-topoi}).
\end{situation}

\begin{lemma}
\label{lemma-verify-iv-ringed-topoi}
Dans la Situation \ref{situation-ses-flat-thickenings-ringed-topoi}, les
modules $\pi^*\mathcal{F}$ et $h^*\mathcal{F}'_2$ sont des
$\mathcal{O}'_1$-modules plats sur $\mathcal{O}_{\mathcal{B}'_1}$ qui se
restreignent à $\mathcal{F}$ sur $(\Sh(\mathcal{C}), \mathcal{O})$.
Leur différence (Lemme \ref{lemma-flat-ringed-topoi}) est un élément $\theta$
de
$\Ext^1_\mathcal{O}(\mathcal{F},
f^*\mathcal{J}_1 \otimes_\mathcal{O} \mathcal{F})$
dont le cobord dans
$\Ext^2_\mathcal{O}(\mathcal{F},
f^*\mathcal{J}_3 \otimes_\mathcal{O} \mathcal{F})$
est égal à l'obstruction (Lemme \ref{lemma-flat-ringed-topoi}) au relèvement de
$\mathcal{F}$ en un $\mathcal{O}'_3$-module plat sur
$\mathcal{O}_{\mathcal{B}'_3}$.
\end{lemma}

\begin{proof}
Remarquons que $\pi^*\mathcal{F}$ et $h^*\mathcal{F}'_2$ se restreignent tous
deux à $\mathcal{F}$ sur $(\Sh(\mathcal{C}), \mathcal{O})$, et que les noyaux
de $\pi^*\mathcal{F} \to \mathcal{F}$ et
$h^*\mathcal{F}'_2 \to \mathcal{F}$ sont donnés par
$f^*\mathcal{J}_1 \otimes_\mathcal{O} \mathcal{F}$.
Ils sont donc plats d'après le
Lemme \ref{lemma-deform-module-ringed-topoi}.
Prendre le cobord a un sens, car la suite de modules
$$
0 \to f^*\mathcal{J}_3 \otimes_\mathcal{O} \mathcal{F} \to
f^*\mathcal{J}_2 \otimes_\mathcal{O} \mathcal{F} \to
f^*\mathcal{J}_1 \otimes_\mathcal{O} \mathcal{F} \to 0
$$
est exacte courte par les hypothèses de la
Situation \ref{situation-ses-flat-thickenings-ringed-topoi} et le fait que
$\mathcal{F}$ est plat sur $\mathcal{O}_\mathcal{B}$.
L'assertion sur la classe d'obstruction est une traduction directe du résultat
de la Remarque \ref{remark-complex-thickenings-and-ses-modules-ringed-topoi}
dans cette situation particulière.
\end{proof}

\section{Déformations de topos annelés et complexe cotangent naïf}
\label{section-deformations-ringed-topoi}

\noindent
Dans cette section, nous utilisons le complexe cotangent naïf pour faire un peu
de théorie des déformations. Nous partons d'un épaississement du premier ordre
$t : (\Sh(\mathcal{B}), \mathcal{O}_\mathcal{B}) \to
(\Sh(\mathcal{B}'), \mathcal{O}_{\mathcal{B}'})$ de topos annelés.
Nous notons $\mathcal{J} = \Ker(t^\sharp)$ et identifions les topos sous-jacents
de $\mathcal{B}$ et $\mathcal{B}'$.
Supposons en outre donnés un morphisme de topos annelés
$f : (\Sh(\mathcal{C}), \mathcal{O}) \to
(\Sh(\mathcal{B}), \mathcal{O}_\mathcal{B})$, un $\mathcal{O}$-module
$\mathcal{G}$ et une application $f^{-1}\mathcal{J} \to \mathcal{G}$ de
faisceaux de $f^{-1}\mathcal{O}_\mathcal{B}$-modules.
Dans cette section, nous nous demandons si nous pouvons trouver le point
d'interrogation qui complète le diagramme suivant
\begin{equation}
\label{equation-to-solve-ringed-topoi}
\vcenter{
\xymatrix{
0 \ar[r] & \mathcal{G} \ar[r] & {?} \ar[r] & \mathcal{O} \ar[r] & 0 \\
0 \ar[r] & f^{-1}\mathcal{J} \ar[u]^c \ar[r] &
f^{-1}\mathcal{O}_{\mathcal{B}'} \ar[u] \ar[r] &
f^{-1}\mathcal{O}_\mathcal{B} \ar[u] \ar[r] & 0
}
}
\end{equation}
et, de plus, à quel point la solution est unique (si elle existe). Plus
précisément, nous cherchons un épaississement du premier ordre
$i : (\Sh(\mathcal{C}), \mathcal{O}) \to (\Sh(\mathcal{C}'), \mathcal{O}')$
et un morphisme d'épaississements $(f, f')$ comme dans
(\ref{equation-morphism-thickenings-ringed-topoi}), où
$\Ker(i^\sharp)$ est identifié à $\mathcal{G}$, tels que $(f')^\sharp$ induise
l'application donnée $c$.
Nous dirons que $(\Sh(\mathcal{C}'), \mathcal{O}')$ est une {\it solution} de
(\ref{equation-to-solve-ringed-topoi}).


\begin{lemma}
\label{lemma-huge-diagram-ringed-topoi}
Supposons donné un diagramme commutatif de morphismes de topos annelés
\begin{equation}
\label{equation-huge-1-ringed-topoi}
\vcenter{
\xymatrix{
& (\Sh(\mathcal{C}_2), \mathcal{O}_2) \ar[r]_{i_2} \ar[d]_{f_2} \ar[ddl]_g &
(\Sh(\mathcal{C}'_2), \mathcal{O}'_2) \ar[d]^{f'_2} \\
&
(\Sh(\mathcal{B}_2), \mathcal{O}_{\mathcal{B}_2}) \ar[r]^{t_2} \ar[ddl]|\hole &
(\Sh(\mathcal{B}'_2), \mathcal{O}_{\mathcal{B}'_2}) \ar[ddl] \\
(\Sh(\mathcal{C}_1), \mathcal{O}_1) \ar[r]_{i_1} \ar[d]_{f_1} &
(\Sh(\mathcal{C}'_1), \mathcal{O}'_1) \ar[d]^{f'_1} \\
(\Sh(\mathcal{B}_1), \mathcal{O}_{\mathcal{B}_1}) \ar[r]^{t_1} &
(\Sh(\mathcal{B}'_1), \mathcal{O}_{\mathcal{B}'_1})
}
}
\end{equation}
dont les flèches horizontales sont des épaississements du premier ordre.
Posons $\mathcal{G}_j = \Ker(i_j^\sharp)$ et supposons donnée une application
de $g^{-1}\mathcal{O}_1$-modules
$\nu : g^{-1}\mathcal{G}_1 \to \mathcal{G}_2$ qui donne lieu au diagramme
commutatif
\begin{equation}
\label{equation-huge-2-ringed-topoi}
\vcenter{
\xymatrix{
& 0 \ar[r] & \mathcal{G}_2 \ar[r] &
\mathcal{O}'_2 \ar[r] &
\mathcal{O}_2 \ar[r] & 0 \\
& 0 \ar[r]|\hole &
f_2^{-1}\mathcal{J}_2 \ar[u]_{c_2} \ar[r] &
f_2^{-1}\mathcal{O}_{\mathcal{B}'_2} \ar[u] \ar[r]|\hole &
f_2^{-1}\mathcal{O}_{\mathcal{B}_2} \ar[u] \ar[r] & 0 \\
0 \ar[r] &
\mathcal{G}_1 \ar[ruu] \ar[r] &
\mathcal{O}'_1 \ar[r] &
\mathcal{O}_1 \ar[ruu] \ar[r] & 0 \\
0 \ar[r] &
f_1^{-1}\mathcal{J}_1 \ar[ruu]|\hole \ar[u]^{c_1} \ar[r] &
f_1^{-1}\mathcal{O}_{\mathcal{B}'_1} \ar[ruu]|\hole \ar[u] \ar[r] &
f_1^{-1}\mathcal{O}_{\mathcal{B}_1} \ar[ruu]|\hole \ar[u] \ar[r] & 0
}
}
\end{equation}
dont les faces avant et arrière sont des solutions de
(\ref{equation-to-solve-ringed-topoi}).
(Les flèches vers le nord-nord-ouest sont des applications sur
$\mathcal{C}_2$ après application de $g^{-1}$ à la source.)
\begin{enumerate}
\item Il existe un élément canonique de
$\Ext^1_{\mathcal{O}_2}(
Lg^*\NL_{\mathcal{O}_1/\mathcal{O}_{\mathcal{B}_1}}, \mathcal{G}_2)$
dont l'annulation est une condition nécessaire et suffisante pour qu'il existe
un morphisme de topos annelés
$(\Sh(\mathcal{C}'_2), \mathcal{O}'_2) \to
(\Sh(\mathcal{C}'_1), \mathcal{O}'_1)$ s'insérant dans
(\ref{equation-huge-1-ringed-topoi}) de façon compatible avec $\nu$.
\item S'il existe un morphisme
$(\Sh(\mathcal{C}'_2), \mathcal{O}'_2) \to
(\Sh(\mathcal{C}'_1), \mathcal{O}'_1)$ s'insérant dans
(\ref{equation-huge-1-ringed-topoi}) de façon compatible avec $\nu$, l'ensemble
de tous ces morphismes est un espace principal homogène sous
$$
\Hom_{\mathcal{O}_1}(
\Omega_{\mathcal{O}_1/\mathcal{O}_{\mathcal{B}_1}}, g_*\mathcal{G}_2) =
\Hom_{\mathcal{O}_2}(
g^*\Omega_{\mathcal{O}_1/\mathcal{O}_{\mathcal{B}_1}}, \mathcal{G}_2) =
\Ext^0_{\mathcal{O}_2}(
Lg^*\NL_{\mathcal{O}_1/\mathcal{O}_{\mathcal{B}_1}}, \mathcal{G}_2).
$$
\end{enumerate}
\end{lemma}

\begin{proof}
La démonstration de ce lemme est identique à celle du
Lemme \ref{lemma-huge-diagram-ringed-spaces}.
Nous conseillons vivement au lecteur de lire cette dernière plutôt que celle
qui suit. Nous identifierons les topos sous-jacents de tous les
épaississements considérés (nous avons déjà employé cette convention dans
l'énoncé). Les égalités de la dernière assertion du lemme résultent
immédiatement des définitions. Nous travaillerons donc avec les groupes
$\Ext^k_{\mathcal{O}_2}(
Lg^*\NL_{\mathcal{O}_1/\mathcal{O}_{\mathcal{B}_1}}, \mathcal{G}_2)$,
$k = 0, 1$ dans la suite de la démonstration. Montrons d'abord que nous pouvons
nous ramener au cas où le topos sous-jacent de tous les topos annelés du lemme
est le même.

\medskip\noindent
Pour cela, observons que
$g^{-1}\NL_{\mathcal{O}_1/\mathcal{O}_{\mathcal{B}_1}}$ est égal au complexe
cotangent naïf de l'homomorphisme de faisceaux d'anneaux
$g^{-1}f_1^{-1}\mathcal{O}_{\mathcal{B}_1} \to g^{-1}\mathcal{O}_1$ ; voir
Modules sur les sites,
Lemme \ref{sites-modules-lemma-pullback-differentials}.
De plus, le terme de degré $0$ de
$\NL_{\mathcal{O}_1/\mathcal{O}_{\mathcal{B}_1}}$ est un
$\mathcal{O}_1$-module plat ; l'application canonique
$$
Lg^*\NL_{\mathcal{O}_1/\mathcal{O}_{\mathcal{B}_1}}
\longrightarrow
g^{-1}\NL_{\mathcal{O}_1/\mathcal{O}_{\mathcal{B}_1}}
\otimes_{g^{-1}\mathcal{O}_1} \mathcal{O}_2
$$
induit donc un isomorphisme sur les faisceaux de cohomologie en degrés $0$ et
$-1$. Nous pouvons ainsi remplacer les groupes Ext du lemme par
$$
\Ext^k_{g^{-1}\mathcal{O}_1}(
g^{-1}\NL_{\mathcal{O}_1/\mathcal{O}_{\mathcal{B}_1}}, \mathcal{G}_2) =
\Ext^k_{g^{-1}\mathcal{O}_1}(
\NL_{g^{-1}\mathcal{O}_1/g^{-1}f_1^{-1}\mathcal{O}_{\mathcal{B}_1}},
\mathcal{G}_2)
$$
L'ensemble des morphismes de topos annelés
$(\Sh(\mathcal{C}'_2), \mathcal{O}'_2) \to
(\Sh(\mathcal{C}'_1), \mathcal{O}'_1)$ s'insérant dans
(\ref{equation-huge-1-ringed-topoi}) de façon compatible avec $\nu$ est en
bijection avec l'ensemble des homomorphismes de
$g^{-1}f_1^{-1}\mathcal{O}_{\mathcal{B}'_1}$-algèbres
$g^{-1}\mathcal{O}'_1 \to \mathcal{O}'_2$ compatibles avec $f^\sharp$ et
$\nu$. Nous voyons ainsi que nous pouvons supposer donné un diagramme
(\ref{equation-huge-2-ringed-topoi}) de faisceaux sur un site $\mathcal{C}$
(avec $f_1 = f_2 = \text{id}$ sur les topos sous-jacents) et chercher un
homomorphisme de faisceaux d'anneaux $\mathcal{O}'_1 \to \mathcal{O}'_2$ qui
s'y insère.

\medskip\noindent
Dans la suite de la démonstration du lemme, nous supposons que tous les topos
sous-jacents sont les mêmes, c'est-à-dire que nous avons un diagramme
(\ref{equation-huge-2-ringed-topoi}) de faisceaux sur un site $\mathcal{C}$
(avec $f_1 = f_2 = \text{id}$ sur les topos sous-jacents) et que nous cherchons
des homomorphismes de faisceaux d'anneaux
$\mathcal{O}'_1 \to \mathcal{O}'_2$ s'y insérant.
Comme groupes Ext, nous utiliserons
$\Ext^k_{\mathcal{O}_1}(
\NL_{\mathcal{O}_1/\mathcal{O}_{\mathcal{B}_1}}, \mathcal{G}_2)$,
$k = 0, 1$.

\medskip\noindent
Étape 1. Construction de la classe d'obstruction. Considérons le faisceau
d'ensembles
$$
\mathcal{E} = \mathcal{O}'_1 \times_{\mathcal{O}_2} \mathcal{O}'_2
$$
Il est muni d'une application surjective
$\alpha : \mathcal{E} \to \mathcal{O}_1$ ; nous pouvons donc employer
$\NL(\alpha)$ au lieu de
$\NL_{\mathcal{O}_1/\mathcal{O}_{\mathcal{B}_1}}$ ; voir Modules sur les sites,
Lemme \ref{sites-modules-lemma-NL-up-to-qis}.
Posons
$$
\mathcal{I}' =
\Ker(\mathcal{O}_{\mathcal{B}'_1}[\mathcal{E}] \to \mathcal{O}_1)
\quad\text{et}\quad
\mathcal{I} =
\Ker(\mathcal{O}_{\mathcal{B}_1}[\mathcal{E}] \to \mathcal{O}_1)
$$
Il existe une surjection $\mathcal{I}' \to \mathcal{I}$ dont le noyau est
$\mathcal{J}_1\mathcal{O}_{\mathcal{B}'_1}[\mathcal{E}]$.
Nous obtenons deux homomorphismes de $\mathcal{O}_{\mathcal{B}'_2}$-algèbres
$$
a : \mathcal{O}_{\mathcal{B}'_1}[\mathcal{E}] \to \mathcal{O}'_1
\quad\text{et}\quad
b : \mathcal{O}_{\mathcal{B}'_1}[\mathcal{E}] \to \mathcal{O}'_2
$$
qui induisent des applications
$a|_{\mathcal{I}'} : \mathcal{I}' \to \mathcal{G}_1$ et
$b|_{\mathcal{I}'} : \mathcal{I}' \to \mathcal{G}_2$. Les deux applications
$a$ et $b$ annulent $(\mathcal{I}')^2$. De plus, $a$ et $b$ coïncident sur
$\mathcal{J}_1\mathcal{O}_{\mathcal{B}'_1}[\mathcal{E}]$ comme applications
à valeurs dans $\mathcal{G}_2$, car le carré de gauche de
(\ref{equation-huge-2-ringed-topoi}) est commutatif. La différence
$b|_{\mathcal{I}'} - \nu \circ a|_{\mathcal{I}'}$ induit donc une application
$\mathcal{O}_1$-linéaire bien définie
$$
\xi : \mathcal{I}/\mathcal{I}^2 \longrightarrow \mathcal{G}_2
$$
qui envoie la classe d'une section locale $f$ de $\mathcal{I}$ sur
$a(f') - \nu(b(f'))$, où $f'$ est un relèvement de $f$ en une section locale
de $\mathcal{I}'$. Nous notons
$[\xi] \in \Ext^1_{\mathcal{O}_1}(\NL(\alpha), \mathcal{G}_2)$ son image
(voir ci-dessous).

\medskip\noindent
Étape 2. L'annulation de $[\xi]$ est nécessaire. Écrivons
$\Omega =
\Omega_{\mathcal{O}_{\mathcal{B}_1}[\mathcal{E}]/\mathcal{O}_{\mathcal{B}_1}}
\otimes_{\mathcal{O}_{\mathcal{B}_1}[\mathcal{E}]} \mathcal{O}_1$.
Observons que $\NL(\alpha) = (\mathcal{I}/\mathcal{I}^2 \to \Omega)$ s'insère
dans un triangle distingué
$$
\Omega[0] \to
\NL(\alpha) \to
\mathcal{I}/\mathcal{I}^2[1] \to
\Omega[1]
$$
Nous voyons donc que $[\xi]$ est nul si et seulement si $\xi$ est une composée
$\mathcal{I}/\mathcal{I}^2 \to \Omega \to \mathcal{G}_2$ pour une certaine
application $\Omega \to \mathcal{G}_2$.
Supposons qu'il existe un homomorphisme de faisceaux d'anneaux
$\varphi : \mathcal{O}'_1 \to \mathcal{O}'_2$ s'insérant dans
(\ref{equation-huge-2-ringed-topoi}). Considérons alors l'application
$\mathcal{O}'_1[\mathcal{E}] \to \mathcal{G}_2$,
$f' \mapsto b(f') - \varphi(a(f'))$. Un calcul montre qu'elle annule
$\mathcal{J}_1\mathcal{O}_{\mathcal{B}'_1}[\mathcal{E}]$ et induit une
dérivation $\mathcal{O}_{\mathcal{B}_1}[\mathcal{E}] \to \mathcal{G}_2$.
L'application linéaire $\Omega \to \mathcal{G}_2$ ainsi obtenue témoigne de
l'égalité $[\xi] = 0$ dans ce cas.

\medskip\noindent
Étape 3. L'annulation de $[\xi]$ est suffisante. Soit
$\theta : \Omega \to \mathcal{G}_2$ une application $\mathcal{O}_1$-linéaire
telle que $\xi$ soit égale à
$\theta \circ (\mathcal{I}/\mathcal{I}^2 \to \Omega)$.
Un calcul montre alors que
$$
b + \theta \circ d :
\mathcal{O}_{\mathcal{B}'_1}[\mathcal{E}]
\longrightarrow
\mathcal{O}'_2
$$
annule $\mathcal{I}'$ et définit donc une application
$\mathcal{O}'_1 \to \mathcal{O}'_2$ s'insérant dans
(\ref{equation-huge-2-ringed-topoi}).

\medskip\noindent
Démonstration de (2) dans le cas particulier ci-dessus. Omise. Indication :
c'est exactement la même démonstration que celle de la partie (2) du
Lemme \ref{lemma-huge-diagram}.
\end{proof}

\begin{lemma}
\label{lemma-NL-represent-ext-class-ringed-topoi}
Soit $\mathcal{C}$ un site. Soit $\mathcal{A} \to \mathcal{B}$ un
homomorphisme de faisceaux d'anneaux sur $\mathcal{C}$.
Soit $\mathcal{G}$ un $\mathcal{B}$-module.
Soit
$\xi \in \Ext^1_\mathcal{B}(\NL_{\mathcal{B}/\mathcal{A}}, \mathcal{G})$.
Il existe une application de faisceaux d'ensembles
$\alpha : \mathcal{E} \to \mathcal{B}$ telle que
$\xi \in \Ext^1_\mathcal{B}(\NL(\alpha), \mathcal{G})$ soit la classe d'une
application $\mathcal{I}/\mathcal{I}^2 \to \mathcal{G}$ (voir les notations
dans la démonstration).
\end{lemma}

\begin{proof}
Rappelons qu'étant donnée $\alpha : \mathcal{E} \to \mathcal{B}$ telle que
$\mathcal{A}[\mathcal{E}] \to \mathcal{B}$ soit surjective de noyau
$\mathcal{I}$, le complexe
$\NL(\alpha) = (\mathcal{I}/\mathcal{I}^2 \to 
\Omega_{\mathcal{A}[\mathcal{E}]/\mathcal{A}}
\otimes_{\mathcal{A}[\mathcal{E}]} \mathcal{B})$ est canoniquement isomorphe
à $\NL_{\mathcal{B}/\mathcal{A}}$ ; voir Modules sur les sites,
Lemme \ref{sites-modules-lemma-NL-up-to-qis}.
Observons en outre que
$\Omega = \Omega_{\mathcal{A}[\mathcal{E}]/\mathcal{A}}
\otimes_{\mathcal{A}[\mathcal{E}]} \mathcal{B}$ est le faisceau associé au
préfaisceau
$U \mapsto \bigoplus_{e \in \mathcal{E}(U)} \mathcal{B}(U)$.
Autrement dit, $\Omega$ est le $\mathcal{B}$-module libre sur le faisceau
d'ensembles $\mathcal{E}$ ; il existe en particulier une application
canonique $\mathcal{E} \to \Omega$.

\medskip\noindent
Ceci étant dit, choisissons un $\mathcal{E}$ (par exemple
$\mathcal{E} = \mathcal{B}$ comme dans la définition du complexe cotangent
naïf). L'obstruction à écrire $\xi$ comme la classe d'une application
$\mathcal{I}/\mathcal{I}^2 \to \mathcal{G}$ est un élément de
$\Ext^1_\mathcal{B}(\Omega, \mathcal{G})$. Supposons qu'il soit représenté par
l'extension $0 \to \mathcal{G} \to \mathcal{H} \to \Omega \to 0$ de
$\mathcal{B}$-modules. Considérons le faisceau d'ensembles
$\mathcal{E}' = \mathcal{E} \times_\Omega \mathcal{H}$, muni d'une application
induite $\alpha' : \mathcal{E}' \to \mathcal{B}$.
Posons $\mathcal{I}' = \Ker(\mathcal{A}[\mathcal{E}'] \to \mathcal{B})$ et
$\Omega' = \Omega_{\mathcal{A}[\mathcal{E}']/\mathcal{A}}
\otimes_{\mathcal{A}[\mathcal{E}']} \mathcal{B}$.
L'image inverse de $\xi$ par le quasi-isomorphisme
$\NL(\alpha') \to \NL(\alpha)$ s'envoie sur zéro dans
$\Ext^1_\mathcal{B}(\Omega', \mathcal{G})$, car l'image inverse de l'extension
$\mathcal{H}$ par l'application $\Omega' \to \Omega$ est scindée : en effet,
$\Omega'$ est le $\mathcal{B}$-module libre sur le faisceau d'ensembles
$\mathcal{E}'$ et, par construction, il existe un diagramme commutatif
$$
\xymatrix{
\mathcal{E}' \ar[r] \ar[d] & \mathcal{E} \ar[d] \\
\mathcal{H} \ar[r] & \Omega
}
$$
Cela achève la démonstration.
\end{proof}

\begin{lemma}
\label{lemma-choices-ringed-topoi}
S'il existe une solution de (\ref{equation-to-solve-ringed-topoi}), alors
l'ensemble des classes d'isomorphisme de solutions est un espace principal
homogène sous
$\Ext^1_\mathcal{O}(
\NL_{\mathcal{O}/\mathcal{O}_\mathcal{B}}, \mathcal{G})$.
\end{lemma}

\begin{proof}
Observons d'emblée qu'étant données deux solutions $\mathcal{O}'_1$ et
$\mathcal{O}'_2$ de (\ref{equation-to-solve-ringed-topoi}), le
Lemme \ref{lemma-huge-diagram-ringed-topoi} fournit un élément d'obstruction
$o(\mathcal{O}'_1, \mathcal{O}'_2) \in \Ext^1_\mathcal{O}(
\NL_{\mathcal{O}/\mathcal{O}_\mathcal{B}}, \mathcal{G})$ à l'existence d'une
application $\mathcal{O}'_1 \to \mathcal{O}'_2$.
Cet élément est manifestement l'obstruction à l'existence d'un isomorphisme ;
il sépare donc les classes d'isomorphisme. Pour achever la démonstration, il
suffit par conséquent de montrer qu'étant donnés une solution $\mathcal{O}'$ et
un élément
$\xi \in \Ext^1_\mathcal{O}(
\NL_{\mathcal{O}/\mathcal{O}_\mathcal{B}}, \mathcal{G})$, nous pouvons trouver
une seconde solution $\mathcal{O}'_\xi$ telle que
$o(\mathcal{O}', \mathcal{O}'_\xi) = \xi$.

\medskip\noindent
Choisissons $\alpha : \mathcal{E} \to \mathcal{O}$ comme dans le
Lemme \ref{lemma-NL-represent-ext-class-ringed-topoi} pour la classe $\xi$.
Considérons la surjection
$f^{-1}\mathcal{O}_\mathcal{B}[\mathcal{E}] \to \mathcal{O}$ de noyau
$\mathcal{I}$ et le complexe cotangent naïf correspondant
$\NL(\alpha) = (\mathcal{I}/\mathcal{I}^2 \to
\Omega_{f^{-1}\mathcal{O}_\mathcal{B}[\mathcal{E}]/
f^{-1}\mathcal{O}_\mathcal{B}}
\otimes_{f^{-1}\mathcal{O}_\mathcal{B}[\mathcal{E}]} \mathcal{O})$.
D'après le lemme, $\xi$ est la classe d'un morphisme
$\delta : \mathcal{I}/\mathcal{I}^2 \to \mathcal{G}$.
Après avoir remplacé $\mathcal{E}$ par
$\mathcal{E} \times_\mathcal{O} \mathcal{O}'$, nous pouvons aussi supposer que
$\alpha$ se factorise par une application
$\alpha' : \mathcal{E} \to \mathcal{O}'$.

\medskip\noindent
Ces choix déterminent un homomorphisme de
$f^{-1}\mathcal{O}_{\mathcal{B}'}$-algèbres
$\varphi : \mathcal{O}_{\mathcal{B}'}[\mathcal{E}] \to \mathcal{O}'$.
Posons $\mathcal{I}' = \Ker(\varphi)$.
Observons que $\varphi$ induit une application
$\varphi|_{\mathcal{I}'} : \mathcal{I}' \to \mathcal{G}$ et que
$\mathcal{O}'$ est la somme amalgamée, comme dans le diagramme suivant
$$
\xymatrix{
0 \ar[r] & \mathcal{G} \ar[r] & \mathcal{O}' \ar[r] &
\mathcal{O} \ar[r] & 0 \\
0 \ar[r] & \mathcal{I}' \ar[u]^{\varphi|_{\mathcal{I}'}} \ar[r] &
f^{-1}\mathcal{O}_{\mathcal{B}'}[\mathcal{E}] \ar[u] \ar[r] &
\mathcal{O} \ar[u]_{=} \ar[r] & 0
}
$$
Soit $\psi : \mathcal{I}' \to \mathcal{G}$ la somme de l'application
$\varphi|_{\mathcal{I}'}$ et de la composée
$$
\mathcal{I}' \to \mathcal{I}'/(\mathcal{I}')^2 \to
\mathcal{I}/\mathcal{I}^2 \xrightarrow{\delta} \mathcal{G}.
$$
La somme amalgamée suivant $\psi$ est alors une autre extension d'anneaux
$\mathcal{O}'_\xi$ s'insérant dans un diagramme comme ci-dessus.
Un calcul (omis) montre que $o(\mathcal{O}', \mathcal{O}'_\xi) = \xi$, comme
voulu.
\end{proof}

\begin{lemma}
\label{lemma-extensions-of-relative-ringed-topoi}
Soit $f : (\Sh(\mathcal{C}), \mathcal{O}) \to
(\Sh(\mathcal{B}), \mathcal{O}_\mathcal{B})$ un morphisme de topos annelés.
Soit $\mathcal{G}$ un $\mathcal{O}$-module.
L'ensemble des classes d'isomorphisme d'extensions de
$f^{-1}\mathcal{O}_\mathcal{B}$-algèbres
$$
0 \to \mathcal{G} \to \mathcal{O}' \to \mathcal{O} \to 0
$$
où $\mathcal{G}$ est un idéal de carré nul\footnote{Autrement dit, l'ensemble
des classes d'isomorphisme d'épaississements du premier ordre
$i : (\Sh(\mathcal{C}), \mathcal{O}) \to (\Sh(\mathcal{C}), \mathcal{O}')$
au-dessus de $(\Sh(\mathcal{B}), \mathcal{O}_\mathcal{B})$, munis d'un
isomorphisme $\mathcal{G} \to \Ker(i^\sharp)$ de $\mathcal{O}$-modules.}
est canoniquement en bijection avec
$\Ext^1_\mathcal{O}(\NL_{\mathcal{O}/\mathcal{O}_\mathcal{B}}, \mathcal{G})$.
\end{lemma}

\begin{proof}
Pour le démontrer, nous appliquons les résultats précédents au cas où
(\ref{equation-to-solve-ringed-topoi}) est donné par le diagramme
$$
\xymatrix{
0 \ar[r] &
\mathcal{G} \ar[r] &
{?} \ar[r] &
\mathcal{O} \ar[r] & 0 \\
0 \ar[r] &
0 \ar[u] \ar[r] &
f^{-1}\mathcal{O}_\mathcal{B} \ar[u] \ar[r]^{\text{id}} &
f^{-1}\mathcal{O}_\mathcal{B} \ar[u] \ar[r] & 0
}
$$
Notre lemme résulte donc du Lemme \ref{lemma-choices-ringed-topoi} et du fait
qu'il existe une solution, à savoir $\mathcal{G} \oplus \mathcal{O}$.
(Voir la remarque ci-dessous pour une construction directe de la bijection.)
\end{proof}

\begin{remark}
\label{remark-extensions-of-relative-ringed-topoi}
Soient $f : (\Sh(\mathcal{C}), \mathcal{O}) \to
(\mathcal{B}, \mathcal{O}_\mathcal{B})$ et $\mathcal{G}$ comme dans le
Lemme \ref{lemma-extensions-of-relative-ringed-topoi}.
Considérons une extension
$0 \to \mathcal{G} \to \mathcal{O}' \to \mathcal{O} \to 0$ comme dans le
lemme. Nous pouvons choisir un faisceau d'ensembles $\mathcal{E}$ et un
diagramme commutatif
$$
\xymatrix{
\mathcal{E} \ar[d]_{\alpha'} \ar[rd]^\alpha \\
\mathcal{O}' \ar[r] & \mathcal{O}
}
$$
tels que $f^{-1}\mathcal{O}_\mathcal{B}[\mathcal{E}] \to \mathcal{O}$ soit
surjective de noyau $\mathcal{J}$.
(On peut par exemple prendre n'importe quel faisceau d'ensembles se surjectant
sur $\mathcal{O}'$.) Alors
$$
\NL_{\mathcal{O}/\mathcal{O}_\mathcal{B}} \cong \NL(\alpha) = 
\left(
\mathcal{J}/\mathcal{J}^2
\longrightarrow
\Omega_{f^{-1}\mathcal{O}_\mathcal{B}[\mathcal{E}]/
f^{-1}\mathcal{O}_\mathcal{B}}
\otimes_{f^{-1}\mathcal{O}_\mathcal{B}[\mathcal{E}]} \mathcal{O}\right)
$$
Voir Modules sur les sites,
section \ref{sites-modules-section-netherlander}, et en particulier le
Lemme \ref{sites-modules-lemma-NL-up-to-qis}.
Bien entendu, $\alpha'$ détermine une application
$f^{-1}\mathcal{O}_\mathcal{B}[\mathcal{E}] \to \mathcal{O}'$, qui détermine
à son tour une application
$$
\mathcal{J}/\mathcal{J}^2 \longrightarrow \mathcal{G}
$$
qui détermine à son tour l'élément de
$\Ext^1_\mathcal{O}(\NL(\alpha), \mathcal{G}) =
\Ext^1_\mathcal{O}(\NL_{\mathcal{O}/\mathcal{O}_\mathcal{B}}, \mathcal{G})$
correspondant à $\mathcal{O}'$ par la bijection du lemme.
\end{remark}

\begin{lemma}
\label{lemma-extensions-of-relative-ringed-topoi-functorial}
Soient $f : (\Sh(\mathcal{C}), \mathcal{O}_\mathcal{C}) \to
(\Sh(\mathcal{B}), \mathcal{O}_\mathcal{B})$ et
$g : (\Sh(\mathcal{D}), \mathcal{O}_\mathcal{D}) \to
(\Sh(\mathcal{C}), \mathcal{O}_\mathcal{C})$ des morphismes de topos annelés.
Soit $\mathcal{F}$ un $\mathcal{O}_\mathcal{C}$-module.
Soit $\mathcal{G}$ un $\mathcal{O}_\mathcal{D}$-module. Soit
$c : g^*\mathcal{F} \to \mathcal{G}$ une application
$\mathcal{O}_\mathcal{D}$-linéaire. Considérons enfin
\begin{enumerate}
\item[(a)]
$0 \to \mathcal{F} \to \mathcal{O}_{\mathcal{C}'} \to
\mathcal{O}_\mathcal{C} \to 0$, une extension de
$f^{-1}\mathcal{O}_\mathcal{B}$-algèbres correspondant à
$\xi \in \Ext^1_{\mathcal{O}_\mathcal{C}}(
\NL_{\mathcal{O}_\mathcal{C}/\mathcal{O}_\mathcal{B}}, \mathcal{F})$, et
\item[(b)]
$0 \to \mathcal{G} \to \mathcal{O}_{\mathcal{D}'} \to
\mathcal{O}_\mathcal{D} \to 0$, une extension de
$g^{-1}f^{-1}\mathcal{O}_\mathcal{B}$-algèbres correspondant à
$\zeta \in \Ext^1_{\mathcal{O}_\mathcal{D}}(
\NL_{\mathcal{O}_\mathcal{D}/\mathcal{O}_\mathcal{B}}, \mathcal{G})$.
\end{enumerate}
Voir le Lemme \ref{lemma-extensions-of-relative-ringed-topoi}.
Il existe alors un morphisme
$$
g' :
(\Sh(\mathcal{D}), \mathcal{O}_{\mathcal{D}'})
\longrightarrow
(\Sh(\mathcal{C}), \mathcal{O}_{\mathcal{C}'})
$$
de topos annelés au-dessus de
$(\Sh(\mathcal{B}), \mathcal{O}_\mathcal{B})$, compatible avec $g$ et $c$, si
et seulement si $\xi$ et $\zeta$ s'envoient sur le même élément de
$\Ext^1_{\mathcal{O}_\mathcal{D}}(
Lg^*\NL_{\mathcal{O}_\mathcal{C}/\mathcal{O}_\mathcal{B}}, \mathcal{G})$.
\end{lemma}

\begin{proof}
L'énoncé a un sens, car nous disposons des applications
$$
\Ext^1_{\mathcal{O}_\mathcal{C}}(
\NL_{\mathcal{O}_\mathcal{C}/\mathcal{O}_\mathcal{B}}, \mathcal{F}) \to
\Ext^1_{\mathcal{O}_\mathcal{D}}(
Lg^*\NL_{\mathcal{O}_\mathcal{C}/\mathcal{O}_\mathcal{B}}, Lg^*\mathcal{F}) \to
\Ext^1_{\mathcal{O}_\mathcal{D}}
(Lg^*\NL_{\mathcal{O}_\mathcal{C}/\mathcal{O}_\mathcal{B}}, \mathcal{G})
$$
au moyen de l'application
$Lg^*\mathcal{F} \to g^*\mathcal{F} \xrightarrow{c} \mathcal{G}$, et
$$
\Ext^1_{\mathcal{O}_Y}(
\NL_{\mathcal{O}_\mathcal{D}/\mathcal{O}_\mathcal{B}}, \mathcal{G}) \to
\Ext^1_{\mathcal{O}_Y}(
Lg^*\NL_{\mathcal{O}_\mathcal{C}/\mathcal{O}_\mathcal{B}}, \mathcal{G})
$$
au moyen de l'application
$Lg^*\NL_{\mathcal{O}_\mathcal{C}/\mathcal{O}_\mathcal{B}} \to
\NL_{\mathcal{O}_\mathcal{D}/\mathcal{O}_\mathcal{B}}$.
L'énoncé du lemme se déduit du
Lemme \ref{lemma-huge-diagram-ringed-topoi}, appliqué au diagramme
$$
\xymatrix{
& 0 \ar[r] &
\mathcal{G} \ar[r] &
\mathcal{O}_{\mathcal{D}'} \ar[r] &
\mathcal{O}_\mathcal{D} \ar[r] & 0 \\
& 0 \ar[r]|\hole & 0 \ar[u] \ar[r] &
g^{-1}f^{-1}\mathcal{O}_\mathcal{B} \ar[u] \ar[r]|\hole &
g^{-1}f^{-1}\mathcal{O}_\mathcal{B} \ar[u] \ar[r] & 0 \\
0 \ar[r] &
\mathcal{F} \ar[ruu] \ar[r] &
\mathcal{O}_{\mathcal{C}'} \ar[r] &
\mathcal{O}_\mathcal{C} \ar[ruu] \ar[r] & 0 \\
0 \ar[r] & 0 \ar[ruu]|\hole \ar[u] \ar[r] &
f^{-1}\mathcal{O}_\mathcal{B} \ar[ruu]|\hole \ar[u] \ar[r] &
f^{-1}\mathcal{O}_\mathcal{B} \ar[ruu]|\hole \ar[u] \ar[r] & 0
}
$$
ainsi que d'une compatibilité entre les constructions des démonstrations des
Lemmes \ref{lemma-extensions-of-relative-ringed-topoi} et
\ref{lemma-huge-diagram-ringed-topoi}, dont nous omettons l'énoncé et la
démonstration. (Voir la remarque ci-dessous pour un argument direct.)
\end{proof}

\begin{remark}
\label{remark-extensions-of-relative-ringed-topoi-functorial}
Soient $f : (\Sh(\mathcal{C}), \mathcal{O}_\mathcal{C}) \to
(\Sh(\mathcal{B}), \mathcal{O}_\mathcal{B})$,
$g : (\Sh(\mathcal{D}), \mathcal{O}_\mathcal{D}) \to
(\Sh(\mathcal{C}), \mathcal{O}_\mathcal{C})$,
$\mathcal{F}$,
$\mathcal{G}$,
$c : g^*\mathcal{F} \to \mathcal{G}$,
$0 \to \mathcal{F} \to \mathcal{O}_{\mathcal{C}'} \to
\mathcal{O}_\mathcal{C} \to 0$,
$\xi \in \Ext^1_{\mathcal{O}_\mathcal{C}}(
\NL_{\mathcal{O}_\mathcal{C}/\mathcal{O}_\mathcal{B}}, \mathcal{F})$,
$0 \to \mathcal{G} \to \mathcal{O}_{\mathcal{D}'} \to
\mathcal{O}_\mathcal{D} \to 0$, et
$\zeta \in \Ext^1_{\mathcal{O}_\mathcal{D}}(
\NL_{\mathcal{O}_\mathcal{D}/\mathcal{O}_\mathcal{B}}, \mathcal{G})$
comme dans le
Lemme \ref{lemma-extensions-of-relative-ringed-topoi-functorial}.
En prenant la somme amalgamée suivant
$c : g^{-1}\mathcal{F} \to \mathcal{G}$, nous pouvons construire une extension
$$
\xymatrix{
0 \ar[r] &
\mathcal{G} \ar[r] &
\mathcal{O}'_1 \ar[r] &
g^{-1}\mathcal{O}_\mathcal{C} \ar[r] & 0 \\
0 \ar[r] &
g^{-1}\mathcal{F} \ar[u]^c \ar[r] &
g^{-1}\mathcal{O}_{\mathcal{C}'} \ar[u] \ar[r] &
g^{-1}\mathcal{O}_\mathcal{C} \ar@{=}[u] \ar[r] & 0
}
$$
En prenant l'image réciproque suivant
$g^\sharp : g^{-1}\mathcal{O}_\mathcal{C} \to \mathcal{O}_\mathcal{D}$, nous
pouvons construire une extension
$$
\xymatrix{
0 \ar[r] &
\mathcal{G} \ar[r] &
\mathcal{O}_{\mathcal{D}'} \ar[r] &
\mathcal{O}_\mathcal{D} \ar[r] & 0 \\
0 \ar[r] &
\mathcal{G} \ar@{=}[u] \ar[r] &
\mathcal{O}'_2 \ar[u] \ar[r] &
g^{-1}\mathcal{O}_\mathcal{C} \ar[u] \ar[r] & 0
}
$$
Une chasse au diagramme montre qu'il existe un morphisme
$g' : (\Sh(\mathcal{D}), \mathcal{O}_{\mathcal{D}'}) \to
(\Sh(\mathcal{C}), \mathcal{O}_{\mathcal{C}'})$ au-dessus de
$(\Sh(\mathcal{B}), \mathcal{O}_\mathcal{B})$, compatible avec $g$ et $c$, si
et seulement si $\mathcal{O}'_1$ est isomorphe à $\mathcal{O}'_2$ comme
extension de $g^{-1}f^{-1}\mathcal{O}_\mathcal{B}$-algèbres de
$g^{-1}\mathcal{O}_\mathcal{C}$ par $\mathcal{G}$.
D'après le Lemme \ref{lemma-extensions-of-relative-ringed-topoi}, ces extensions
sont classifiées par le membre de gauche de
$$
\Ext^1_{g^{-1}\mathcal{O}_\mathcal{C}}(
\NL_{g^{-1}\mathcal{O}_\mathcal{C}/g^{-1}f^{-1}\mathcal{O}_\mathcal{B}},
\mathcal{G}) =
\Ext^1_{\mathcal{O}_\mathcal{D}}(
Lg^*\NL_{\mathcal{O}_\mathcal{C}/\mathcal{O}_\mathcal{B}}, \mathcal{G})
$$
Ici, l'égalité résulte de l'adjonction entre produit tensoriel et Hom, et des
égalités
$$
\NL_{g^{-1}\mathcal{O}_\mathcal{C}/g^{-1}f^{-1}\mathcal{O}_\mathcal{B}} =
g^{-1}\NL_{\mathcal{O}_\mathcal{C}/\mathcal{O}_\mathcal{B}}
\quad\text{et}\quad
Lg^*\NL_{\mathcal{O}_\mathcal{C}/\mathcal{O}_\mathcal{B}} =
g^{-1}\NL_{\mathcal{O}_\mathcal{C}/\mathcal{O}_\mathcal{B}}
\otimes_{g^{-1}\mathcal{O}_X}^\mathbf{L} \mathcal{O}_Y
$$
Pour la première, voir Modules sur les sites,
Lemme \ref{sites-modules-lemma-pullback-NL} ; la seconde résulte de la
définition de l'image réciproque dérivée.
Ainsi, pour voir que le
Lemme \ref{lemma-extensions-of-relative-ringed-topoi-functorial} est vrai, il
suffit de montrer que $\mathcal{O}'_1$ correspond à l'image de $\xi$ et que
$\mathcal{O}'_2$ correspond à l'image de $\zeta$.
La correspondance entre $\xi$ et $\mathcal{O}'_1$ résulte immédiatement de la
construction de la classe $\xi$ dans la
Remarque \ref{remark-extensions-of-relative-ringed-topoi}.
Pour la correspondance entre $\zeta$ et $\mathcal{O}'_2$, choisissons d'abord
un diagramme commutatif
$$
\xymatrix{
\mathcal{E} \ar[d]_{\beta'} \ar[rd]^\beta \\
\mathcal{O}_{\mathcal{D}'} \ar[r] & \mathcal{O}_\mathcal{D}
}
$$
tel que $g^{-1}f^{-1}\mathcal{O}_\mathcal{B}[\mathcal{E}] \to
\mathcal{O}_\mathcal{D}$ soit surjective de noyau $\mathcal{K}$. Choisissons
ensuite un diagramme commutatif
$$
\xymatrix{
\mathcal{E} \ar[d]_{\beta'} &
\mathcal{E}' \ar[l]^\varphi \ar[d]_{\alpha'} \ar[rd]^\alpha \\
\mathcal{O}_{\mathcal{D}'} &
\mathcal{O}'_2 \ar[l] \ar[r] &
g^{-1}\mathcal{O}_\mathcal{C}
}
$$
tel que $g^{-1}f^{-1}\mathcal{O}_\mathcal{B}[\mathcal{E}'] \to
g^{-1}\mathcal{O}_\mathcal{C}$ soit surjective de noyau $\mathcal{J}$.
(On peut par exemple prendre simplement
$\mathcal{E}' = \mathcal{E} \amalg \mathcal{O}'_2$ comme faisceau
d'ensembles.)
L'application $\varphi$ induit une application de complexes
$\NL(\alpha) \to \NL(\beta)$ (notations de Modules,
section \ref{modules-section-netherlander}), et en particulier
$\bar\varphi : \mathcal{J}/\mathcal{J}^2 \to \mathcal{K}/\mathcal{K}^2$.
Alors $\NL(\alpha) \cong \NL_{\mathcal{O}_\mathcal{D}/\mathcal{O}_\mathcal{B}}$
et
$\NL(\beta) \cong
\NL_{g^{-1}\mathcal{O}_\mathcal{C}/g^{-1}f^{-1}\mathcal{O}_\mathcal{B}}$, et
l'application de complexes $\NL(\alpha) \to \NL(\beta)$ représente
l'application
$Lg^*\NL_{\mathcal{O}_\mathcal{C}/\mathcal{O}_\mathcal{B}} \to
\NL_{\mathcal{O}_\mathcal{D}/\mathcal{O}_\mathcal{B}}$ employée dans l'énoncé
du Lemme \ref{lemma-extensions-of-relative-ringed-topoi-functorial} (voir la
première partie de sa démonstration). Or $\zeta$ correspond à la classe de
l'application $\mathcal{K}/\mathcal{K}^2 \to \mathcal{G}$ induite par
$\beta'$ ; voir la Remarque \ref{remark-extensions-of-relative-ringed-topoi}.
De même, l'extension $\mathcal{O}'_2$ correspond à l'application
$\mathcal{J}/\mathcal{J}^2 \to \mathcal{G}$ induite par $\alpha'$.
Le diagramme commutatif ci-dessus montre que cette application est la composée
de l'application $\mathcal{K}/\mathcal{K}^2 \to \mathcal{G}$ induite par
$\beta'$ avec l'application
$\bar\varphi : \mathcal{J}/\mathcal{J}^2 \to \mathcal{K}/\mathcal{K}^2$.
C'est la compatibilité recherchée.
\end{remark}

\begin{lemma}
\label{lemma-parametrize-solutions-ringed-topoi}
Soient $t : (\Sh(\mathcal{B}), \mathcal{O}_\mathcal{B}) \to
(\Sh(\mathcal{B}'), \mathcal{O}_{\mathcal{B}'})$,
$\mathcal{J} = \Ker(t^\sharp)$,
$f : (\Sh(\mathcal{C}), \mathcal{O}) \to
(\Sh(\mathcal{B}), \mathcal{O}_\mathcal{B})$, $\mathcal{G}$ et
$c : \mathcal{J} \to \mathcal{G}$ comme dans
(\ref{equation-to-solve-ringed-topoi}).
Notons $\xi \in \Ext^1_{\mathcal{O}_\mathcal{B}}(
\NL_{\mathcal{O}_\mathcal{B}/\mathcal{O}_{\mathcal{B}'}}, \mathcal{J})$
l'élément correspondant à l'extension $\mathcal{O}_{\mathcal{B}'}$ de
$\mathcal{O}_\mathcal{B}$ par $\mathcal{J}$ au moyen du
Lemme \ref{lemma-extensions-of-relative-ringed-topoi}.
L'ensemble des classes d'isomorphisme de solutions est canoniquement en
bijection avec la fibre de
$$
\Ext^1_\mathcal{O}(\NL_{\mathcal{O}/\mathcal{O}_{\mathcal{B}'}},
\mathcal{G})\to
\Ext^1_\mathcal{O}(
Lf^*\NL_{\mathcal{O}_\mathcal{B}/\mathcal{O}_{\mathcal{B}'}}, \mathcal{G})
$$
au-dessus de l'image de $\xi$.
\end{lemma}

\begin{proof}
D'après le Lemme \ref{lemma-extensions-of-relative-ringed-topoi}, appliqué à
$t \circ f : (\Sh(\mathcal{C}), \mathcal{O}) \to
(\Sh(\mathcal{B}'), \mathcal{O}_{\mathcal{B}'})$ et au $\mathcal{O}$-module
$\mathcal{G}$, les éléments $\zeta$ de
$\Ext^1_\mathcal{O}(\NL_{\mathcal{O}/\mathcal{O}_{\mathcal{B}'}},
\mathcal{G})$ paramètrent les extensions
$0 \to \mathcal{G} \to \mathcal{O}' \to \mathcal{O} \to 0$ de
$f^{-1}\mathcal{O}_{\mathcal{B}'}$-algèbres. D'après le
Lemme \ref{lemma-extensions-of-relative-ringed-topoi-functorial}, appliqué à
$$
(\Sh(\mathcal{C}), \mathcal{O}) \xrightarrow{f}
(\Sh(\mathcal{B}), \mathcal{O}_\mathcal{B}) \xrightarrow{t}
(\Sh(\mathcal{B}'), \mathcal{O}_{\mathcal{B}'})
$$
et à $c : \mathcal{J} \to \mathcal{G}$, il existe un morphisme
$$
f' : 
(\Sh(\mathcal{C}), \mathcal{O}')
\longrightarrow
(\Sh(\mathcal{B}'), \mathcal{O}_{\mathcal{B}'})
$$
au-dessus de $(\Sh(\mathcal{B}'), \mathcal{O}_{\mathcal{B}'})$, compatible
avec $c$ et $f$, si et seulement si $\zeta$ s'envoie sur $\xi$.
Bien entendu, cela revient à dire que $\mathcal{O}'$ est une solution de
(\ref{equation-to-solve-ringed-topoi}).
\end{proof}

\section{Déformations d'espaces algébriques}
\label{section-deformations-spaces}

\noindent
Dans cette section, nous explicitons la signification des résultats de la
section \ref{section-deformations-ringed-topoi} pour les déformations d'espaces
algébriques.

\begin{lemma}
\label{lemma-match-thickenings}
Soit $S$ un schéma. Soit $i : Z \to Z'$ un morphisme d'espaces algébriques
au-dessus de $S$. Les assertions suivantes sont équivalentes :
\begin{enumerate}
\item $i$ est un épaississement d'espaces algébriques tel qu'il est défini dans
Compléments sur les morphismes d'espaces, section
\ref{spaces-more-morphisms-section-thickenings}, et
\item le morphisme associé
$i_{small} : (\Sh(Z_\etale), \mathcal{O}_Z) \to
(\Sh(Z'_\etale), \mathcal{O}_{Z'})$
de topos annelés (Propriétés des espaces, Lemme
\ref{spaces-properties-lemma-morphism-ringed-topoi}) est un épaississement au
sens de la section \ref{section-thickenings-ringed-topoi}.
\end{enumerate}
\end{lemma}

\begin{proof}
Soulignons qu'il ne s'agit pas d'une trivialité.

\medskip\noindent
Supposons (1). D'après Compléments sur les morphismes d'espaces,
Lemme \ref{spaces-more-morphisms-lemma-thickening-equivalence}, le morphisme
$i$ induit une équivalence des petits sites étales, et donc en particulier des
topos. Bien entendu, $i^\sharp$ est surjective à noyau localement nilpotent par
définition des épaississements.

\medskip\noindent
Supposons (2). (Cette implication est moins importante et relève davantage de
la curiosité.) Pour tout morphisme étale $Y' \to Z'$, nous voyons que
$Y = Z \times_{Z'} Y'$ a le même topos étale que $Y'$. En particulier, $Y'$ est
quasi-compact si et seulement si $Y$ est quasi-compact, car la quasi-compacité
est une notion toposique (Sites, Lemme \ref{sites-lemma-quasi-compact}).
Ceci étant dit, nous voyons que $Y'$ est quasi-compact et quasi-séparé si et
seulement si $Y$ est quasi-compact et quasi-séparé (car on peut caractériser la
quasi-séparation de $Y'$ en demandant que, pour tous espaces algébriques
quasi-compacts $Y'_1, Y'_2$ étales sur $Y'$, le produit
$Y'_1 \times_{Y'} Y'_2$ soit quasi-compact).
Prenons $Y'$ affine. Alors l'espace algébrique $Y$ est quasi-compact et
quasi-séparé. Pour tout $\mathcal{O}_Y$-module quasi-cohérent $\mathcal{F}$,
nous avons
$H^q(Y, \mathcal{F}) = H^q(Y', (Y \to Y')_*\mathcal{F})$, car les topos étales
sont les mêmes. Ensuite,
$H^q(Y', (Y \to Y')_*\mathcal{F}) = 0$, car l'image directe est
quasi-cohérente (Morphismes d'espaces,
Lemme \ref{spaces-morphisms-lemma-pushforward}) et $Y$ est affine.
Il s'ensuit que $Y'$ est affine d'après Cohomologie des espaces, Proposition
\ref{spaces-cohomology-proposition-vanishing-affine} (il existe sûrement une
démonstration de cette implication qui évite cette référence).
Ainsi, $i$ est un morphisme affine. Dans le cas affine, les conditions de la
section \ref{section-thickenings-ringed-topoi} entraînent aisément que $i$ est
un épaississement d'espaces algébriques.
\end{proof}

\begin{lemma}
\label{lemma-deform-spaces}
Soit $S$ un schéma.
Soit $Y \subset Y'$ un épaississement du premier ordre d'espaces algébriques
au-dessus de $S$.
Soit $f : X \to Y$ un morphisme plat d'espaces algébriques au-dessus de $S$.
S'il existe un morphisme plat $f' : X' \to Y'$ d'espaces algébriques au-dessus
de $S$ et un isomorphisme $a : X \to X' \times_{Y'} Y$ au-dessus de $Y$, alors
\begin{enumerate}
\item l'ensemble des classes d'isomorphisme de couples
$(f' : X' \to Y', a)$ est un espace principal homogène sous
$\Ext^1_{\mathcal{O}_X}(\NL_{X/Y}, f^*\mathcal{C}_{Y/Y'})$, et
\item l'ensemble des automorphismes $\varphi : X' \to X'$ au-dessus de $Y'$
qui se réduisent à l'identité sur $X' \times_{Y'} Y$ est
$\Ext^0_{\mathcal{O}_X}(\NL_{X/Y}, f^*\mathcal{C}_{Y/Y'})$.
\end{enumerate}
\end{lemma}

\begin{proof}
Nous appliquerons les résultats sur les déformations de topos annelés aux
petits topos étales des espaces algébriques du lemme.
Nous pouvons considérer $X$ comme un sous-espace fermé de $X'$ de sorte que
$(f, f') : (X \subset X') \to (Y \subset Y')$ soit un morphisme
d'épaississements du premier ordre.
D'après le Lemme \ref{lemma-match-thickenings}, cela se traduit par un
morphisme d'épaississements de topos annelés.
Il résulte alors de Compléments sur les morphismes d'espaces,
Lemme \ref{spaces-more-morphisms-lemma-deform} (ou du
Lemme \ref{lemma-deform-module-ringed-topoi}, plus général) que le faisceau
d'idéaux de $X$ dans $X'$ est égal à $f^*\mathcal{C}_{Y'/Y}$, et ceci équivaut
en fait à la platitude de $X'$ sur $Y'$.
Nous avons donc un diagramme commutatif
$$
\xymatrix{
0 \ar[r] & f^*\mathcal{C}_{Y/Y'} \ar[r] &
\mathcal{O}_{X'} \ar[r] &
\mathcal{O}_X \ar[r] & 0 \\
0 \ar[r] &
f_{small}^{-1}\mathcal{C}_{Y/Y'} \ar[u] \ar[r] &
f_{small}^{-1}\mathcal{O}_{Y'} \ar[u] \ar[r] &
f_{small}^{-1}\mathcal{O}_Y \ar[u] \ar[r] & 0
}
$$
Comparer avec (\ref{equation-to-solve-ringed-topoi}).
Observons que les automorphismes $\varphi$ de (2) donnent des automorphismes
$\varphi^\sharp : \mathcal{O}_{X'} \to \mathcal{O}_{X'}$ qui s'insèrent dans
le diagramme ci-dessus. Réciproquement, un automorphisme
$\alpha : \mathcal{O}_{X'} \to \mathcal{O}_{X'}$ s'insérant dans le diagramme
de faisceaux ci-dessus est égal à $\varphi^\sharp$ pour un certain
automorphisme $\varphi$ de (2), d'après Compléments sur les morphismes
d'espaces, Lemme
\ref{spaces-more-morphisms-lemma-first-order-thickening-maps}.
Enfin, d'après Compléments sur les morphismes d'espaces,
Lemme \ref{spaces-more-morphisms-lemma-first-order-thickening}, si nous trouvons
un autre faisceau d'anneaux $\mathcal{A}$ sur $X_\etale$ s'insérant dans le
diagramme
$$
\xymatrix{
0 \ar[r] & f^*\mathcal{C}_{Y/Y'} \ar[r] &
\mathcal{A} \ar[r] &
\mathcal{O}_X \ar[r] & 0 \\
0 \ar[r] &
f_{small}^{-1}\mathcal{C}_{Y/Y'} \ar[u] \ar[r] &
f_{small}^{-1}\mathcal{O}_{Y'} \ar[u] \ar[r] &
f_{small}^{-1}\mathcal{O}_Y \ar[u] \ar[r] & 0
}
$$
alors il existe un épaississement du premier ordre $X \subset X''$ tel que
$\mathcal{O}_{X''} = \mathcal{A}$ et, en appliquant encore Compléments sur les
morphismes d'espaces, Lemme
\ref{spaces-more-morphisms-lemma-first-order-thickening-maps}, nous obtenons un
morphisme $(f, f'') : (X \subset X'') \to (Y \subset Y')$ ayant toutes les
propriétés voulues.
La partie (1) résulte donc du Lemme \ref{lemma-choices-ringed-topoi}, et la
partie (2), de la partie (2) du
Lemme \ref{lemma-huge-diagram-ringed-topoi}.
(Notons que $\NL_{X/Y}$, tel qu'il est défini pour un morphisme d'espaces
algébriques dans Compléments sur les morphismes d'espaces, section
\ref{spaces-more-morphisms-section-netherlander}, coïncide avec $\NL_{X/Y}$ tel
qu'il est employé dans la section \ref{section-deformations-ringed-topoi}.)
\end{proof}

\noindent
Soit $S$ un schéma. Soit $f : X \to B$ un morphisme d'espaces algébriques
au-dessus de $S$.
Soit $\mathcal{F} \to \mathcal{G}$ un homomorphisme de
$\mathcal{O}_X$-modules (non nécessairement quasi cohérents).
Considérons le foncteur
$$
F :
\left\{
\begin{matrix}
\text{extensions of }f^{-1}\mathcal{O}_B\text{ algebras}\\
0 \to \mathcal{F} \to \mathcal{O}' \to \mathcal{O}_X \to 0\\
\text{où }\mathcal{F}\text{ est un idéal de carré nul}
\end{matrix}
\right\}
\longrightarrow
\left\{
\begin{matrix}
\text{extensions of }f^{-1}\mathcal{O}_B\text{ algebras}\\
0 \to \mathcal{G} \to \mathcal{O}' \to \mathcal{O}_X \to 0\\
\text{où }\mathcal{G}\text{ est un idéal de carré nul}
\end{matrix}
\right\}
$$
défini par somme amalgamée.

\begin{lemma}
\label{lemma-thickening-space-quasi-coherent}
Dans la situation ci-dessus, supposons que $X$ soit quasi-compact et
quasi-séparé, et que $DQ_X(\mathcal{F}) \to DQ_X(\mathcal{G})$
(Catégories dérivées des espaces, section
\ref{spaces-perfect-section-better-coherator}) soit un isomorphisme.
Alors le foncteur $F$ est une équivalence de catégories.
\end{lemma}

\begin{proof}
Rappelons que $\NL_{X/B}$ est un objet de $D_\QCoh(\mathcal{O}_X)$ ; voir
Compléments sur les morphismes d'espaces, Lemme
\ref{spaces-more-morphisms-lemma-netherlander-quasi-coherent}.
Notre hypothèse implique donc que les applications
$$
\Ext^i_X(\NL_{X/B}, \mathcal{F}) \longrightarrow
\Ext^i_X(\NL_{X/B}, \mathcal{G})
$$
sont des isomorphismes pour tout $i$. Le
Lemme \ref{lemma-huge-diagram-ringed-topoi} implique que notre foncteur est
pleinement fidèle. D'autre part, le foncteur est essentiellement surjectif par
le Lemme \ref{lemma-choices-ringed-topoi}, car nous avons les solutions
$\mathcal{O}_X \oplus \mathcal{F}$ et $\mathcal{O}_X \oplus \mathcal{G}$ dans
les deux catégories.
\end{proof}

\noindent
Soit $S$ un schéma. Soit $B \subset B'$ un épaississement du premier ordre
d'espaces algébriques au-dessus de $S$, de faisceau d'idéaux $\mathcal{J}$,
que nous considérons soit comme un $\mathcal{O}_B$-module quasi-cohérent, soit
comme un faisceau quasi-cohérent d'idéaux sur $B'$ ; voir Compléments sur les
morphismes d'espaces, section
\ref{spaces-more-morphisms-section-thickenings}.
Soit $f : X \to B$ un morphisme d'espaces algébriques au-dessus de $S$.
Soit $\mathcal{F} \to \mathcal{G}$ un homomorphisme de
$\mathcal{O}_X$-modules (non nécessairement quasi cohérents).
Soit $c : f^{-1}\mathcal{J} \to \mathcal{F}$ une application de
$f^{-1}\mathcal{O}_B$-modules, et notons
$c' : f^{-1}\mathcal{J} \to \mathcal{G}$ la composée.
Considérons le foncteur
$$
FT :
\{\text{solutions de }(\ref{equation-to-solve-ringed-topoi})
\text{ pour }\mathcal{F}\text{ et }c\}
\longrightarrow
\{\text{solutions de }(\ref{equation-to-solve-ringed-topoi})
\text{ pour }\mathcal{G}\text{ et }c'\}
$$
défini par somme amalgamée.

\begin{lemma}
\label{lemma-thickening-over-thickening-space-quasi-coherent}
Dans la situation ci-dessus, supposons que $X$ soit quasi-compact et
quasi-séparé, et que $DQ_X(\mathcal{F}) \to DQ_X(\mathcal{G})$
(Catégories dérivées des espaces, section
\ref{spaces-perfect-section-better-coherator}) soit un isomorphisme.
Alors le foncteur $FT$ est une équivalence de catégories.
\end{lemma}

\begin{proof}
Une solution de (\ref{equation-to-solve-ringed-topoi}) pour $\mathcal{F}$ donne
en particulier une extension de $f^{-1}\mathcal{O}_{B'}$-algèbres
$$
0 \to \mathcal{F} \to \mathcal{O}' \to \mathcal{O}_X \to 0
$$
où $\mathcal{F}$ est un idéal de carré nul. De même pour $\mathcal{G}$.
De plus, une telle extension fournit une application
$c_{\mathcal{O}'} : f^{-1}\mathcal{J} \to \mathcal{F}$.
Nous considérons donc la sous-catégorie pleine de ces extensions de
$f^{-1}\mathcal{O}_{B'}$-algèbres telles que $c = c_{\mathcal{O}'}$.
Manifestement, si $\mathcal{O}'' = F(\mathcal{O}')$, où $F$ est l'équivalence
du Lemme \ref{lemma-thickening-space-quasi-coherent} (appliqué cette fois à
$X \to B'$), alors $c_{\mathcal{O}''}$ est la composée de
$c_{\mathcal{O}'}$ et de l'application $\mathcal{F} \to \mathcal{G}$.
Cela démontre le lemme.
\end{proof}

\section{Déformations de complexes}
\label{section-deformations-complexes}

\noindent
Cette section sert de préparation à la suivante.
Nous utiliserons autant que possible les résultats des chapitres consacrés à
l'algèbre commutative.

\begin{lemma}
\label{lemma-canonical-class-algebra}
Soit $R' \to R$ une surjection d'anneaux dont le noyau est un idéal $I$ de
carré nul. Pour tout $K \in D^-(R)$, il existe une application canonique
$$
\omega(K) : K \longrightarrow K \otimes_R^\mathbf{L} I[2]
$$
dans $D(R)$ possédant les propriétés suivantes :
\begin{enumerate}
\item $\omega(K) = 0$ si et seulement s'il existe
$K' \in D(R')$ tel que $K' \otimes_{R'}^\mathbf{L} R = K$,
\item étant donnée une application $K \to L$ dans $D^-(R)$, le diagramme
$$
\xymatrix{
K \ar[d] \ar[rr]_-{\omega(K)} & &
K \otimes^\mathbf{L}_R I[2] \ar[d] \\
L \ar[rr]^-{\omega(L)} & &
L \otimes^\mathbf{L}_R I[2]
}
$$
est commutatif, et
\item la formation de $\omega(K)$ est compatible aux homomorphismes d'anneaux
$R' \to S'$ (voir la démonstration pour un énoncé précis).
\end{enumerate}
\end{lemma}

\begin{proof}
Choisissons un complexe $K^\bullet$ de $R$-modules libres, borné
supérieurement, représentant $K$. Nous pouvons alors choisir des
$R'$-modules libres $(K')^n$ relevant $K^n$.
Nous pouvons choisir des applications de $R'$-modules
$(d')^n_K : (K')^n \to (K')^{n + 1}$ relevant les différentielles
$d^n_K : K^n \to K^{n + 1}$ de $K^\bullet$.
Bien que les composées
$$
(d')^{n + 1}_K \circ (d')^n_K : (K')^n \to (K')^{n + 2}
$$
ne soient pas nécessairement nulles, elles se factorisent sous la forme
$$
(K')^n \to K^n \xrightarrow{\omega^n_K}
K^{n + 2} \otimes_R I = I(K')^{n + 2} \to (K')^{n + 2}
$$
car $d^{n + 1} \circ d^n = 0$.
Un calcul montre que $\omega^n_K$ définit une application de complexes.
Cette application de complexes définit $\omega(K)$.

\medskip\noindent
Montrons que cette construction est compatible avec une application de
complexes $\alpha^\bullet : K^\bullet \to L^\bullet$ de $R$-modules libres
bornés supérieurement et avec des choix de relèvements
$(K')^n, (L')^n, (d')^n_K, (d')^n_L$.
Choisissons pour cela $(\alpha')^n : (K')^n \to (L')^n$ relevant les
composantes $\alpha^n : K^n \to L^n$. Comme ci-dessus, nous obtenons une
factorisation
$$
(K')^n \to K^n \xrightarrow{h^n}
L^{n + 1} \otimes_R I = I(L')^{n + 1} \to (L')^{n + 2}
$$
de $(d')^n_L \circ (\alpha')^n - (\alpha')^{n + 1} \circ (d')_K^n$.
Un calcul agréable montre alors que
$$
\omega^n_L \circ \alpha^n =
(d_L^{n + 1} \otimes \text{id}_I) \circ h^n + h^{n + 1} \circ d_K^n +
(\alpha^{n + 2} \otimes \text{id}_I) \circ \omega^n_K
$$
Cela démontre la commutativité du diagramme de la partie (2) du lemme dans ce
cas particulier. En appliquant ceci à deux choix différents de complexes
libres bornés supérieurement représentant $K$, nous voyons que $\omega(K)$ est
bien définie. Bien entendu, (2) vaut alors aussi en général.

\medskip\noindent
Si $K$ se relève en $K'$ dans $D^-(R')$, nous pouvons représenter $K'$ par un
complexe de $R'$-modules libres borné supérieurement et voyons immédiatement
que $\omega(K) = 0$.
Réciproquement, revenons à nos choix $K^\bullet$,
$(K')^n$, $(d')^n_K$. Si $\omega(K) = 0$, nous pouvons trouver
$g^n : K^n \to K^{n + 1} \otimes_R I$ tel que
$$
\omega^n = (d_K^{n + 1} \otimes \text{id}_I) \circ g^n +
g^{n + 1} \circ d_K^n
$$
En prenant comme différentielles
$(d')^n_K - g^n : (K')^n \to (K')^{n + 1}$, nous obtenons donc un complexe de
$R'$-modules libres relevant $K^\bullet$. Cela démontre (1).

\medskip\noindent
Enfin, la partie (3) signifie ce qui suit. Soit $R' \to S'$ un homomorphisme
d'anneaux. Posons $S = S' \otimes_{R'} R$ et notons
$J = IS' \subset S'$ le noyau de carré nul de $S' \to S$. Étant donné
$K \in D^-(R)$, l'assertion est que nous obtenons un diagramme commutatif
$$
\xymatrix{
K \otimes_R^\mathbf{L} S \ar[d] \ar[rr]_-{\omega(K) \otimes \text{id}} & &
(K \otimes^\mathbf{L}_R I[2]) \otimes_R^\mathbf{L} S \ar[d] \\
K \otimes_R^\mathbf{L} S \ar[rr]^-{\omega(K \otimes_R^\mathbf{L} S)} & &
(K \otimes_R^\mathbf{L} S) \otimes^\mathbf{L}_S J[2]
}
$$
Ici, la flèche verticale de droite provient de
$$
(K \otimes^\mathbf{L}_R I[2]) \otimes_R^\mathbf{L} S =
(K \otimes_R^\mathbf{L} S) \otimes_S^\mathbf{L}
(I \otimes_R^\mathbf{L} S)[2] \longrightarrow
(K \otimes_R^\mathbf{L} S) \otimes_S^\mathbf{L} J[2]
$$
Choisissons $K^\bullet$, $(K')^n$ et $(d')^n_K$ comme ci-dessus.
Nous pouvons alors employer $K^\bullet \otimes_R S$,
$(K')^n \otimes_{R'} S'$ et $(d')^n_K \otimes \text{id}_{S'}$ pour la
construction de $\omega(K \otimes_R^\mathbf{L} S)$.
Avec ces choix, la commutativité se vérifie immédiatement au niveau des
applications de complexes.
\end{proof}

\section{Déformations de complexes sur des topos annelés}
\label{section-thickenings-complexes}

\noindent
Ce matériau est tiré de \cite{lieblich-complexes}.

\medskip\noindent
Les résultats de cette section valent dans le cadre d'un épaississement du
premier ordre de topos annelés tel qu'il est défini dans la
section \ref{section-thickenings-ringed-topoi}.
Cependant, afin de simplifier les notations, nous supposerons que les sites
sous-jacents $\mathcal{C}$ et $\mathcal{D}$ sont les mêmes.
De plus, l'homomorphisme surjectif $\mathcal{O}' \to \mathcal{O}$ de faisceaux
d'anneaux sera noté $\mathcal{O} \to \mathcal{O}_0$, conformément à l'usage
sans doute plus courant dans la littérature.

\begin{lemma}
\label{lemma-lift-complex}
Soit $\mathcal{C}$ un site. Soit $\mathcal{O} \to \mathcal{O}_0$ une
surjection de faisceaux d'anneaux. Supposons données les informations
suivantes :
\begin{enumerate}
\item des $\mathcal{O}$-modules plats $\mathcal{G}^n$,
\item des applications de $\mathcal{O}$-modules
$\mathcal{G}^n \to \mathcal{G}^{n + 1}$,
\item un complexe $\mathcal{K}_0^\bullet$ de $\mathcal{O}_0$-modules,
\item des applications de $\mathcal{O}$-modules
$\mathcal{G}^n \to \mathcal{K}_0^n$
\end{enumerate}
telles que
\begin{enumerate}
\item[(a)] $H^n(\mathcal{K}_0^\bullet) = 0$ pour $n \gg 0$,
\item[(b)] $\mathcal{G}^n = 0$ pour $n \gg 0$,
\item[(c)] en posant
$\mathcal{G}^n_0 = \mathcal{G}^n \otimes_\mathcal{O} \mathcal{O}_0$, les
applications induites déterminent un complexe $\mathcal{G}_0^\bullet$ et une
application de complexes
$\mathcal{G}_0^\bullet \to \mathcal{K}_0^\bullet$.
\end{enumerate}
Il existe alors
\begin{enumerate}
\item[(\romannumeral1)]
des $\mathcal{O}$-modules plats $\mathcal{F}^n$,
\item[(\romannumeral2)]
des applications de $\mathcal{O}$-modules
$\mathcal{F}^n \to \mathcal{F}^{n + 1}$,
\item[(\romannumeral3)]
des applications de $\mathcal{O}$-modules
$\mathcal{F}^n \to \mathcal{K}_0^n$,
\item[(\romannumeral4)]
des applications de $\mathcal{O}$-modules
$\mathcal{G}^n \to \mathcal{F}^n$,
\end{enumerate}
tels que $\mathcal{F}^n = 0$ pour $n \gg 0$, que les diagrammes
$$
\xymatrix{
\mathcal{G}^n \ar[r] \ar[d] & \mathcal{G}^{n + 1} \ar[d] \\
\mathcal{F}^n \ar[r] & \mathcal{F}^{n + 1}
}
$$
soient commutatifs pour tout $n$, que la composée
$\mathcal{G}^n \to \mathcal{F}^n \to \mathcal{K}_0^n$ soit l'application
donnée $\mathcal{G}^n \to \mathcal{K}_0^n$, et qu'en posant
$\mathcal{F}^n_0 = \mathcal{F}^n \otimes_\mathcal{O} \mathcal{O}_0$, nous
obtenions un complexe $\mathcal{F}_0^\bullet$ et une application de complexes
$\mathcal{F}_0^\bullet \to \mathcal{K}_0^\bullet$ qui soit un
quasi-isomorphisme.
\end{lemma}

\begin{proof}
Nous allons démontrer par récurrence descendante sur $e$ que nous pouvons
trouver $\mathcal{F}^n$,
$\mathcal{G}^n \to \mathcal{F}^n$ et
$\mathcal{F}^n \to \mathcal{F}^{n + 1}$ pour $n \geq e$, s'insérant dans un
diagramme commutatif
$$
\xymatrix{
\ldots \ar[r] &
\mathcal{G}^{e - 1} \ar[r] \ar@/_2pc/[dd] &
\mathcal{G}^e \ar[d] \ar[r] \ar@/_2pc/[dd] &
\mathcal{G}^{e + 1} \ar[d] \ar[r] \ar@/_2pc/[dd]|\hole &
\ldots \\
& &
\mathcal{F}^e \ar[d] \ar[r] &
\mathcal{F}^{e + 1} \ar[d] \ar[r] & \ldots \\
\ldots  \ar[r] &
\mathcal{K}_0^{e - 1} \ar[r] &
\mathcal{K}_0^e \ar[r] &
\mathcal{K}_0^{e + 1} \ar[r] & \ldots
}
$$
tel que $\mathcal{F}_0^\bullet$ soit un complexe et que l'application induite
$\mathcal{F}_0^\bullet \to \mathcal{K}_0^\bullet$ induise un isomorphisme sur
$H^n$ pour $n > e$ et une surjection pour $n = e$. Pour $e \gg 0$, ceci est
vrai, car les hypothèses (a) et (b) permettent de prendre
$\mathcal{F}^n = 0$ pour $n \geq e$.

\medskip\noindent
Étape de récurrence. Nous devons construire $\mathcal{F}^{e - 1}$ et les
applications $\mathcal{G}^{e - 1} \to \mathcal{F}^{e - 1}$,
$\mathcal{F}^{e - 1} \to \mathcal{F}^e$ et
$\mathcal{F}^{e - 1} \to \mathcal{K}_0^{e - 1}$.
Nous choisirons $\mathcal{F}^{e - 1} = A \oplus B \oplus C$ comme somme
directe de trois morceaux.

\medskip\noindent
Pour le premier, nous prenons $A = \mathcal{G}^{e - 1}$ et choisissons pour
$\mathcal{G}^{e - 1} \to \mathcal{F}^{e - 1}$ l'inclusion du premier facteur.
Les applications $A \to \mathcal{K}^{e - 1}_0$ et
$A \to \mathcal{F}^e$ seront les applications évidentes.

\medskip\noindent
Pour choisir $B$, considérons la surjection (par hypothèse de récurrence)
$$
\gamma :
\Ker(\mathcal{F}^e_0 \to \mathcal{F}^{e + 1}_0)
\longrightarrow
\Ker(\mathcal{K}^e_0 \to \mathcal{K}^{e + 1}_0)/
\Im(\mathcal{K}^{e - 1}_0 \to \mathcal{K}^e_0)
$$
Nous pouvons choisir un ensemble $I$, pour chaque $i \in I$ un objet $U_i$ de
$\mathcal{C}$, et des sections
$s_i \in \mathcal{F}^e(U_i)$,
$t_i \in \mathcal{K}^{e - 1}_0(U_i)$ telles que
\begin{enumerate}
\item $s_i$ s'envoie sur une section de $\Ker(\gamma) \subset
\Ker(\mathcal{F}^e_0 \to \mathcal{F}^{e + 1}_0)$,
\item $s_i$ et $t_i$ s'envoient sur la même section de
$\mathcal{K}^e_0$,
\item les sections $s_i$ engendrent $\Ker(\gamma)$ comme
$\mathcal{O}_0$-module.
\end{enumerate}
Nous omettons la justification complète ; on utilise le fait que
$\mathcal{F}^e \to \mathcal{F}^e_0$ est une application surjective de
faisceaux d'ensembles. Nous posons alors
$$
B = \bigoplus\nolimits_{i \in I} j_{U_i!}\mathcal{O}_{U_i}
$$
et définissons les applications $B \to \mathcal{F}^e$ et
$B \to \mathcal{K}_0^{e - 1}$ en utilisant $s_i$ et $t_i$ pour déterminer où
envoyer le facteur $j_{U_i!}\mathcal{O}_{U_i}$.

\medskip\noindent
Avec $\mathcal{F}^{e - 1} = A \oplus B$ et les applications ci-dessus, nous
obtenons un diagramme du type précédent pour $e - 1$, tel que
$\mathcal{F}_0^\bullet \to \mathcal{K}_0^\bullet$ induise un isomorphisme sur
$H^n$ pour $n \geq e$.
Pour que l'application soit surjective sur $H^{e - 1}$, nous choisissons le
facteur $C$ comme suit.
Choisissons un ensemble $J$, pour chaque $j \in J$ un objet $U_j$ de
$\mathcal{C}$ et une section $t_j$ de
$\Ker(\mathcal{K}^{e - 1}_0 \to \mathcal{K}^e_0)$ au-dessus de $U_j$, de sorte
que ces sections engendrent ce noyau sur $\mathcal{O}_0$. Posons alors
$$
C = \bigoplus\nolimits_{j \in J} j_{U_j!}\mathcal{O}_{U_j}
$$
et prenons l'application nulle $C \to \mathcal{F}^e$ ainsi que l'application
$C \to \mathcal{K}_0^{e - 1}$ définie en utilisant $s_j$ pour déterminer où
envoyer le facteur $j_{U_j!}\mathcal{O}_{U_j}$. Ceci achève l'étape de
récurrence en prenant $\mathcal{F}^{e - 1} = A \oplus B \oplus C$ et les
applications indiquées.
\end{proof}

\begin{lemma}
\label{lemma-canonical-class}
Soit $\mathcal{C}$ un site. Soit $\mathcal{O} \to \mathcal{O}_0$ une
surjection de faisceaux d'anneaux dont le noyau est un faisceau d'idéaux
$\mathcal{I}$ de carré nul. Pour tout objet $K_0$ de
$D^-(\mathcal{O}_0)$, il existe une application canonique
$$
\omega(K_0) :
K_0 \longrightarrow
K_0 \otimes_{\mathcal{O}_0}^\mathbf{L} \mathcal{I}[2]
$$
dans $D(\mathcal{O}_0)$ telle que, pour toute application
$K_0 \to L_0$ dans $D^-(\mathcal{O}_0)$, le diagramme
$$
\xymatrix{
K_0 \ar[d] \ar[rr]_-{\omega(K_0)} & &
(K_0 \otimes^\mathbf{L}_{\mathcal{O}_0} \mathcal{I})[2] \ar[d] \\
L_0 \ar[rr]^-{\omega(L_0)} & &
(L_0 \otimes^\mathbf{L}_{\mathcal{O}_0} \mathcal{I})[2]
}
$$
soit commutatif.
\end{lemma}

\begin{proof}
Représentons $K_0$ par un complexe quelconque
$\mathcal{K}_0^\bullet$ de $\mathcal{O}_0$-modules.
Appliquons le Lemme \ref{lemma-lift-complex} avec $\mathcal{G}^n = 0$ pour tout
$n$. Notons $d : \mathcal{F}^n \to \mathcal{F}^{n + 1}$ les applications
produites par le lemme. Nous voyons alors que
$d \circ d : \mathcal{F}^n \to \mathcal{F}^{n + 2}$ est nul modulo
$\mathcal{I}$. Comme $\mathcal{F}^n$ est plat, nous voyons que
$\mathcal{I}\mathcal{F}^n =
\mathcal{F}^n \otimes_{\mathcal{O}} \mathcal{I} =
\mathcal{F}^n_0 \otimes_{\mathcal{O}_0} \mathcal{I}$.
Nous obtenons donc une application canonique de complexes
$$
d \circ d : \mathcal{F}_0^\bullet
\longrightarrow
(\mathcal{F}_0^\bullet \otimes_{\mathcal{O}_0} \mathcal{I})[2]
$$
Puisque $\mathcal{F}_0^\bullet$ est un complexe borné supérieurement de
$\mathcal{O}_0$-modules plats, il est K-plat et peut servir à calculer le
produit tensoriel dérivé. De plus, l'application de complexes
$\mathcal{F}_0^\bullet \to \mathcal{K}_0^\bullet$ est un quasi-isomorphisme par
construction. La source et le but de l'application que nous venons de
construire représentent donc $K_0$ et
$K_0 \otimes_{\mathcal{O}_0}^\mathbf{L} \mathcal{I}[2]$ ; nous obtenons ainsi
notre application $\omega(K_0)$.

\medskip\noindent
Montrons que ce procédé est compatible avec les applications de complexes.
Soit $\mathcal{L}_0^\bullet$ un complexe représentant un autre objet de
$D^-(\mathcal{O}_0)$, et supposons que
$$
\mathcal{K}_0^\bullet \longrightarrow \mathcal{L}_0^\bullet
$$
soit une application de complexes. Appliquons le
Lemme \ref{lemma-lift-complex} au complexe $\mathcal{L}_0^\bullet$, aux modules
plats $\mathcal{F}^n$, aux applications
$\mathcal{F}^n \to \mathcal{F}^{n + 1}$ et aux composées
$\mathcal{F}^n \to \mathcal{K}_0^n \to \mathcal{L}_0^n$ (nous nous excusons
de l'inversion des lettres employées).
Nous obtenons des modules plats $\mathcal{G}^n$, des applications
$\mathcal{F}^n \to \mathcal{G}^n$, des applications
$\mathcal{G}^n \to \mathcal{G}^{n + 1}$ et des applications
$\mathcal{G}^n \to \mathcal{L}_0^n$ possédant toutes les propriétés du lemme.
Il est alors clair que
$$
\xymatrix{
\mathcal{F}_0^\bullet \ar[d] \ar[r] &
(\mathcal{F}_0^\bullet \otimes_{\mathcal{O}_0} \mathcal{I})[2] \ar[d] \\
\mathcal{G}_0^\bullet \ar[r] &
(\mathcal{G}_0^\bullet \otimes_{\mathcal{O}_0} \mathcal{I})[2]
}
$$
est un diagramme commutatif de complexes.

\medskip\noindent
Pour voir que $\omega(K_0)$ est bien définie, supposons que nous disposions de
deux complexes $\mathcal{K}_0^\bullet$ et
$(\mathcal{K}'_0)^\bullet$ de $\mathcal{O}_0$-modules représentant $K_0$, et de
deux systèmes
$(\mathcal{F}^n, d : \mathcal{F}^n \to \mathcal{F}^{n + 1},
\mathcal{F}^n \to \mathcal{K}_0^n)$
et
$((\mathcal{F}')^n, d : (\mathcal{F}')^n \to (\mathcal{F}')^{n + 1},
(\mathcal{F}')^n \to \mathcal{K}_0^n)$
comme ci-dessus. Nous pouvons choisir un complexe
$(\mathcal{K}''_0)^\bullet$ et des quasi-isomorphismes
$\mathcal{K}_0^\bullet  \to (\mathcal{K}''_0)^\bullet$
et
$(\mathcal{K}'_0)^\bullet  \to (\mathcal{K}''_0)^\bullet$
réalisant le fait que les deux complexes représentent $K_0$ dans la catégorie
dérivée. Appliquons ensuite le résultat du paragraphe précédent à
$$
(\mathcal{K}_0)^\bullet \oplus (\mathcal{K}'_0)^\bullet
\longrightarrow
(\mathcal{K}''_0)^\bullet
$$
Nous obtenons ainsi un diagramme commutatif
$$
\xymatrix{
\mathcal{F}_0^\bullet \oplus (\mathcal{F}'_0)^\bullet
\ar[d] \ar[r] &
(\mathcal{F}_0^\bullet \otimes_{\mathcal{O}_0} \mathcal{I})[2] \oplus
((\mathcal{F}'_0)^\bullet \otimes_{\mathcal{O}_0} \mathcal{I})[2] \ar[d] \\
\mathcal{G}_0^\bullet \ar[r] &
(\mathcal{G}_0^\bullet \otimes_{\mathcal{O}_0} \mathcal{I})[2]
}
$$
Puisque les flèches verticales donnent des quasi-isomorphismes sur les facteurs,
nous en déduisons la commutativité voulue dans $D(\mathcal{O}_0)$.

\medskip\noindent
La bonne définition étant établie, l'assertion de compatibilité aux
applications découle du résultat du deuxième paragraphe.
\end{proof}

\begin{lemma}
\label{lemma-induced-map}
Soit $(\mathcal{C}, \mathcal{O})$ un site annelé.
Soit $\alpha : K \to L$ une application dans $D^-(\mathcal{O})$.
Soit $\mathcal{F}$ un faisceau de $\mathcal{O}$-modules.
Soit $n \in \mathbf{Z}$.
\begin{enumerate}
\item Si $H^i(\alpha)$ est un isomorphisme pour $i \geq n$, alors
$H^i(\alpha \otimes_\mathcal{O}^\mathbf{L} \text{id}_\mathcal{F})$ est un
isomorphisme pour $i \geq n$.
\item Si $H^i(\alpha)$ est un isomorphisme pour $i > n$ et surjectif pour
$i = n$, alors
$H^i(\alpha \otimes_\mathcal{O}^\mathbf{L} \text{id}_\mathcal{F})$ est un
isomorphisme pour $i > n$ et surjectif pour $i = n$.
\end{enumerate}
\end{lemma}

\begin{proof}
Choisissons un triangle distingué
$$
K \to L \to C \to K[1]
$$
Dans le cas (2), nous voyons que $H^i(C) = 0$ pour $i \geq n$.
Par conséquent,
$H^i(C \otimes_\mathcal{O}^\mathbf{L} \mathcal{F}) = 0$ pour $i \geq n$
d'après le dual de Catégories dérivées,
Lemme \ref{derived-lemma-negative-vanishing}.
Cela montre à son tour que
$H^i(\alpha \otimes_\mathcal{O}^\mathbf{L} \text{id}_\mathcal{F})$ est un
isomorphisme pour $i > n$ et surjectif pour $i = n$.
Dans le cas (1), nous voyons en outre que
$H^{n - 1}(L) \to H^{n - 1}(C)$ est surjectif. En considérant le diagramme
$$
\xymatrix{
H^{n - 1}(L) \otimes_\mathcal{O} \mathcal{F} \ar[r] \ar[d] &
H^{n - 1}(C) \otimes_\mathcal{O} \mathcal{F} \ar@{=}[d] \\
H^{n - 1}(L \otimes_\mathcal{O}^\mathbf{L} \mathcal{F}) \ar[r] &
H^{n - 1}(C \otimes_\mathcal{O}^\mathbf{L} \mathcal{F})
}
$$
nous concluons que la flèche horizontale inférieure est surjective.
Combiné à ce qui précède, cela implique que
$H^n(\alpha \otimes_\mathcal{O}^\mathbf{L} \text{id}_\mathcal{F})$ est un
isomorphisme.
\end{proof}

\begin{lemma}
\label{lemma-canonical-class-obstruction}
Soit $\mathcal{C}$ un site. Soit $\mathcal{O} \to \mathcal{O}_0$ une
surjection de faisceaux d'anneaux dont le noyau est un faisceau d'idéaux
$\mathcal{I}$ de carré nul. Pour tout objet $K_0$ de
$D^-(\mathcal{O}_0)$, les assertions suivantes sont équivalentes :
\begin{enumerate}
\item la classe
$\omega(K_0) \in
\Ext^2_{\mathcal{O}_0}(K_0, K_0 \otimes_{\mathcal{O}_0} \mathcal{I})$
construite dans le Lemme \ref{lemma-canonical-class} est nulle,
\item il existe $K \in D^-(\mathcal{O})$ tel que
$K \otimes_\mathcal{O}^\mathbf{L} \mathcal{O}_0 = K_0$
dans $D(\mathcal{O}_0)$.
\end{enumerate}
\end{lemma}

\begin{proof}
Soit $K$ comme dans (2). Nous pouvons représenter $K$ par un complexe
$\mathcal{F}^\bullet$ de $\mathcal{O}$-modules plats, borné supérieurement.
Alors $\mathcal{F}_0^\bullet =
\mathcal{F}^\bullet \otimes_{\mathcal{O}} \mathcal{O}_0$ représente $K_0$ dans
$D(\mathcal{O}_0)$.
Puisque $d_{\mathcal{F}^\bullet} \circ d_{\mathcal{F}^\bullet} = 0$, car
$\mathcal{F}^\bullet$ est un complexe, la construction même de
$\omega(K_0)$ montre que cette classe est nulle.

\medskip\noindent
Supposons (1). Soient $\mathcal{F}^n$ et
$d : \mathcal{F}^n \to \mathcal{F}^{n + 1}$ comme dans la construction de
$\omega(K_0)$. La nullité de $\omega(K_0)$ implique que l'application
$$
\omega = d \circ d : \mathcal{F}_0^\bullet
\longrightarrow
(\mathcal{F}_0^\bullet \otimes_{\mathcal{O}_0} \mathcal{I})[2]
$$
est nulle dans $D(\mathcal{O}_0)$. Par définition de la catégorie dérivée
comme localisation de la catégorie homotopique des complexes de
$\mathcal{O}_0$-modules, il existe un quasi-isomorphisme
$\alpha : \mathcal{G}_0^\bullet \to \mathcal{F}_0^\bullet$ et des applications
de $\mathcal{O}_0$-modules
$h^n : \mathcal{G}_0^n \to
\mathcal{F}_0^{n + 1} \otimes_\mathcal{O} \mathcal{I}$ telles que
$$
\omega \circ \alpha =
d_{\mathcal{F}_0^\bullet \otimes \mathcal{I}} \circ h +
h \circ d_{\mathcal{G}_0^\bullet}
$$
Posons
$$
\mathcal{H}^n = \mathcal{F}^n \times_{\mathcal{F}^n_0} \mathcal{G}_0^n
$$
et définissons
$$
d' : \mathcal{H}^n \longrightarrow \mathcal{H}^{n + 1},\quad
(f^n, g_0^n) \longmapsto (d(f^n) - h^n(g_0^n), d(g_0^n))
$$
avec les notations évidentes, en utilisant
$\mathcal{F}_0^{n + 1} \otimes_{\mathcal{O}_0} \mathcal{I} =
\mathcal{F}^{n + 1} \otimes_\mathcal{O} \mathcal{I} =
\mathcal{I}\mathcal{F}^{n + 1} \subset \mathcal{F}^{n + 1}$.
On vérifie alors que $d' \circ d' = 0$ grâce au choix de $h^n$
et à la définition de $\omega$.
Ainsi, $\mathcal{H}^\bullet$ définit un objet de $D(\mathcal{O})$.
D'autre part, il existe une suite exacte courte de complexes de
$\mathcal{O}$-modules
$$
0 \to \mathcal{F}_0^\bullet \otimes_{\mathcal{O}_0} \mathcal{I} \to
\mathcal{H}^\bullet \to \mathcal{G}_0^\bullet \to 0
$$
Il reste à montrer que
$\mathcal{H}^\bullet \otimes_\mathcal{O}^\mathbf{L} \mathcal{O}_0$ est
isomorphe à $K_0$.
Choisissons un quasi-isomorphisme
$\mathcal{E}^\bullet \to \mathcal{H}^\bullet$, où
$\mathcal{E}^\bullet$ est un complexe borné supérieurement de
$\mathcal{O}$-modules plats. Nous obtenons un diagramme commutatif
$$
\xymatrix{
0 \ar[r] &
\mathcal{E}^\bullet \otimes_\mathcal{O} \mathcal{I} \ar[d]^\beta \ar[r] &
\mathcal{E}^\bullet \ar[d]^\gamma \ar[r] &
\mathcal{E}_0^\bullet \ar[d]^\delta \ar[r] &
0 \\
0 \ar[r] &
\mathcal{F}_0^\bullet \otimes_{\mathcal{O}_0} \mathcal{I} \ar[r] &
\mathcal{H}^\bullet \ar[r] &
\mathcal{G}_0^\bullet \ar[r] &
0
}
$$
Nous affirmons que $\delta$ est un quasi-isomorphisme. Comme
$H^i(\delta)$ est un isomorphisme pour $i \gg 0$, nous pouvons procéder par
récurrence descendante sur $n$, en supposant que $H^i(\delta)$ est un
isomorphisme pour $i \geq n$.
Observons que
$\mathcal{E}^\bullet \otimes_\mathcal{O} \mathcal{I}$ représente
$\mathcal{E}_0^\bullet \otimes_{\mathcal{O}_0}^\mathbf{L} \mathcal{I}$, que
$\mathcal{F}_0^\bullet \otimes_{\mathcal{O}_0} \mathcal{I}$ représente
$\mathcal{G}_0^\bullet \otimes_{\mathcal{O}_0}^\mathbf{L} \mathcal{I}$, et que
$\beta = \delta \otimes_{\mathcal{O}_0}^\mathbf{L} \text{id}_\mathcal{I}$
comme applications dans $D(\mathcal{O}_0)$. Cela est vrai parce que
$\beta =
(\alpha \otimes \text{id}_\mathcal{I})
\circ
(\delta \otimes \text{id}_\mathcal{I})$.
Supposons que $H^i(\delta)$ soit un isomorphisme en degrés $\geq n$.
Il en est alors de même pour $\beta$, d'après ce que nous venons de dire et le
Lemme \ref{lemma-induced-map}.
Nous pouvons ensuite considérer le diagramme
$$
\xymatrix{
H^{n - 1}(\mathcal{E}^\bullet \otimes_\mathcal{O} \mathcal{I})
\ar[r] \ar[d]^{H^{n - 1}(\beta)} &
H^{n - 1}(\mathcal{E}^\bullet) \ar[r] \ar[d] &
H^{n - 1}(\mathcal{E}_0^\bullet) \ar[r] \ar[d]^{H^{n - 1}(\delta)} &
H^n(\mathcal{E}^\bullet \otimes_\mathcal{O} \mathcal{I})
\ar[r] \ar[d]^{H^n(\beta)} &
H^n(\mathcal{E}^\bullet) \ar[d] \\
H^{n - 1}(\mathcal{F}_0^\bullet \otimes_\mathcal{O} \mathcal{I}) \ar[r] &
H^{n - 1}(\mathcal{H}^\bullet) \ar[r] &
H^{n - 1}(\mathcal{G}_0^\bullet) \ar[r] &
H^n(\mathcal{F}_0^\bullet \otimes_\mathcal{O} \mathcal{I}) \ar[r] &
H^n(\mathcal{H}^\bullet)
}
$$
D'après Homologie, Lemme \ref{homology-lemma-four-lemma}, nous voyons que
$H^{n - 1}(\delta)$ est surjectif.
Cela implique à son tour que $H^{n - 1}(\beta)$ est surjectif, d'après le
Lemme \ref{lemma-induced-map}.
En appliquant de nouveau Homologie, Lemme \ref{homology-lemma-four-lemma}, nous
voyons que $H^{n - 1}(\delta)$ est un isomorphisme.
L'affirmation résulte par récurrence ; ainsi, $\delta$ est un
quasi-isomorphisme, comme nous voulions le montrer.
\end{proof}

\begin{lemma}
\label{lemma-lift-map-complexes}
Soit $\mathcal{C}$ un site. Soit $\mathcal{O} \to \mathcal{O}_0$ une
surjection de faisceaux d'anneaux. Supposons données les informations
suivantes :
\begin{enumerate}
\item un complexe de $\mathcal{O}$-modules $\mathcal{F}^\bullet$,
\item un complexe $\mathcal{K}_0^\bullet$ de $\mathcal{O}_0$-modules,
\item un quasi-isomorphisme $\mathcal{K}_0^\bullet \to
\mathcal{F}^\bullet \otimes_\mathcal{O} \mathcal{O}_0$,
\end{enumerate}
Il existe alors un quasi-isomorphisme
$\mathcal{G}^\bullet \to \mathcal{F}^\bullet$ tel que l'application de
complexes
$\mathcal{G}^\bullet  \otimes_\mathcal{O} \mathcal{O}_0 \to
\mathcal{F}^\bullet \otimes_\mathcal{O} \mathcal{O}_0$ se factorise par
$\mathcal{K}_0^\bullet$ dans la catégorie homotopique des complexes de
$\mathcal{O}_0$-modules.
\end{lemma}

\begin{proof}
Posons $\mathcal{F}_0^\bullet =
\mathcal{F}^\bullet \otimes_\mathcal{O} \mathcal{O}_0$.
D'après Catégories dérivées, Lemme \ref{derived-lemma-make-surjective}, il
existe une factorisation
$$
\mathcal{K}_0^\bullet \to \mathcal{L}_0^\bullet \to \mathcal{F}_0^\bullet
$$
de l'application donnée, telle que la première flèche admette un inverse à
homotopie près et que la seconde soit surjective et scindée terme à terme.
Nous pouvons donc supposer que
$\mathcal{K}_0^\bullet \to \mathcal{F}_0^\bullet$ est surjective terme à
terme. Dans ce cas, nous prenons
$$
\mathcal{G}^n = \mathcal{F}^n \times_{\mathcal{F}^n_0} \mathcal{K}_0^n
$$
et tout est clair.
\end{proof}

\begin{lemma}
\label{lemma-inf-obs-map-defo-complex}
Soit $\mathcal{C}$ un site. Soit $\mathcal{O} \to \mathcal{O}_0$ une
surjection de faisceaux d'anneaux dont le noyau est un faisceau d'idéaux
$\mathcal{I}$ de carré nul. Soient $K, L \in D^-(\mathcal{O})$.
Posons $K_0 = K \otimes_\mathcal{O}^\mathbf{L} \mathcal{O}_0$ et
$L_0 = L \otimes_\mathcal{O}^\mathbf{L} \mathcal{O}_0$ dans
$D^-(\mathcal{O}_0)$. Étant donnée
$\alpha_0 : K_0 \to L_0$ dans $D(\mathcal{O}_0)$, il existe un élément
canonique
$$
o(\alpha_0) \in \Ext^1_{\mathcal{O}_0}(K_0,
L_0 \otimes_{\mathcal{O}_0}^\mathbf{L} \mathcal{I})
$$
dont l'annulation est une condition nécessaire et suffisante pour qu'il existe
une application $\alpha : K \to L$ dans $D(\mathcal{O})$ telle que
$\alpha_0 = \alpha \otimes_\mathcal{O}^\mathbf{L} \text{id}$.
\end{lemma}

\begin{proof}
Trouver une application $\alpha : K \to L$ relevant $\alpha_0$ revient à
trouver $\alpha : K \to L$ telle que la composée
$K \xrightarrow{\alpha} L \to L_0$ soit égale à la composée
$K \to K_0 \xrightarrow{\alpha_0} L_0$.
La suite exacte courte
$0 \to \mathcal{I} \to \mathcal{O} \to \mathcal{O}_0 \to 0$ donne lieu à un
triangle distingué canonique
$$
L \otimes_\mathcal{O}^\mathbf{L} \mathcal{I} \to
L \to
L_0 \to
(L \otimes_\mathcal{O}^\mathbf{L} \mathcal{I})[1]
$$
dans $D(\mathcal{O})$.
D'après Catégories dérivées,
Lemme \ref{derived-lemma-representable-homological}, la composée
$$
K \to K_0 \xrightarrow{\alpha_0} L_0 \to
(L \otimes_\mathcal{O}^\mathbf{L} \mathcal{I})[1]
$$
est nulle si et seulement si nous pouvons trouver
$\alpha : K \to L$ relevant $\alpha_0$. Cette composée est un élément de
$$
\Hom_{D(\mathcal{O})}(K, (L \otimes_\mathcal{O}^\mathbf{L} \mathcal{I})[1]) =
\Hom_{D(\mathcal{O}_0)}(K_0,
(L \otimes_\mathcal{O}^\mathbf{L} \mathcal{I})[1]) =
\Ext^1_{\mathcal{O}_0}(K_0,
L_0 \otimes_{\mathcal{O}_0}^\mathbf{L} \mathcal{I})
$$
par adjonction.
\end{proof}

\begin{lemma}
\label{lemma-first-order-defos-complex}
Soit $\mathcal{C}$ un site. Soit $\mathcal{O} \to \mathcal{O}_0$ une
surjection de faisceaux d'anneaux dont le noyau est un faisceau d'idéaux
$\mathcal{I}$ de carré nul. Soit $K_0 \in D^-(\mathcal{O})$.
Un relèvement de $K_0$ est un couple $(K, \alpha_0)$ formé d'un objet $K$ de
$D^-(\mathcal{O})$ et d'un isomorphisme
$\alpha_0 : K \otimes_\mathcal{O}^\mathbf{L} \mathcal{O}_0 \to K_0$ dans
$D(\mathcal{O}_0)$.
\begin{enumerate}
\item Étant donné un relèvement $(K, \alpha)$, le groupe des automorphismes du
couple est canoniquement le conoyau d'une application
$$
\Ext^{-1}_{\mathcal{O}_0}(K_0, K_0)
\longrightarrow
\Hom_{\mathcal{O}_0}(K_0, K_0 \otimes_{\mathcal{O}_0}^\mathbf{L} \mathcal{I})
$$
\item S'il existe un relèvement, alors l'ensemble des classes d'isomorphisme de
relèvements est un espace principal homogène sous
$\Ext^1_{\mathcal{O}_0}(K_0,
K_0 \otimes_{\mathcal{O}_0}^\mathbf{L} \mathcal{I})$.
\end{enumerate}
\end{lemma}

\begin{proof}
Un automorphisme de $(K, \alpha)$ est une application
$\varphi : K \to K$ dans $D(\mathcal{O})$ telle que
$\varphi \otimes_\mathcal{O} \text{id}_{\mathcal{O}_0} = \text{id}$.
Cela revient à dire que
$$
K \xrightarrow{\varphi - \text{id}} K \to
K \otimes_\mathcal{O}^\mathbf{L} \mathcal{O}_0
$$
est nulle. Nous en concluons que le groupe des automorphismes est le conoyau
d'une application
$$
\Hom_\mathcal{O}(K, K_0[-1])
\longrightarrow
\Hom_\mathcal{O}(K, K_0 \otimes_{\mathcal{O}_0}^\mathbf{L} \mathcal{I})
$$
en vertu du triangle distingué
$$
K \otimes_\mathcal{O}^\mathbf{L} \mathcal{I} \to
K \to
K \otimes_\mathcal{O}^\mathbf{L} \mathcal{O}_0 \to
(K \otimes_\mathcal{O}^\mathbf{L} \mathcal{I})[1]
$$
dans $D(\mathcal{O})$ et de Catégories dérivées,
Lemme \ref{derived-lemma-representable-homological}.
Pour traduire ceci dans les groupes du lemme, nous utilisons l'adjonction entre
le foncteur de restriction $D(\mathcal{O}_0) \to D(\mathcal{O})$ et
$- \otimes_\mathcal{O} \mathcal{O}_0 : D(\mathcal{O}) \to D(\mathcal{O}_0)$.
Cela démontre (1).

\medskip\noindent
Démonstration de (2).
Supposons que $K_0 = K \otimes_\mathcal{O}^\mathbf{L} \mathcal{O}_0$ dans
$D(\mathcal{O})$. D'après le
Lemme \ref{lemma-inf-obs-map-defo-complex}, l'application qui associe à un
relèvement $(K', \alpha_0)$ l'obstruction $o(\alpha_0)$ au relèvement de
$\alpha_0$ définit une injection canonique de l'ensemble des classes
d'isomorphisme de couples dans
$\Ext^1_{\mathcal{O}_0}(K_0,
K_0 \otimes_{\mathcal{O}_0}^\mathbf{L} \mathcal{I})$.
Pour achever la démonstration, montrons qu'elle est surjective.
Choisissons
$\xi : K_0 \to (K_0 \otimes_{\mathcal{O}_0}^\mathbf{L} \mathcal{I})[1]$
dans le groupe $\Ext^1$ du lemme.
Choisissons un complexe $\mathcal{F}^\bullet$ de $\mathcal{O}$-modules plats,
borné supérieurement, représentant $K$.
L'application $\xi$ peut être représentée sous la forme $t \circ s^{-1}$, où
$s : \mathcal{K}_0^\bullet \to \mathcal{F}_0^\bullet$ est un
quasi-isomorphisme et
$t : \mathcal{K}_0^\bullet \to
\mathcal{F}_0^\bullet \otimes_{\mathcal{O}_0} \mathcal{I}[1]$ est une
application de complexes.
D'après le Lemme \ref{lemma-lift-map-complexes}, nous pouvons supposer qu'il
existe un quasi-isomorphisme
$\mathcal{G}^\bullet \to \mathcal{F}^\bullet$ de complexes de
$\mathcal{O}$-modules tel que
$\mathcal{G}_0^\bullet \to \mathcal{F}_0^\bullet$ se factorise par $s$ à
homotopie près.
Nous pouvons et allons remplacer $\mathcal{G}^\bullet$ par un complexe borné
supérieurement de $\mathcal{O}$-modules plats (en choisissant un
quasi-isomorphisme d'un tel complexe vers $\mathcal{G}^\bullet$, puis en
effectuant le remplacement).
Nous voyons alors que $\xi$ est représenté par une application de complexes
$t : \mathcal{G}_0^\bullet \to
\mathcal{F}_0^\bullet \otimes_{\mathcal{O}_0} \mathcal{I}[1]$ et par le
quasi-isomorphisme
$\mathcal{G}_0^\bullet \to \mathcal{F}_0^\bullet$.
Posons
$$
\mathcal{H}^n = \mathcal{F}^n \times_{\mathcal{F}_0^n} \mathcal{G}_0^n
$$
avec les différentielles
$$
\mathcal{H}^n \to \mathcal{H}^{n + 1},\quad
(f^n, g_0^n) \mapsto (d(f^n) + t(g_0^n), d(g_0^n))
$$
Cela a un sens, car
$\mathcal{F}_0^{n + 1} \otimes_{\mathcal{O}_0} \mathcal{I} =
\mathcal{F}^{n + 1} \otimes_\mathcal{O} \mathcal{I} =
\mathcal{I}\mathcal{F}^{n + 1} \subset \mathcal{F}^{n + 1}$.
Nous omettons le calcul montrant que $\mathcal{H}^\bullet$ est un complexe de
$\mathcal{O}$-modules. Par construction, il existe une suite exacte courte
$$
0 \to \mathcal{F}_0^\bullet \otimes_{\mathcal{O}_0} \mathcal{I} \to
\mathcal{H}^\bullet \to \mathcal{G}_0^\bullet \to 0
$$
de complexes de $\mathcal{O}$-modules.
Exactement comme dans la démonstration du
Lemme \ref{lemma-canonical-class-obstruction}, on montre que cette suite induit
un isomorphisme
$\alpha_0 :
\mathcal{H}^\bullet \otimes_\mathcal{O}^\mathbf{L} \mathcal{O}_0 \to
\mathcal{G}_0^\bullet$ dans $D(\mathcal{O}_0)$.
Autrement dit, nous avons construit un couple
$(\mathcal{H}^\bullet, \alpha_0)$.
Nous omettons la vérification de $o(\alpha_0) = \xi$ ; indication :
$o(\alpha_0)$ peut être calculé explicitement dans ce cas, car nous avons des
applications $\mathcal{H}^n \to \mathcal{F}^n$ (non compatibles aux
différentielles) relevant les composantes de $\alpha_0$.
Cela achève la démonstration.
\end{proof}













\begin{multicols}{2}[\section{Autres chapitres}]
\noindent
Préliminaires
\begin{enumerate}
\item \hyperref[introduction-section-phantom]{Introduction}
\item \hyperref[conventions-section-phantom]{Conventions}
\item \hyperref[sets-section-phantom]{Théorie des ensembles}
\item \hyperref[categories-section-phantom]{Catégories}
\item \hyperref[topology-section-phantom]{Topologie}
\item \hyperref[sheaves-section-phantom]{Faisceaux sur les espaces}
\item \hyperref[sites-section-phantom]{Sites et faisceaux}
\item \hyperref[stacks-section-phantom]{Champs}
\item \hyperref[fields-section-phantom]{Corps}
\item \hyperref[algebra-section-phantom]{Algèbre commutative}
\item \hyperref[brauer-section-phantom]{Groupes de Brauer}
\item \hyperref[homology-section-phantom]{Algèbre homologique}
\item \hyperref[derived-section-phantom]{Catégories dérivées}
\item \hyperref[simplicial-section-phantom]{Méthodes simpliciales}
\item \hyperref[more-algebra-section-phantom]{Compléments d'algèbre}
\item \hyperref[smoothing-section-phantom]{Lissage des morphismes d'anneaux}
\item \hyperref[modules-section-phantom]{Faisceaux de modules}
\item \hyperref[sites-modules-section-phantom]{Modules sur les sites}
\item \hyperref[injectives-section-phantom]{Objets injectifs}
\item \hyperref[cohomology-section-phantom]{Cohomologie des faisceaux}
\item \hyperref[sites-cohomology-section-phantom]{Cohomologie sur les sites}
\item \hyperref[dga-section-phantom]{Algèbre différentielle graduée}
\item \hyperref[dpa-section-phantom]{Algèbre à puissances divisées}
\item \hyperref[sdga-section-phantom]{Faisceaux différentiels gradués}
\item \hyperref[hypercovering-section-phantom]{Hyperrecouvrements}
\end{enumerate}
Schémas
\begin{enumerate}
\setcounter{enumi}{25}
\item \hyperref[schemes-section-phantom]{Schémas}
\item \hyperref[constructions-section-phantom]{Constructions de schémas}
\item \hyperref[properties-section-phantom]{Propriétés des schémas}
\item \hyperref[morphisms-section-phantom]{Morphismes de schémas}
\item \hyperref[coherent-section-phantom]{Cohomologie des schémas}
\item \hyperref[divisors-section-phantom]{Diviseurs}
\item \hyperref[limits-section-phantom]{Limites de schémas}
\item \hyperref[varieties-section-phantom]{Variétés}
\item \hyperref[topologies-section-phantom]{Topologies sur les schémas}
\item \hyperref[descent-section-phantom]{Descente}
\item \hyperref[perfect-section-phantom]{Catégories dérivées de schémas}
\item \hyperref[more-morphisms-section-phantom]{Compléments sur les morphismes}
\item \hyperref[flat-section-phantom]{Compléments sur la platitude}
\item \hyperref[groupoids-section-phantom]{Schémas en groupoïdes}
\item \hyperref[more-groupoids-section-phantom]{Compléments sur les schémas en groupoïdes}
\item \hyperref[etale-section-phantom]{Morphismes étales de schémas}
\end{enumerate}
Thèmes de théorie des schémas
\begin{enumerate}
\setcounter{enumi}{41}
\item \hyperref[chow-section-phantom]{Homologie de Chow}
\item \hyperref[intersection-section-phantom]{Théorie de l'intersection}
\item \hyperref[pic-section-phantom]{Schémas de Picard des courbes}
\item \hyperref[weil-section-phantom]{Théories cohomologiques de Weil}
\item \hyperref[adequate-section-phantom]{Modules adéquats}
\item \hyperref[dualizing-section-phantom]{Complexes dualisants}
\item \hyperref[duality-section-phantom]{Dualité pour les schémas}
\item \hyperref[discriminant-section-phantom]{Discriminants et différentes}
\item \hyperref[derham-section-phantom]{Cohomologie de de Rham}
\item \hyperref[local-cohomology-section-phantom]{Cohomologie locale}
\item \hyperref[algebraization-section-phantom]{Géométrie algébrique et formelle}
\item \hyperref[curves-section-phantom]{Courbes algébriques}
\item \hyperref[resolve-section-phantom]{Résolution des surfaces}
\item \hyperref[models-section-phantom]{Réduction semi-stable}
\item \hyperref[functors-section-phantom]{Foncteurs et morphismes}
\item \hyperref[equiv-section-phantom]{Catégories dérivées de variétés}
\item \hyperref[pione-section-phantom]{Groupes fondamentaux de schémas}
\item \hyperref[etale-cohomology-section-phantom]{Cohomologie étale}
\item \hyperref[crystalline-section-phantom]{Cohomologie cristalline}
\item \hyperref[proetale-section-phantom]{Cohomologie pro-étale}
\item \hyperref[relative-cycles-section-phantom]{Cycles relatifs}
\item \hyperref[more-etale-section-phantom]{Compléments de cohomologie étale}
\item \hyperref[trace-section-phantom]{La formule des traces}
\end{enumerate}
Espaces algébriques
\begin{enumerate}
\setcounter{enumi}{64}
\item \hyperref[spaces-section-phantom]{Espaces algébriques}
\item \hyperref[spaces-properties-section-phantom]{Propriétés des espaces algébriques}
\item \hyperref[spaces-morphisms-section-phantom]{Morphismes d'espaces algébriques}
\item \hyperref[decent-spaces-section-phantom]{Espaces algébriques décents}
\item \hyperref[spaces-cohomology-section-phantom]{Cohomologie des espaces algébriques}
\item \hyperref[spaces-limits-section-phantom]{Limites d'espaces algébriques}
\item \hyperref[spaces-divisors-section-phantom]{Diviseurs sur les espaces algébriques}
\item \hyperref[spaces-over-fields-section-phantom]{Espaces algébriques sur des corps}
\item \hyperref[spaces-topologies-section-phantom]{Topologies sur les espaces algébriques}
\item \hyperref[spaces-descent-section-phantom]{Descente et espaces algébriques}
\item \hyperref[spaces-perfect-section-phantom]{Catégories dérivées d'espaces}
\item \hyperref[spaces-more-morphisms-section-phantom]{Compléments sur les morphismes d'espaces}
\item \hyperref[spaces-flat-section-phantom]{Platitude sur les espaces algébriques}
\item \hyperref[spaces-groupoids-section-phantom]{Groupoïdes dans les espaces algébriques}
\item \hyperref[spaces-more-groupoids-section-phantom]{Compléments sur les groupoïdes dans les espaces}
\item \hyperref[bootstrap-section-phantom]{Amorçage}
\item \hyperref[spaces-pushouts-section-phantom]{Sommes amalgamées d'espaces algébriques}
\end{enumerate}
Thèmes de géométrie
\begin{enumerate}
\setcounter{enumi}{81}
\item \hyperref[spaces-chow-section-phantom]{Groupes de Chow des espaces}
\item \hyperref[groupoids-quotients-section-phantom]{Quotients de groupoïdes}
\item \hyperref[spaces-more-cohomology-section-phantom]{Compléments sur la cohomologie des espaces}
\item \hyperref[spaces-simplicial-section-phantom]{Espaces simpliciaux}
\item \hyperref[spaces-duality-section-phantom]{Dualité pour les espaces}
\item \hyperref[formal-spaces-section-phantom]{Espaces algébriques formels}
\item \hyperref[restricted-section-phantom]{Algébrisation des espaces formels}
\item \hyperref[spaces-resolve-section-phantom]{Résolution des surfaces revisitée}
\end{enumerate}
Théorie des déformations
\begin{enumerate}
\setcounter{enumi}{89}
\item \hyperref[formal-defos-section-phantom]{Théorie formelle des déformations}
\item \hyperref[defos-section-phantom]{Théorie des déformations}
\item \hyperref[cotangent-section-phantom]{Le complexe cotangent}
\item \hyperref[examples-defos-section-phantom]{Problèmes de déformation}
\end{enumerate}
Champs algébriques
\begin{enumerate}
\setcounter{enumi}{93}
\item \hyperref[algebraic-section-phantom]{Champs algébriques}
\item \hyperref[examples-stacks-section-phantom]{Exemples de champs}
\item \hyperref[stacks-sheaves-section-phantom]{Faisceaux sur les champs algébriques}
\item \hyperref[criteria-section-phantom]{Critères de représentabilité}
\item \hyperref[artin-section-phantom]{Axiomes d'Artin}
\item \hyperref[quot-section-phantom]{Espaces de Quot et de Hilbert}
\item \hyperref[stacks-properties-section-phantom]{Propriétés des champs algébriques}
\item \hyperref[stacks-morphisms-section-phantom]{Morphismes de champs algébriques}
\item \hyperref[stacks-limits-section-phantom]{Limites de champs algébriques}
\item \hyperref[stacks-cohomology-section-phantom]{Cohomologie des champs algébriques}
\item \hyperref[stacks-perfect-section-phantom]{Catégories dérivées de champs}
\item \hyperref[stacks-introduction-section-phantom]{Introduction aux champs algébriques}
\item \hyperref[stacks-more-morphisms-section-phantom]{Compléments sur les morphismes de champs}
\item \hyperref[stacks-geometry-section-phantom]{La géométrie des champs}
\end{enumerate}
Thèmes de théorie des espaces de modules
\begin{enumerate}
\setcounter{enumi}{107}
\item \hyperref[moduli-section-phantom]{Champs de modules}
\item \hyperref[moduli-curves-section-phantom]{Espaces de modules de courbes}
\end{enumerate}
Divers
\begin{enumerate}
\setcounter{enumi}{109}
\item \hyperref[examples-section-phantom]{Exemples}
\item \hyperref[exercises-section-phantom]{Exercices}
\item \hyperref[guide-section-phantom]{Guide bibliographique}
\item \hyperref[desirables-section-phantom]{Desiderata}
\item \hyperref[coding-section-phantom]{Style de programmation}
\item \hyperref[obsolete-section-phantom]{Obsolète}
\item \hyperref[fdl-section-phantom]{Licence de documentation libre GNU}
\item \hyperref[index-section-phantom]{Index généré automatiquement}
\end{enumerate}
\end{multicols}


\bibliography{my}
\bibliographystyle{amsalpha}

\end{document}
